\documentclass{article}

\usepackage[preprint]{neurips_2026}

\usepackage[utf8]{inputenc}
\usepackage[T2A,T1]{fontenc}
\usepackage[english,russian]{babel}
\usepackage{tempora} % Times-совместимый шрифт с кириллицей (T2A)
\renewcommand{\sfdefault}{PTSans-TLF} % sans с кириллицей вместо phv
\usepackage{hyperref}
\hypersetup{unicode=true}
\usepackage{url}
\usepackage{booktabs}
\usepackage{amsfonts}
\usepackage{amssymb}
\usepackage{amsmath}
\usepackage{nicefrac}
\usepackage{microtype}
\usepackage{xcolor}
\usepackage{graphicx}
\usepackage{array}
\usepackage{enumitem}
\usepackage{placeins}
\usepackage{tikz}
\usetikzlibrary{arrows.meta,positioning}

\newcommand{\method}[1]{\texttt{#1}}
\newcommand{\wstar}{W^{\star}}
\newcommand{\rwstar}{RW^{\star}}
\newcommand{\mstar}{M^{\star}}
\newcommand{\cstar}{C^{\star}}
\newcommand{\phiroute}{\varphi_{\mathrm{route}}}

\setcounter{tocdepth}{2}

\title{Songlines: стадийное измерение и коллективная семантическая
память для многоагентной навигации\\
\large Объединённая серия (части I--V)}

\begin{document}
\maketitle
\begin{abstract}
Этот том объединяет пять статей исследовательской программы
Songlines в один документ. Программа спрашивает, как агенты с
символической семантической памятью выполняют навигационные задачи
--- и что меняется, когда память становится коллективной.
\textbf{Часть I} вводит стадийную декомпозицию Q/R/M/C (Query $\to$
Retrieval $\to$ Materialization $\to$ Completion): один
logging-контракт, локализующий отказы, которых не видят метрики
успеха, переносимый на пять классов систем --- от символических
одиночных агентов до LLM-стеков; острейший одноагентный результат ---
связывающее ограничение не в ранжировании, а в \emph{генерации}
кандидатов ($91\%$ пустых извлечений при точности ранжирования
$96{,}9\%$). \textbf{Часть II} делает «навигацию по чужой памяти»
измеряемым событием: каждое обязательство аннотировано чужой долей
$\varphi$ его свидетельств; сырые чужие локи почти никогда не
завершаются напрямую, но окупаются как транспорт; закон дистанции в
замкнутой форме предсказывает горизонты намерения с точностью до
тика; а двухстрочный координационный контракт (Warp Drive) делает
самый быстрый режим обмена и лучшим. \textbf{Часть III} находит
переносимую единицу: не точку (координату), а песню --- маршруты с
провенансом рёбер, наследующие завершение и безопасность свидетеля и
рвущиеся на предсказанном ребре; идентичность мест, восстановленную
из одного смысла (отказывающую закрыто там, где координаты отказывают
открыто); и результат анатомии: узлы несут идентичность, рёбра ---
перенос, за $6{,}4\%$ бит снапшота. \textbf{Часть IV} даёт
коллективной памяти механику --- слияние, доверие и устаревание как
три явных правила, занимающие новый режим на плоскости $M{\times}C$
--- и математику: категорную рамку, в которой морфизмы выравнивания
конструируются, а не предполагаются, а фазовая диаграмма
доверие$\times$каденция решает, охватывает ли копредел коллектив.
\textbf{Часть V} обращается к \emph{онтогенезу} памяти: формирование
как двухосевое решение (контрфактическая маргинальная полезность
$\times$ простота аналогии) с операцией исключения, защищающей от
ложной генерализации, --- измерено детерминированно, выучено
контекстным бандитом, обыгравшим собственного проектировщика, и
эволюционировано к конечному бюджету куплетов. Затем Часть V собирает
все механизмы в один рантайм и закрывает acceptance-тест рецензента:
runtime-оцениватель полезности, совпадающий с точным контрфактом
(Спирмен $0{,}989$), так что контроллер не нуждается в оракуле;
допуск с валидацией на стороне приёмника, делающий социальный обмен
нетто-полезным на горизонте, где нерегулируемый обмен вреден;
сравнение при равных бюджетах, где сырая история не догоняет и при
$4\times$ памяти; тест на adversarial-провенанс; safety-гейты,
доводящие fail-open ниже $1\%$ на сеточном и континуальном
субстратах; и полный учёт стоимости провода. Части писались как
самостоятельные подачи; перекрёстные ссылки между ними появляются
здесь как части одного тома.
\end{abstract}

\section*{Как читать этот том}

\paragraph{По одному предложению на часть.} Часть I строит инструмент
(стадийное измерение); часть II им расценивает чужие свидетельства
(провенанс и его закон); часть III находит правильную единицу
переноса (песня: ориентиры $+$ такты); часть IV даёт коллективную
механику и математику «иметь в виду одно и то же»; часть V
спрашивает, откуда вообще берётся память, которую координируют части
I--IV, и делает правило её формирования выучиваемым и
эволюционируемым.

\paragraph{Происхождение.} Части I и IV выросли из одной ранней
рукописи (\emph{Q/R/M/C Framework $+$ Minimal Collective Semantic
Memory}, хранится в архиве серии как \texttt{symbolic\_memory});
каждый эксперимент предка присутствует здесь, в уточнённом виде, в
частях I и IV. Части II и III --- компаньоны, построенные на
инструментарии части I. Там, где текст части говорит
«сопутствующая статья», имеется в виду другая часть этого тома.

\paragraph{Измерительная дисциплина.} Каждое утверждение каждой части
следует одному протоколу: предсказания регистрируются на диск до
исполнения, отказы относительно регистрации публикуются с
механизмами, а сильнейшие результаты --- точные совпадения с
предсказаниями в замкнутой форме ($6/6$ hold-out закона дистанции,
$110/110$ закона разрыва, $80/80$ восстановления $SE(2)$, $0$
fail-open в $252$ ячейках абляции словаря).

\paragraph{Код и данные.} Все эксперименты живут в одном репозитории
(\texttt{experiments/}, \texttt{songline\_drive/}); приложение о
воспроизводимости каждой части перечисляет точные команды. Сырые
артефакты прогонов --- в \texttt{tmp/} с файлами \texttt{RESULTS.md}
у ключевых экспериментов.

\clearpage
\tableofcontents
\clearpage
\part{Измерительный фреймворк Q/R/M/C}
\noindent\textit{Стадийная оценка навигации на основе памяти: один logging-контракт, пять классов систем и диагноз candidate generation.}

\paragraph{Аннотация части.} \textbf{Постановка.} Метрики, учитывающие только успех, схлопывают качественно различные отказы навигации на основе памяти в один скаляр, скрывая, \emph{какая стадия} конвейера извлечения (retrieval) за них ответственна. Данная статья вносит вклад в виде \emph{протокола оценивания} (evaluation protocol): постадийной декомпозиции Q/R/M/C (формирование запроса --- \textbf{Q}uery, удовлетворение извлечения --- \textbf{R}etrieval, материализация цели --- \textbf{M}aterialization, завершение после извлечения --- \textbf{C}ompletion), операционализированной так, что четыре условные частоты вычислимы непосредственно из стандартных логов траекторий (Определение~1). Арифметика --- это цепное правило, и это сделано намеренно: как декомпозиция accuracy на precision и recall не добавила математики, но изменила то, что поле смогло \emph{увидеть}, так и здесь ценность --- в идентифицируемости из логов и в том, что декомпозиция вскрывает, а не в теореме.
\textbf{Одноагентные измерения.} На 10-сидовом бенчмарке MiniGrid (FourRooms, LavaGapS7) и проверке переносимости на BabyAI фреймворк локализует бутылочные горлышки: восстановление после опасности ограничено извлечением (оракульное извлечение поднимает успех $0.39\to0.60$), возврат в целевой регион ограничен нижележащими стадиями. Инспекция логов уточняет атрибуцию: ранжирование кандидатов почти оптимально ($96.9\%$), тогда как генерация кандидатов отказывает в 91\% точек принятия решений --- это предел накопления памяти, а не ранжирования.
\textbf{Многоагентные измерения.} На многоагентном свипе из 35{,}640 прогонов по четырём коммуникационным архитектурам (независимая, общая шина, центральный агрегатор, широковещательная рассылка между пирами при восьми каденциях $K{\in}\{1,2,4,8,16,32,48,64\}$, 40 сидов) смещение бутылочного горлышка (bottleneck shift) между звеньями M и C проявляется напрямую как отрицательная внутрикаденционная корреляция с пиком при $K{=}4$ (Пирсон $-0.61$); среднее время до успеха имеет внутренний минимум при $K{=}8$ (устойчивый по положению, но мелкий по запасу --- не отличим от соседних быстрых каденций). Оба знака наклонов (Эмпирическое утверждение~1) подтверждаются как умеренные, кластерно-робастные размеры эффекта (Спирмен $-0.53$ и $+0.43$ от $K$, 95\%-е ДИ $[-0.59,-0.47]$ и $[+0.37,+0.50]$ по 81 конфигурационной ячейке; Приложение~\ref{p1:app:effect-sizes}). Не-утверждения (в частности, об условном произведении $\Phi$) собраны в одном месте (\S1, «Границы применимости и не-утверждения»).
\textbf{Диагностическое применение.} Те же четыре частоты также показывают, почему обученные политики с неявной коммуникацией могут быть конкурентоспособны по успеху, при этом схлопывая наблюдаемость Q/R/M/C: их коммуникационные векторы никогда не бывают информативно пустыми, и у них нет явного интерфейса фиксации цели. На двух локальных LLM-агентах (qwen3:4b, llama3.1 на TextNav) тот же логгер локализует отказ, ограниченный завершением: при жёстком бюджете шагов оба агента формируют запрос, извлекают и фиксируют цель ($Q^\star{=}R^\star{=}M^\star{=}1$), но ни один не завершает ($C^\star{=}0$), так что в остальном непрозрачный провал успеха разрешается в конкретную отказывающую стадию, а не в число в лидерборде. Таким образом, вклад --- это измерительный протокол, делающий отказы памяти сопоставимыми между символьными, обученными, LLM-, одноагентными и многоагентными системами.
\textbf{Границы применимости.} Основные эксперименты проведены в grid-world-субстратах (MiniGrid, BabyAI, собственная сетка $12{\times}10$). Вклад --- диагностический фреймворк и его валидация; предварительная проверка в непрерывном пространстве состояний на VMAS воспроизводит оба знака наклонов направленно (Приложение~\ref{p1:app:vmas}), а полная валидация в непрерывных/3D-средах с кластерно-робастными ДИ (Habitat, ProcTHOR) --- предмет будущей работы.

\FloatBarrier
\section{Введение}

\begin{figure}[t]
  \centering
  \includegraphics[width=\linewidth]{figures/fig_hero_qrmc.pdf}
  \caption{Q/R/M/C раскладывает один эпизод семантической навигации на четыре
  наблюдаемые стадии. \textbf{Сверху:} четыре стадии и их определения на уровне
  событий. \textbf{(a)}~Успешный эпизод поиска воды закрепляется
  на одном событии на стадию, порождая логируемые «звёздные» события
  $Q^\star, R^\star, M^\star, C^\star$. \textbf{(b)}~Одна и та же скалярная
  сводка «сквозного успеха» маскирует качественно различные отказы
  на четырёх одноагентных семействах задач: вода и отдых сильны;
  возврат в целевой регион (FourRooms) ограничен нижележащими стадиями; восстановление
  после опасности (LavaGapS7) ограничено извлечением (оракульные интервенции, \S\ref{p1:sec:oracle}).}
  \label{p1:fig:hero}
\end{figure}

Классические навигационные бенчмарки часто предполагают, что агенту дана
явная координата пункта назначения или плотная награда за задачу. Многие практические
навигационные задачи устроены иначе: агент должен достичь места определённого
\emph{типа} --- безопасной зоны отдыха, водоподобного ресурса,
выхода для восстановления после опасности или региона, ассоциированного с целью, ---
заданного семантическим паттерном, а не фиксированной путевой точкой $(x, y)$. Следовательно, место
должно быть \emph{запомнено и извлечено по смыслу}. Эта
постановка открывает методологический вопрос: можем ли мы измерить, \emph{где}
ломается семантическая навигация на основе памяти, а не только \emph{ломается ли}
она?

Мы отвечаем на этот вопрос четырёхстадийным диагностическим фреймворком,
применимым как к одиночным агентам, так и к многоагентным коллективам.
Декомпозиция --- Q/R/M/C: формирование запроса (\textbf{Q}uery formation), удовлетворение
извлечения (\textbf{R}etrieval satisfaction), материализация цели (target \textbf{M}aterialization) и завершение
после извлечения (\textbf{C}ompletion after retrieval). Декомпозиция факторизует поэпизодное завершение семантического
пути, и при условном постадийно-марковском свойстве её факторы
оцениваемы из логов траекторий (Определение~1).
Остальная часть статьи развивает единый аргумент в трёх частях.

\textbf{Одиночный агент.} Постадийная атрибуция нетривиальна.
Метрики только успеха схлопывают качественно различные отказы в
один скаляр. В 10-сидовом бенчмарке MiniGrid задачи воды и отдыха
сильны сквозным образом (успех 0.95, 0.93), но возврат в целевой регион и
восстановление после опасности существенно труднее, причём по \emph{разным
причинам}. Оракульные постадийные интервенции показывают, что восстановление после опасности
ограничено извлечением (оракульное извлечение поднимает успех $0.39 \to 0.60$),
тогда как возврат в целевой регион ограничен нижележащими стадиями (оракульное извлечение
спасает завершение семантического пути, не меняя сквозного успеха).
Инспекция собственных логов фреймворка затем заостряет атрибуцию:
бутылочное горлышко извлечения --- не \emph{ранжирование} кандидатов (точность
$96.9\%$, когда кандидаты существуют), а их \emph{генерация}:
память не содержит подходящих символов в 91\% точек принятия решений.
Одноагентное бутылочное горлышко --- предел \emph{накопления памяти}:
свидетельства наблюдаются, но так и не консолидируются в структуру, доступную запросам.

\textbf{Много агентов.} Консолидация обретает темп.
Когда $N$ агентов разделяют память через широковещательный протокол с каденцией (cadence)
$K$, факторизация расширяется укрупнёнными множествами обусловливания
(Определение~2), и $K$ становится структурной осью: механически это
ручка, регулирующая, \emph{насколько быстро приватные памяти консолидируются
в коллективную}. Утверждение на уровне наклонов
(Эмпирическое утверждение~1) предсказывает смещение бутылочного горлышка: с уменьшением $K$
$P(M^\star|R^\star)$ растёт (более свежая память пиров позволяет
надёжно коммититься), а $P(C^\star|M^\star)$ падает (пиры соревнуются за
одни и те же цели). Оба знака наклонов подтверждены на свипе из 35{,}640 прогонов
как умеренные, кластерно-робастные размеры эффекта (Приложение~\ref{p1:app:effect-sizes}); сдвиг напрямую виден как отрицательная
внутрикаденционная корреляция между двумя звеньями, а результирующий
компромисс порождает внутренний оптимум среднего времени до успеха при
$K\!=\!8$: у консолидации есть нетривиальный оптимальный темп.

\textbf{Диагностический синтез.} Оба режима указывают на один и тот же
объект измерения. Одноагентный диагноз говорит, что консолидация
свидетельств в символы может быть связывающим ограничением; многоагентный
диагноз говорит, что темп и цена обмена
консолидированной памятью определяют оптимум. Q/R/M/C --- общий
инструмент: системы-кандидаты можно сравнивать по тому, как они сдвигают
$P(R^\star|Q^\star)$, $P(M^\star|R^\star)$ и
$P(C^\star|M^\star)$, даже когда их сквозные показатели успеха выглядят
похоже.

\begin{center}
\fbox{\parbox{0.94\linewidth}{
\textbf{Границы применимости и не-утверждения (консолидировано).} Все оговорки этой
статьи в одном месте; основной текст их больше не повторяет.
(i)~Аргумент состоятельности декомпозиции --- это цепное правило плюс
частотное оценивание --- математическая новизна не заявляется; вклад ---
дисциплина логирования и идентифицируемость
(Определения~1--2). (ii)~Эмпирическое утверждение~1 --- проверенное наблюдение
на уровне наклонов об условных частотах, а не теорема.
(iii)~Мы не делаем утверждений о форме условного произведения
$\Phi$; в наших данных оно монотонно по $K$, а внутренний оптимум
живёт на среднем времени до успеха (\S\ref{p1:sec:multi-agent-results}).
(iv)~Основные субстраты --- grid worlds; VMAS --- направленная
проверка в непрерывном пространстве состояний, а LLM-исследования малы и контролируемы
(две локальные модели на TextNav; коллектив $N{=}3$; внешняя валидация
на трёх фреймворках при $n{=}8$ на ячейку, \S\ref{p1:sec:external}). (v)~Идентичность мест в многоагентном случае предполагает
общую систему координат (\S\ref{p1:sec:limitations}). Постадийная
операционализация проверяема на конкретных сводках и может быть
\emph{неверной} на других --- (iii) как раз такой случай, --- что и
отличает измерительный протокол от догадки, безразличной к стадиям.}}
\end{center}

\textbf{Почему такая рамка своевременна.} Воплощённые и целеобусловленные
агенты всё чаще извлекают места, объекты или инструменты по смыслу
\citep{andrychowicz2017hindsight,ghosh2021learning,savva2019habitat}.
Системы LLM-агентов опираются на явные слои памяти, но оцениваются в основном
сквозным образом \citep{wang2023voyager,shinn2023reflexion,packer2023memgpt}.
Многоагентное обучение с подкреплением также всё чаще вовлекает
коммуницирующие популяции \citep{lowe2017maddpg,foerster2018ctde}, однако
вопросу «какой агент должен знать что и когда?» недостаёт
единого диагностического словаря. Наш фреймворк его даёт.

\textbf{Вклады.} Это методологическая статья. Мы явно разделяем, что здесь бухгалтерия, что эмпирическое утверждение, а что инстанцированный артефакт.
\begin{enumerate}[leftmargin=20pt,itemsep=2pt,topsep=2pt]
\item \textbf{Измерительный протокол (один агент).} Декомпозиция Q/R/M/C операционализирована так, что условные частоты оцениваемы из логов траекторий (Определение~1); вклад --- \emph{дисциплина логирования}, делающая декомпозицию применимой к стандартным навигационным трассам (не-утверждения: см. бокс выше).

\item \textbf{Измерительный протокол (много агентов).} Расширение на $N$ коммуницирующих агентов с двумя каналами связывания (память $\mathcal{M}_{-i}^{(K)}$, занятость $\mathcal{O}_{-i}$) при протоколе широковещательной рассылки между пирами с каденцией $K$ (Определение~2).

\item \textbf{Эмпирическое утверждение~1 (смещение бутылочного горлышка).} На многоагентном свипе из 35{,}640 прогонов оба стилизованных условия на наклоны \textbf{(F)} и \textbf{(O)} подтверждены как умеренные, кластерно-робастные размеры эффекта (Спирмен $-0.53$ и $+0.43$ от $K$, 95\%-е ДИ, исключающие ноль, по 81 конфигурационной ячейке; Приложение~\ref{p1:app:effect-sizes}; полный бенчмарк в \S\ref{p1:sec:multi-agent-results}), с механизмом на уровне провенанса (provenance) (\S\ref{p1:sec:multi-agent-results}).

\item \textbf{Система и бенчмарки.} Символьно-семантический стек памяти с четырьмя коммуникационными архитектурами (независимая, общая шина, центральный агрегатор, широковещательная рассылка peer-to-peer); одноагентный оракульный бенчмарк (10-сидовый MiniGrid, переносимость на BabyAI); многоагентный свип из 35{,}640 прогонов при $N\!\in\!\{3,5,8\}$ + масштабирование на 8{,}640 прогонов при $N\!=\!12$; свип переносимости из 100 прогонов на обёртке стандартного MiniGrid.

\item \textbf{Диагностическое сравнение архитектур.} Один и тот же контракт логирования отделяет символьные системы памяти от обученных политик с неявной коммуникацией: последние могут быть конкурентоспособны по успеху, схлопывая при этом постадийную наблюдаемость Q/R/M/C. Это результат оценивания, а не утверждение, что один класс архитектур должен всегда доминировать.

\item \textbf{Границы и честность.} Основные эксперименты --- в grid-world-субстратах. Мы рассматриваем это как выбор, характерный для методологической статьи, и заявляем его явно (\S\ref{p1:sec:limitations}); включена предварительная проверка VMAS на непрерывном субстрате (Приложение~\ref{p1:app:vmas}), а полная валидация в непрерывных/3D-средах (Habitat, ProcTHOR, MineDojo) --- будущая работа.
\end{enumerate}

\section{Декомпозиция Q/R/M/C}
\label{p1:sec:framework}

\subsection{Постановка задачи и нотация}
\label{p1:sec:notation}

Агент действует в частично наблюдаемой навигационной среде.
На каждом шаге он получает наблюдение $o_t$, поддерживает состояние памяти
$\mathcal{M}_t$ и формирует запрос $q_t$, кодирующий текущее
намерение задачи. Шаг извлечения возвращает из $\mathcal{M}_t$ множество мест-кандидатов,
соответствующих $q_t$; выбранный кандидат
материализуется в действенную цель; локальный контроллер исполняет
действия, пока цель не достигнута или эпизод не истёк по времени.
Цель --- не координата, а семантический паттерн, поэтому успех есть
функция семантического профиля выбранного кандидата, а не какой-либо
заранее заданной координаты.

Habitat, AI2-THOR и ProcTHOR сделали визуально богатое оценивание навигации
стандартом
\citep{savva2019habitat,kolve2017ai2thor,deitke2022procthor}. Мы
намеренно фокусируемся на средах класса MiniGrid, чтобы получить чистые и
воспроизводимые постадийные измерения, а BabyAI используем только как проверку
переносимости. Мы не претендуем здесь на решение проблемы памяти LLM-агентов; скорее, мы
утверждаем, что та же декомпозиция Q/R/M/C --- полезный шаблон для
оценивания решений агентов, опирающихся на память, за пределами навигации.

Центральный эмпирический вопрос, таким образом, не просто в том, успешен ли агент,
а в том, \emph{какая стадия отказывает первой}, когда он неуспешен. Используются два
связанных измерительных слоя. \textbf{Операциональные частоты} считают,
на каждую попытку запроса, дал ли запрос непустое множество кандидатов,
удовлетворил ли выбранный кандидат семантическому предикату
и был ли он материализован в действенную цель;
\emph{завершение после извлечения} --- индикатор уровня эпизода.
\textbf{Вложенные события уровня эпизода}
$Q^\star \Rightarrow R^\star \Rightarrow M^\star \Rightarrow C^\star$
используются в факторизации ниже и аудируются в
\S\ref{p1:sec:validation}. Полные операциональные определения даны в
Приложении~\ref{p1:app:estimators}.

\textbf{Нотация.} Четыре стадии --- $Q$, $R$, $M$, $C$ со «звёздными» событиями уровня эпизода $Q^\star,R^\star,M^\star,C^\star\!\in\!\{0,1\}$, вложенными по построению. $P(\cdot)$ обозначает популяционную вероятность; $\hat{P}(\cdot)$ --- её эмпирико-частотную оценку. \textbf{Буква $M$ всегда обозначает стадию материализации; число целей в многоагентной среде записывается как $T$ (не $M$, ограничение $T{\le}N$, «дефицит» = $N{>}T$).} Прочие символы: $\mathcal{C}_K$ --- протокол широковещательной рассылки между пирами с каденцией $K$, $\mathcal{M}_{-i}^{(K)}$ --- снимки памяти пиров, $\mathcal{O}_{-i}$ --- занятость пиров, $\nu(\cdot)$ --- функционал свежести, $\Phi(K){=}P(M^\star|R^\star;K)\!\cdot\!P(C^\star|M^\star;K)$, $t_{\mathrm{succ}}$ --- среднее время до успеха. Полная таблица символов в Приложении~\ref{p1:app:notation}.

\subsection{Одноагентная факторизация}
\label{p1:sec:single-theory}

Пусть $Y$ обозначает общий успех эпизода, а
$Q^\star, R^\star, M^\star, C^\star$ --- вложенные события семантического пути
уровня эпизода. Фреймворк факторизует \emph{завершение семантического
пути}, а не сырой успех задачи:
\[
Q^\star \rightarrow R^\star \rightarrow M^\star \rightarrow C^\star.
\]
Для любого запроса $q$ и состояния памяти $\mathcal{M}$
\begin{align*}
P(C^\star = 1 \mid q, \mathcal{M})
&=
P(Q^\star = 1 \mid q, \mathcal{M})\;
P(R^\star = 1 \mid Q^\star = 1, q, \mathcal{M}) \\
&\quad \cdot
P(M^\star = 1 \mid R^\star = 1, q, \mathcal{M})\;
P(C^\star = 1 \mid M^\star = 1, q, \mathcal{M}).
\end{align*}
Общий успех задачи $Y$ может превышать $C^\star$, поскольку некоторые эпизоды
успешны без прохождения полного семантического пути извлечения.

\paragraph{Предположение 1 (условное постадийно-марковское свойство).}
При условии успешности текущей стадии вероятность
следующей стадии зависит только от текущей успешной стадии,
текущего запроса и текущего состояния памяти. Формально,
\[
P(R^\star \mid Q^\star, q, \mathcal{M}, \text{история}) = P(R^\star \mid Q^\star, q, \mathcal{M}),
\]
и аналогично для $M^\star$ и $C^\star$.

\paragraph{Предположение 2 (положительность).}
Каждое используемое выше событие обусловливания имеет ненулевую вероятность на
носителе бенчмарка.

\paragraph{Определение 1 (одноагентная оцениваемость стадий).}
При Предположениях 1--2 четыре условные частоты в факторизации выше являются корректно определёнными популяционными величинами, а эмпирико-частотные оценки --- их состоятельными оценками; их произведение сходится к завершению семантического пути. Аргумент элементарен --- факторизация по цепному правилу плюс состоятельность частотных оценок на бернуллиевских событиях --- и ценность его формулирования операциональна, а не математична: он фиксирует \emph{дисциплину логирования} (какие события, при каком обусловливании), при которой декомпозиция читаема из стандартных логов траекторий, и является контрактом, на соответствие которому аудируется каждый эксперимент этой статьи (\S\ref{p1:sec:validation}). Полная формулировка и элементарный аргумент в Приложении~\ref{p1:app:proofs}.

\paragraph{Когда Предположение~1 может нарушаться.}
Предположение~1 здесь выполняется жёстко, потому что $q$ и $\mathcal{M}$
явно логируются при каждом стадийном событии. Остаются два практически значимых
нарушения: адаптивные внутриэпизодные переписывания запроса, не всплывающие в
логгер, и скрытые межстадийные обновления памяти. В этих случаях
факторизацию следует читать как проекцию на наблюдаемые
трассы, а не как полное утверждение о структурной идентификации.
Приложение~\ref{p1:app:assumption1-stress} количественно оценивает смещение оценки
при контролируемом нарушении такого рода: оно пренебрежимо при мягких
нелогируемых связываниях и однознаково (смещение вверх у
$P(C^\star|M^\star)$), что ограничивает уязвимость фреймворка при его
применении к обученным или LLM-агентам.

\paragraph{Механизменное прочтение.}
Та же декомпозиция допускает стадийно-структурированную механизменную
интерпретацию. Рисунок~\ref{p1:fig:causal} трактует исходы стадий как
направленный ациклический граф. При такой интерпретации абляции включения/выключения
ассиста --- это нацеленные на механизм возмущения материализации и
пост-извлеченческого исполнения, а не недифференцированный анализ
чувствительности; в многоагентной постановке тот же граф реплицируется на
агента и связывается через каденционный протокол $\mathcal{C}_K$,
который действует на $\mathcal{M}_{-i}^{(K)}$ на ребре R$\to$M и на
$\mathcal{O}_{-i}$ на ребре M$\to$C --- ровно те два ребра, где
Эмпирическое утверждение~1 предскажет сдвиг наклонов.

\begin{figure}[t]
  \centering
  \begin{tikzpicture}[
    node distance=7mm and 7mm,
    box/.style={draw, rounded corners, align=center, minimum height=8mm, minimum width=15mm, fill=blue!8},
    arrow/.style={-{Latex[length=2.0mm]}, thick},
    assist/.style={draw, rounded corners, align=center, minimum height=7mm, minimum width=18mm, fill=orange!15}
  ]
    \node[box] (q) {$Q$};
    \node[box, right=of q] (r) {$R$};
    \node[box, right=of r] (m) {$M$};
    \node[box, right=of m] (c) {$C$};
    \node[box, right=of c] (s) {успех};
    \node[assist, above=of m] (a) {абляция\\ассиста};
    \draw[arrow] (q) -- (r);
    \draw[arrow] (r) -- (m);
    \draw[arrow] (m) -- (c);
    \draw[arrow] (c) -- (s);
    \draw[arrow] (a) -- (m);
    \draw[arrow] (a) -- (c);
  \end{tikzpicture}
  \caption{Стадийно-структурированная интерпретация фреймворка Q/R/M/C. Абляции включения/выключения ассиста нацелены на промежуточные механизмы; в многоагентном расширении каденционный протокол $\mathcal{C}_K$ действует на рёбра R$\to$M и M$\to$C.}
  \label{p1:fig:causal}
\end{figure}

\paragraph{Связь с предшествующими работами.} Ближайшие направления: (i) фреймворки когнитивных и семантических карт \citep{tolman1948cognitive,okeefe1978hippocampus,kuipers2000spatial,gupta2017cognitive,savinov2018semi}, (ii) целеобусловленный RL и RL с репрезентациями преемника \citep{andrychowicz2017hindsight,ghosh2021learning,dayan1993improving,stachenfeld2017hippocampus,momennejad2017successor}, (iii) многоагентная коммуникация \citep{lowe2017maddpg,foerster2018ctde,sukhbaatar2016commnet,foerster2016learning} и gossip/консенсус \citep{boyd2006gossip,xiao2007distributed}, а также (iv) слои памяти LLM-агентов, оцениваемые сквозным образом \citep{wang2023voyager,shinn2023reflexion,packer2023memgpt}. Наш фреймворк ортогонален всем четырём: он \emph{диагностирует}, где в цепочке Q/R/M/C данная архитектура успешна или отказывает, и поднимает ось каденции с уровня сетевых протоколов до параметра когнитивной архитектуры (полное обсуждение в Приложении~\ref{p1:app:related-work-extended}).

\subsection{Семантика событий, ограниченная рациональность и теоретико-порядковое прочтение}
\label{p1:sec:event-semantics}

Четыре «звёздных» события закрепляются на конкретных логируемых триггерах. Мы формулируем триггеры один раз, поскольку каждое эмпирическое утверждение статьи --- это утверждение о них (формальные определения по режимам в Приложении~\ref{p1:app:estimators}), и даём каждой стадии прочтение в терминах ресурса, который она ограничивает.

\begin{description}[leftmargin=10pt,topsep=2pt,itemsep=2pt]
\item[$Q^\star$ (намерение выпущено).] Планировщик выпускает семантическое намерение $q=(q^{\mathrm{req}},q^{\mathrm{pref}},q^{\mathrm{pen}})$: малые множества обязательных, предпочтительных и штрафных тегов. Событие уровня эпизода срабатывает при выпуске, поэтому оно насыщено всякий раз, когда задача семантична ($\hat P(Q^\star){=}1.00$ в Таблице~\ref{p1:tab:consistency}); многоагентный потиковый логгер дополнительно требует непустого множества кандидатов (Приложение~\ref{p1:app:multi-agent-defs}) --- бухгалтерское различие, которое мы фиксируем, а не прячем. Ресурсное ограничение: внимание --- намерение состоит из нескольких тегов, а не полного состояния.
\item[$R^\star$ (память отвечает).] Извлечение оценивает $q$ по памяти и возвращает множество кандидатов; событие срабатывает, когда множество содержит кандидата, удовлетворяющего предикату задачи (в пределах $\epsilon$ от истинной цели в многоагентном логгере). Отказы с пустым множеством кандидатов (91\%) из \S\ref{p1:sec:refined-attribution} --- события на стороне $R$: ни одно сохранённое место не проходит порог обязательных тегов. Ресурсное ограничение: память --- извлечение может поднять на поверхность лишь то, что консолидация сохранила.
\item[$M^\star$ (фиксация цели).] Планировщик материализует одного кандидата в конкретную действенную клетку. Само кандидатство пропускается через порог (место входит в множество кандидатов, только если его балл проходит гейт планировщика, Приложение~\ref{p1:app:schemas}), и фиксация коммитится на одного кандидата один раз на запрос, а не переоптимизируется по всей памяти на каждом тике. Событие срабатывает при захвате фиксации (в пределах $\epsilon$ от истинной цели в многоагентном логгере). Ресурсное ограничение: делиберация.
\item[$C^\star$ (замкнутое завершение).] Локальный контроллер исполняет, пока зафиксированная клетка не достигнута или эпизод не истёк; событие срабатывает по прибытии. $C$ --- единственная стадия, чьи моды отказа внесемантичны: компетентность контроллера и, в многоагентном режиме, занятость $\mathcal{O}_{-i}$. Ресурсное ограничение: управление, а в коллективе --- разделяемые цели.
\end{description}

\paragraph{Почему детерминированная стратегия: ограниченная и экологическая рациональность.}
Естественное возражение --- что семантический стек является фиксированной детерминированной стратегией, а не обученной политикой. Мы считаем детерминизм принципиальным и называем его: конвейер --- это satisficing-агент в смысле Саймона \citep{simon1955behavioral,simon1956rational}. Кандидатство пропускается через порог, а не исчерпывающе скорится; коммит на цель совершается однажды, а не переоптимизируется на каждом тике; намерение --- несколько тегов, а не полное состояние. Ограниченная рациональность --- это ещё и то, что делает четыре локуса отказа наблюдаемыми: каждая стадия --- отдельное ресурсное ограничение (внимание, память, делиберация, управление), а Q/R/M/C --- аудиторский след того, где ограничение связывает. Обученные бейзлайны эмпирически демонстрируют обратное: политики без явного интерфейса запроса/фиксации схлопывают стадии даже при конкурентоспособном успехе (BC в \S\ref{p1:sec:single-results}; CommNet PPO в \S\ref{p1:sec:multi-agent-results}), то есть наблюдаемость выторговывается ровно там, где выторговывается сквозная оптимизация. То, что \emph{одна и та же} эвристика сильна на воде и отдыхе, голодает по извлечению на восстановлении после опасности и ограничена нижележащими стадиями на целевом регионе (Рисунок~\ref{p1:fig:hero}b), --- это тезис экологической рациональности: успех эвристики есть свойство пары эвристика--среда \citep{simon1956rational,gigerenzer1996reasoning}; фреймворк --- инструмент, локализующий рассогласование. Та же линза читает расщепление $\Phi$-против-$t_{\mathrm{succ}}$ из \S\ref{p1:sec:multi-agent-results}: цепная вероятность $\Phi$ игнорирует цену прохождения цепочки, поэтому при ресурсно-рациональном прочтении \citep{lieder2020resource} оптимум следует искать на метрике, несущей цену; именно там он и появляется ($K{=}8$ на среднем времени до успеха), а не на $\Phi$.

\paragraph{Теоретико-порядковое прочтение (вынесено в Приложение~\ref{p1:app:galois}).}
Интерфейс запрос--извлечение несёт стандартную теоретико-порядковую структуру: места и теги образуют формальный контекст, чья связь Галуа реализует решётку понятий памяти, так что извлечение отказывает пустым множеством ровно тогда, когда объём запрошенного намерения пуст --- «отсутствие свидетельств», режим $91\%/96.9\%$ из \S\ref{p1:sec:refined-attribution}, и режим, в котором рост контекста может только наращивать объёмы. Мы используем это как именующий приём, а не как вклад; формулировка, её четырёхстрочный аргумент и границы применимости --- в Приложении~\ref{p1:app:galois}.

\paragraph{Динамическое прочтение (options / H-MDP) (вынесено в Приложение~\ref{p1:app:hmdp}).}
Стадии также допускают прочтение через темпорально протяжённые действия: каждый запрос инстанцирует опцию $o=\langle\mathcal{I},\pi,\beta\rangle$ \citep{sutton1999between}, чью внутреннюю структуру Q/R/M/C аудирует ребро в ребро (инициация $\to$ политика верхнего уровня над опциями $\to$ исполнение до терминации), а длительность исполнения $\tau$ в полу-MDP структурно объясняет, почему внутренний оптимум появляется на $t_{\mathrm{succ}}$, но не на слепой к $\tau$ цепной вероятности $\Phi$. Как и теоретико-порядковое прочтение, это интерпретирующая линза, а не вклад; полное развитие, глосса belief--desire--intention (общая с компаньоном о коллективной памяти) и многоагентный расчёт столкновений опций для смещения бутылочного горлышка --- в Приложении~\ref{p1:app:hmdp}.

\section{Один агент: постадийная диагностика локализует отказ}
\label{p1:sec:single}

\subsection{Система, задачи и протокол}
\label{p1:sec:single-setup}

\paragraph{Стек памяти.}
Система использует графовый субстрат с интерфейсом семантического извлечения поверх него: наблюдения $o_t$ токенизируются в символьные токены сцены и семантические теги, обновляющие память записей мест, переходов, эпизодических следов и записей понятий/отношений. Запросы планировщика возвращают кандидатов по взвешенному лог-баллу над обязательными, предпочтительными и штрафными тегами. Механизм извлечения переиспользуется во всех задачах и архитектурах; полная схема памяти, правила обновления и формула скоринга --- в Приложении~\ref{p1:app:schemas}.

\paragraph{Задачи, бейзлайны, протокол.}
Оцениваются четыре одноагентных семейства задач (вода, безопасный отдых, целевой регион на Empty-6x6 / FourRooms~\citep{sutton1999between}, восстановление после опасности на LavaGapS7), 10 сидов $\times$ 8 эпизодов каждое, $\times$ 2 режима ассиста. Семантический стек (варианты node, concept, state-conditioned) сравнивается с четырьмя бейзлайнами: random, graph-only, SPTM-like, обученный BC. Отчётные величины: сквозной успех, шаги, четыре операциональные постадийные частоты, дельты ассиста, слабейшая стадия по задаче и по методу. Все средние агрегированы по сидам; 95\%-е бутстреп-ДИ используют $B{=}5{,}000$ ресемплов~\citep{efron1979bootstrap}; сравнения ассиста используют парные знако-ранговые критерии Уилкоксона~\citep{wilcoxon1945individual} с поправкой Холма--Бонферрони~\citep{holm1979simple} (Приложение~\ref{p1:app:assist}). Для двух самых трудных задач (FourRooms, LavaGapS7) выполняются четыре оракульных режима (\emph{base}, \emph{oracle R/M/C}) --- каждый заменяет ровно одну стадию, оставляя остальной конвейер неизменным. Проверка переносимости на BabyAI --- в Приложении~\ref{p1:app:babyai}. Все одноагентные эксперименты выполнялись на 10-ядерном ARM CPU, 32\,ГБ RAM (только CPU); оракульный свип --- $<1$\,мин настенного времени.

\paragraph{Диагностический блок LLM.}
Чтобы проверить, что Q/R/M/C не привязан к реализации символьного планировщика, мы также запускаем фиксированный субстрат TextNav с LLM-управляемым циклом observe/query/retrieve/decide. Задача TextNav, промпты, контракт парсера и определения событий фиксированы, меняется только локальная модель Ollama; бэкенд использует совместимые с моделью настройки completion, когда chat-шаблон Ollama иначе прячет короткие структурированные ответы в отдельном поле размышлений. Этот блок --- контролируемое исследование бюджета шагов (Таблица~\ref{p1:tab:llm-smoke}): свип бюджета завершения загоняет обоих LLM-агентов в отказ, который Q/R/M/C локализует на стадии $C$, тогда как $Q/R/M$ остаются насыщенными, демонстрируя диагностическое свойство фреймворка на реальных языковых агентах. Воспроизведение полных агентов в стиле ReAct/Reflexion/Voyager/MemGPT --- отдельный системный вклад; сравнение с этими системами --- сравнение покрытия в Приложении~\ref{p1:app:sota-coverage}.

\subsection{Валидация постадийных оценок}
\label{p1:sec:validation}

Перед основным бенчмарком валидируется факторизация. На
синтетическом четырёхстадийном процессе с известными вероятностями
произведение вложенных оценок отслеживает эмпирическое завершение семантического пути
по размерам выборки. На реальном бенчмарке то же произведение совпадает с
завершением семантического пути, а не с общим успехом
(Таблица~\ref{p1:tab:consistency}) --- в точности контракт
Определения~1.

\begin{table*}[t]
  \caption{Аудит согласованности постадийных факторов для лучшего семантического метода в каждом семействе задач. Произведение вложенных условных оценок совпадает с завершением семантического пути; общий успех может быть выше, поскольку некоторые эпизоды успешны без прохождения полного семантического пути извлечения.}
  \label{p1:tab:consistency}
  \centering
  \small
  \begin{tabular}{p{1.7cm}p{2.5cm}c c c c c c c}
    \toprule
    Задача & Метод & Общий & SP-заверш. & $\hat P(Q^\star)$ & $\hat P(R^\star|Q^\star)$ & $\hat P(M^\star|R^\star)$ & $\hat P(C^\star|M^\star)$ & Произв. \\
    \midrule
    Вода & \texttt{water-cp} & 0.95 & 0.70 & 1.00 & 0.98 & 0.74 & 0.97 & 0.70 \\
    Отдых & \texttt{rest-s2} & 0.93 & 0.73 & 1.00 & 1.00 & 0.80 & 0.91 & 0.73 \\
    Целевой регион & \texttt{goal-n1} & 0.54 & 0.00 & 1.00 & 0.54 & 0.01 & 0.00 & 0.00 \\
    Восст. после опасности & \texttt{hazard-s7} & 0.40 & 0.16 & 1.00 & 0.56 & 0.58 & 0.50 & 0.16 \\
    \bottomrule
  \end{tabular}
\end{table*}

\subsection{Сводка «только успех» и что она скрывает}
\label{p1:sec:single-results}

Вода и отдых --- сильные сквозные задачи (успех 0.95 и 0.93; Рисунок~\ref{p1:fig:hero}b). Целевой регион достигает успеха 0.54 на уровне семейства задач --- среднее, объединяющее решённую среду (Empty-6x6, $1.00$) с трудной (FourRooms, $0.075$; \S\ref{p1:sec:oracle}), так что число по семейству бимодально, а не равномерно средне-трудно, --- но бутылочное горлышко сидит ниже по потоку от выбора кандидатов (\S\ref{p1:sec:oracle}). Восстановление после опасности --- самая трудная чувствительная к извлечению задача с успехом 0.40, где слабейшая стадия --- удовлетворение извлечения. При оценивании «только успех» два трудных семейства выглядят одинаковым видом слабости; постадийный взгляд покажет, что они противоположны. Позадачные операциональные частоты с 95\%-ми ДИ --- в Таблице~\ref{p1:tab:main-results} (Приложение~\ref{p1:app:single-results-table}).

\paragraph{Сравнение с бейзлайнами.}
Таблица~\ref{p1:tab:baselines} сравнивает семантический стек с бейзлайнами random,
graph-only, SPTM-like и обученным поведенчески клонированным. Обученный
бейзлайн доминирует на самых лёгких ресурсных задачах (1.00 на воде
и отдыхе) и восстановлении после опасности (0.74), тогда как семантический стек остаётся
сильнейшим на насыщенном отношениями семействе целевого региона (0.54).
Однако релевантный для этой статьи контраст --- не столбец успеха:
семантический стек --- единственное семейство методов, чьи решения обнажают
интерпретируемую структуру Q/R/M/C. У BC-бейзлайна стадийные события
вырождены по той же структурной причине, что и у многоагентного
обученного бейзлайна, анализируемого в \S\ref{p1:sec:multi-agent-results}: у
политики без явного интерфейса запроса/фиксации нет возможности выпускать разделимые
стадийные события, так что диагностической декомпозиции не к чему
прикрепиться. Сквозное обучение и постадийная диагностика отвечают на разные
вопросы; эта статья --- о втором.

\begin{table*}[t]
  \caption{Сравнение с бейзлайнами. Обученный поведенчески клонированный бейзлайн доминирует на самых лёгких ресурсных задачах и восстановлении после опасности, тогда как семантический стек остаётся сильнейшим на насыщенном отношениями семействе целевого региона и остаётся единственным семейством с интерпретируемой структурой Q/R/M/C.}
  \label{p1:tab:baselines}
  \centering
  \small
  \begin{tabular}{p{1.9cm}c c c c c}
    \toprule
    Задача & Random & Graph-only & SPTM-like & Обученный & Лучший семантич. \\
    \midrule
    Вода & 0.85 & 0.90 & 0.90 & 1.00 & 0.95 \\
    Отдых & 0.85 & 0.90 & 0.90 & 1.00 & 0.93 \\
    Целевой регион & 0.41 & 0.49 & 0.52 & 0.50 & 0.54 \\
    Восст. после опасности & 0.08 & 0.39 & 0.00 & 0.74 & 0.40 \\
    \bottomrule
  \end{tabular}
\end{table*}

\subsection{Оракульные постадийные интервенции локализуют бутылочное горлышко}
\label{p1:sec:oracle}

На целевом регионе семантические варианты решают Empty-6x6 (успех $1.00$), но проваливаются на FourRooms ($0.075$). Для восстановления после опасности (LavaGapS7) базовый успех $0.39$ поднимается лишь до $0.43$ при оракульном контроллере, но подскакивает до $0.59$ при оракульной материализации и $0.60$ при оракульном извлечении; возврат в целевой регион ограничен нижележащими стадиями (оракульное R поднимает завершение семантического пути до $0.08$, не сдвигая сквозной успех). Позадачная атрибуция в виде каскадной диаграммы --- на Рисунке~\ref{p1:fig:oracle-waterfall} (Приложение~\ref{p1:app:single-results-table}). Две задачи, идентичные при оценивании «только успех», отказывают на противоположных концах цепочки Q/R/M/C.

\subsection{Логи уточняют диагноз: предел накопления памяти}
\label{p1:sec:refined-attribution}

Инспекция собственных логов фреймворка заостряет атрибуцию извлечения. На 10-сидовом прогоне восстановления после опасности залогировано $349$ вызовов извлечения: лишь $32$ ($9.2\%$) вернули хотя бы одного кандидата, и из них $31$ (точность $96.9\%$) выбрал семантически удовлетворяющего кандидата. Бутылочное горлышко, следовательно, не в \emph{ранжировании} кандидатов, а в их \emph{генерации} --- символьная память не содержит узлов с тегами восстановления после опасности в 91\% точек принятия решений. Переобученный двухслойный MLP генерации кандидатов, обученный на попозиционных признаках (тестовая точность $0.59$ против мажоритарного класса $0.81$; Приложение~\ref{p1:app:learned-baselines}), подтверждает, что мгновенные признаки состояния не могут заменить недостающую структуру. Атрибуция фреймворка «ограничено извлечением» тем самым уточняется до предела \emph{накопления памяти}: генерация кандидатов не решается из одних лишь признаков состояния, и связывающим ограничением является то, как свидетельства траекторий консолидируются в персистентную символьную структуру. Эта атрибуция --- не артефакт порога обязательных тегов $\theta$: свип $\theta$ по всему диапазону оставляет долю пустых множеств кандидатов инвариантной (Приложение~\ref{p1:app:theta-robustness}), потому что пустоты --- это отсутствие какого-либо сохранённого кандидата, а не подпороговая фильтрация. Это подводит к многоагентной секции: когда консолидация становится общим процессом между $N$ агентами, она обретает измеримый \emph{темп}.

\section{Много агентов: память как разделяемый ресурс}
\label{p1:sec:multi}

\subsection{Четыре архитектуры памяти и ось каденции}
\label{p1:sec:archs}

Инстанцированы четыре коммуникационные архитектуры; они разделяют
одноагентный субстрат извлечения, но различаются тем, как память
распределена между агентами (Рисунок~\ref{p1:fig:architectures}).

\begin{figure}[t]
  \centering
  \includegraphics[width=\linewidth]{figures/fig_architectures.pdf}
  \caption{Четыре многоагентные архитектуры памяти. \textbf{(a)}~Независимая: приватная память на агента, без обмена. \textbf{(b)}~Общая шина: одна глобальная память, каждый агент читает/пишет каждый тик. \textbf{(c)}~Центральный агрегатор: памяти агентов плюс хаб, сливающий их в консенсусный отчёт. \textbf{(d)}~Широковещательная рассылка peer-to-peer: памяти агентов с попарным обменом снимками каждые $K$ тиков (ось каденции Эмпирического утверждения~1).}
  \label{p1:fig:architectures}
\end{figure}

\begin{description}[leftmargin=10pt,topsep=2pt,itemsep=2pt]
\item[Независимая.] Каждый агент поддерживает приватное хранилище событий и граф понятий. Механизма межагентной коммуникации не существует. Бейзлайн-нижняя граница.
\item[Общая шина.] Все агенты публикуют наблюдения в единую шину событий \emph{коллективной памяти}. Один общий граф понятий перестраивается из объединения событий. Максимально централизовано.
\item[Центральный агрегатор (ConsensusLayer).] У каждого агента свой приватный граф понятий; центральный слой стягивает снимки и вычисляет слияние, взвешенное доверием, порождая единый \emph{консенсусный отчёт}, потребляемый всеми агентами.
\item[Широковещательная рассылка peer-to-peer (каденция $K$).] У каждого агента свой граф понятий и своё слитое представление. Каждые $K$ тиков каждый агент рассылает снимок в пассивную шину; каждый агент независимо сливает свой снимок с последним снимком, полученным от каждого пира. Разные агенты могут держать разные слитые представления; центрального отчёта нет.
\end{description}

Первые три соответствуют трём режимам в литературе по многоагентной
коммуникации: без обмена, агрегация в стиле CTDE и
децентрализованная пировая коммуникация. Четвёртая --- архитектура, к
которой Определение~2 и Эмпирическое утверждение~1 применимы
напрямую, поскольку каденция $K$ --- настраиваемый параметр, а не
фиксированный протокол. Механически $K$ --- это темп, с которым приватные
памяти консолидируются в коллективное представление, поэтому он и есть
структурная ось этой секции.

\subsection{Многоагентная факторизация}
\label{p1:sec:multi-theory}

Расширим факторизацию на популяцию из $N$ агентов, разделяющих
информацию через коммуникационный протокол $\mathcal{C}_K$,
параметризованный каденцией рассылки $K$. Каждый агент $i$ имеет своё
состояние памяти $\mathcal{M}_i$ и выпускает свой запрос $q_i$. Пусть
$\mathcal{M}_i^{(K)}$ обозначает проекцию памяти агента $i$ в
момент запроса, после впитывания каждого пирового сообщения, полученного
при каденции $K$. Разделяются два канала связывания: (i) связывание
памяти $\mathcal{M}_{-i}^{(K)}$ --- последние снимки памяти пиров,
доступные агенту $i$; (ii) связывание занятости $\mathcal{O}_{-i}$ ---
множество целей, которые другие агенты уже заявили к моменту,
когда действует $i$.

\paragraph{Определение 2 (многоагентная оцениваемость стадий).}
Для каждого агента $i$
\[
P(C^\star_i \mid q_i, \mathcal{M}_i, \mathcal{C}_K)
=
P(Q^\star_i \mid q_i, \mathcal{M}_i)\;
P(R^\star_i \mid Q^\star_i, q_i, \mathcal{M}_i, \mathcal{M}_{-i}^{(K)})
\]
\[
\quad \cdot\;
P(M^\star_i \mid R^\star_i, q_i, \mathcal{M}_i, \mathcal{M}_{-i}^{(K)})\;
P(C^\star_i \mid M^\star_i, q_i, \mathcal{O}_{-i}).
\]
При многоагентном постадийно-марковском предположении
(Приложение~\ref{p1:app:estimators}) четыре фактора оцениваемы из
поагентных логов траекторий, а эмпирико-частотные оценки
состоятельны.

\paragraph{Когда Определение~2 может нарушаться.}
Одноагентное постадийно-марковское свойство впитывает историю траектории в
$(q, \mathcal{M})$. В многоагентной постановке аналогичное
впитывание требует, чтобы зависимость от прошлых решений пиров входила только
через $\mathcal{M}_{-i}^{(K)}$ на стадии M и через
$\mathcal{O}_{-i}$ на стадии C. Этот шаг нетривиален:
$\mathcal{O}_{-i}$ суммирует занятость целей в момент решения, но
теряет информацию о том, \emph{как} пиры туда пришли.
Факторизация, следовательно, --- проекция на наблюдаемые поагентные
трассы при такой сводке. Возникают две конкретные моды нарушения: (i)
межагентные коммуникационные побочные каналы, не логируемые через
$\mathcal{C}_K$; (ii) координированные политики пиров, чьё решение на
стадии C зависит от скрытого состояния пиров, не суммируемого
$\mathcal{O}_{-i}$.

\subsection{Смещение бутылочного горлышка}
\label{p1:sec:shift}

\paragraph{Эмпирическое утверждение~1 (смещение бутылочного горлышка M$\leftrightarrow$C, на уровне наклонов).}
Предположим дефицит ($N > T$ целей, см. \S\ref{p1:sec:notation}) и
архитектуру пировой рассылки с каденцией $K \in \mathbb{N}$. Постулируются два
стилизованных предположения:
\textbf{(F)}~свежесть памяти $\nu(\mathcal{M}_{-i}^{(K)})$
поточечно не убывает при уменьшении $K$, и вероятность того, что
планировщик агента $i$ зафиксируется на достижимой цели, не убывает по
$\nu$;
\textbf{(O)}~вероятность того, что произвольная закоммиченная цель
уже находится в $\mathcal{O}_{-i}$, не возрастает при росте $K$.
Обозначим через $P(M^\star|R^\star;K)$ и $P(C^\star|M^\star;K)$
популяционные условные вероятности (усреднённые по агентам и
распределению сидов бенчмарка); $\hat{P}(\cdot)$ обозначает
соответствующую эмпирико-частотную оценку, вычисленную из логов. При
(F) и (O)
\[
\frac{\partial\, P(M^\star|R^\star;K)}{\partial K^{-1}} > 0,
\qquad
\frac{\partial\, P(C^\star|M^\star;K)}{\partial K^{-1}} < 0.
\]
Оба условия на наклоны --- прямые эмпирические предсказания,
проверяемые на логах через $\hat{P}$.

\paragraph{Чем это утверждение является и чем не является.}
Эмпирическое утверждение~1 --- \emph{утверждение об условных частотах на уровне наклонов}: каждый из двух наклонов проверяем и подтверждён на свипе из 35{,}640 прогонов (\S\ref{p1:sec:multi-agent-results}) и независимо на свипе переносимости из 100 прогонов на стандартном MiniGrid (\S\ref{p1:sec:limitations}). Мы сообщаем наклоны как \emph{размеры эффекта с кластерно-робастными доверительными интервалами} (ресемплированными по 81 независимой конфигурационной ячейке, а не по коррелированному числу прогонов) в Приложении~\ref{p1:app:effect-sizes}. Граница утверждения относительно условного произведения $\Phi$ сформулирована один раз, в боксе «Границы применимости и не-утверждения» (\S1) и в \S\ref{p1:sec:multi-agent-results}. Смещение --- сигнатура соперничества за дефицитные цели, а не устройства сетки: при изобилии ($T{>}N$) оба наклона выполаживаются, и антикорреляция M$\leftrightarrow$C ослабевает (Приложение~\ref{p1:app:cadence-robustness}), т.е.\ предсказание ломается ровно там, где устранён механизм соперничества.

\paragraph{Хронология.}
Эмпирическое утверждение~1 является post-hoc. Оно было сформулировано \emph{после} первоначального свипа из 12{,}960 прогонов, проведённого при четырёх безразличных к стадиям гипотезах (Приложение~\ref{p1:app:diagnostics}); затем два условия на наклоны были подтверждены на полном свипе из 35{,}640 прогонов, на масштабировании $N\!=\!12$ (Приложение~\ref{p1:app:babyai}) и на обёртке переносимости MiniGrid (\S\ref{p1:sec:limitations}). Мы заявляем это явно: подтверждённые знаки наклонов --- предсказания на укрупнённых данных, а не предварительно зарегистрированное
предсказание. Полностью формальная версия оставлена для будущей работы.

\subsection{Протокол свипа}
\label{p1:sec:multi-agent-protocol}

Четыре архитектуры оцениваются на собственной многоагентной сетке ($12{\times}10$ клеток); каждый агент должен достичь клетки с тегом \texttt{water\_source}, выбирая действия через общий планировщик на основе памяти, так что различается только слой памяти. Свип скрещивает $N{\in}\{3,5,8\}$, $T{\le}N$, три раскладки (симметричная/асимметричная/случайная), три плотности опасностей, четыре архитектуры, восемь пировых каденций $K{\in}\{1,2,4,8,16,32,48,64\}$, 40 сидов --- 35{,}640 прогонов за $\approx 22$ мин на 8 ядрах. Логируются поагентные потиковые индикаторы Q/R/M/C (Приложение~\ref{p1:app:estimators}); «звёздные» события уровня эпизода --- дизъюнкция по тикам. Отдельный блок оракульных интервенций на случаях дефицита $(N,T){\in}\{(3,2),(5,3),(8,5)\}$ при пяти каденциях оценивает оракульное R (память возвращает истинные позиции воды), оракульное M (планировщик фиксируется на ближайшей реальной воде) и оракульное C (завершение реальной материализованной фиксации гарантировано; Приложение~\ref{p1:app:oracle-c-def}).

\subsection{Результаты: постадийная декомпозиция и смещение бутылочного горлышка}
\label{p1:sec:multi-agent-results}

Свип из 35{,}640 прогонов излагается через четыре наблюдения,
структурированных зеркально Эмпирическому утверждению~1: (i) условная
декомпозиция изолирует стадии M и C как различающие
звенья; (ii) смещение бутылочного горлышка проявляется напрямую как отрицательная
внутрикаденционная корреляция между $P(M^\star|R^\star)$ и
$P(C^\star|M^\star)$, с пиком при $K\!=\!4$ и затуханием до шума при
$K\!=\!64$; (iii) на среднем времени до успеха результирующий компромисс
порождает внутренний оптимум при $K\!=\!8$ --- устойчивый по положению, но мелкий по запасу, с преимуществом над соседними быстрыми каденциями, не разрешённым при этом размере выборки (Приложение~\ref{p1:app:effect-sizes}); (iv) на
условном произведении $\Phi$ тот же компромисс порождает монотонный
тренд, а не внутренний пик --- сводная статистика, о которой мы не делаем формальных утверждений. Бейзлайн с обученной коммуникацией и многоагентный
цикл оракульных интервенций закрывают подсекцию.

\subsubsection{Условная декомпозиция показывает, где отказывает каждая архитектура.}
Рисунок~\ref{p1:fig:conditional-rates} показывает поархитектурные условные частоты Q/R/M/C; полная числовая картина с бутстреп-ДИ --- в Приложении~\ref{p1:app:phi-detail}, Таблица~\ref{p1:tab:multi-agent-cond}. Условная частота извлечения $P(R^\star|Q^\star)$ насыщена; различающий сигнал живёт на M и C. Общая шина, центральный агрегатор и быстрая пировая рассылка ($K\le4$) имеют высокое $P(M^\star|R^\star)$, но низкое $P(C^\star|M^\star)$; независимая и медленная пировая рассылка ($K\ge16$) показывают обратное --- компромисс M$\leftrightarrow$C, предсказанный Эмпирическим утверждением~1.

\begin{figure}[t]
  \centering
  \includegraphics[width=\linewidth]{figures/fig_conditional_rates.pdf}
  \caption{Поархитектурные условные частоты (сверху) и среднее время до успеха (снизу) на свипе из 35{,}640 прогонов. Общая шина, центральный агрегатор и быстрая пировая рассылка ($K \le 4$) живут в режиме насыщения M / коллапса C; независимая и медленная пировая ($K \ge 16$) --- в режиме дефицита M / насыщения C. Среднее время до успеха имеет внутренний минимум при $K{=}8$ ($7.47$ тиков), ниже бейзлайна независимых агентов ($8.02$); положение устойчиво, но мелко --- запас над соседними быстрыми каденциями не разрешён (Приложение~\ref{p1:app:effect-sizes}).}
  \label{p1:fig:conditional-rates}
\end{figure}


\subsubsection{Смещение бутылочного горлышка и внутренний оптимум эффективности.}
По всем пировым прогонам Спирмен $P(M^\star|R^\star)$ от $K$ равен $-0.53$ (95\%-й кластерно-робастный ДИ $[-0.59,-0.47]$), а $P(C^\star|M^\star)$ от $K$ равен $+0.43$ (ДИ $[+0.37,+0.50]$); оба --- умеренные размеры эффекта, чьи ДИ исключают ноль при ресемплировании 81 конфигурационной ячейки (Приложение~\ref{p1:app:effect-sizes}), что соответствует условиям на наклоны Эмпирического утверждения~1. Мы сообщаем именно их, а не $p$-значения, которые неинформативны при таком коррелированном размере выборки. Более сильная эмпирическая сигнатура --- на уровне отдельных конфигураций: Пирсон $P(M^\star|R^\star)$ от $P(C^\star|M^\star)$ \emph{внутри} фиксированного $K$ равен $-0.43, -0.54, -0.61, -0.19, -0.05, -0.04, -0.02, -0.02$ при $K{=}1,2,4,8,16,32,48,64$, с пиком при $K{=}4$ и затуханием до шума к $K{=}64$ (Рисунок~\ref{p1:fig:jensen-scatter}, Приложение~\ref{p1:app:phi-detail}). Конфигурации с высоким M систематически имеют низкое C: бутылочное горлышко мигрирует между двумя стадиями внутри одной и той же ячейки $(N,T,\text{раскладка},\text{опасность},\text{сид})$ при варьировании $K$. На метрике эффективности среднее $t_{\mathrm{succ}}$ имеет U-образную форму с внутренним минимумом при $K{=}8$ (7.47 тиков против 8.02 у независимого бейзлайна); положение минимума устойчиво, но мелко --- его запас над соседними быстрыми каденциями ($K{=}1,4$) не разрешён при ресемплировании на уровне ячеек (парно-бутстреповские ДИ в Приложениях~\ref{p1:app:phi-detail} и \ref{p1:app:effect-sizes}). Масштабирование до $N{=}12$ \emph{заостряет} сигнатуру: пиковый Пирсон растёт до $-0.70$ и
мигрирует наружу к $K\!=\!8$, что согласуется с тем, что более сильное связывание
занятости при большем $N$ задерживает переход между режимами издержек
(Приложение~\ref{p1:app:babyai}, Рисунок~\ref{p1:fig:scale-n12}).

\paragraph{Стратификация по провенансу: механизм за смещением.}
Аннотирование каждой фиксации $M^\star$ \emph{провенансом} свидетельств,
стоящих за ней, делает смещение бутылочного горлышка механизменным, а не
корреляционным. Пусть $\varphi$ --- чужая доля массы свидетельств,
поддерживающих зафиксированного кандидата, вычисленная из тех же весов
доверия/поддержки, которые использовало само слияние, и замороженная в момент фиксации; назовём фиксацию
\emph{варпом} (warp, $W^\star$), когда $\varphi \ge 0.8$: агент
коммитится на цель, которую знает по существу только от пиров. На
выделенном пировом свипе по ячейкам дефицита (случайные раскладки, 20
сидов, $K\in\{1,2,4,8,16\}$) доля варпов $P(W^\star|M^\star)$
падает с $0.55$ при $K{=}1$ до $0.11$ при $K{=}8$, тогда как страта
собственных свидетельств завершает на ровных $0.26$--$0.34$, а страта варпов
завершает ниже $0.05$ на всём диапазоне: \emph{C-коллапс
быстрого обмена сконцентрирован в страте варпов}. То, чем каденция
управляет механически, --- это какая часть материализации коллектива
работает на чужих свидетельствах, --- и именно эта страта проигрывает
гонку занятости. Та же стратификация
воспроизводится на непрерывном мосте VMAS ($0.72\to0.00$ по
$K{=}1/4/16$) и на текстовом LLM-субстрате ($0.42\to0.08$ по
$K{=}2/8$), причём кривые почти совпадают между субстратами. Полный
анализ, обусловленный провенансом (контрфактическая атрибуция маскированным
реплеем, замкнутая форма горизонта устаревания с зарегистрированным
hold-out $6/6$-exact и децентрализованный протокол починки), приведён в
сопутствующей рукописи; здесь стратификация служит
механизменным прочтением Эмпирического утверждения~1.

\paragraph{Чего мы не утверждаем о $\Phi$.} Условное произведение $\Phi = P(M^\star|R^\star)\!\cdot\!P(C^\star|M^\star)$ \emph{не} имеет внутреннего пика в проверенном диапазоне (оценка средним произведений монотонно растёт от $0.569$ при $K{=}4$ до $0.603$ при $K{=}64$, сходясь к границе независимых агентов $0.605$). Утверждение уровня наклонов Эмпирического утверждения~1 --- об условных частотах самих по себе, не о $\Phi$; мы не делаем формальных утверждений о форме $\Phi(K)$. Детали разрыва Йенсена и сравнение POM--MOP --- в Приложении~\ref{p1:app:phi-detail}.

\paragraph{Бейзлайн с обученной коммуникацией: стадийные события структурно вырождены даже при конкурентоспособном успехе.}
Политика в стиле CommNet \citep{sukhbaatar2016commnet}, обученная PPO~\citep{schulman2017proximal} на том же сценарии ($N{=}3$, $T{=}2$), достигает $0.667$ успеха на отложенной выборке ($n{=}100$ сидов) --- вровень с символьной пировой архитектурой при $K{=}8$ ($0.65$). Несмотря на это, профиль Q/R/M/C остаётся структурно вырожденным: $Q^\star{=}1.00$ насыщено по построению; $R^\star{=}0.997$, $M^\star{=}0.997$ схлопнуты вместе ($P(M^\star|R^\star){=}1.00$). Это и есть проверенное структурное утверждение: вещественный коммуникационный вектор никогда не бывает информативно пустым (Q насыщается) и лишён явного интерфейса фиксации цели (R/M сливаются), независимо от бюджета обучения. Многоагентный аналог одноагентного анализа обученного BC (\S\ref{p1:sec:single-results}). Полная архитектура, гиперпараметры PPO и параллельное сравнение REINFORCE и PPO --- в Приложении~\ref{p1:app:learned-baselines}.

\paragraph{Свип LLM-моделей: логгер локализует реальный отказ в языковых агентах.}
Таблица~\ref{p1:tab:llm-smoke} прогоняет тот же контракт событий Q/R/M/C на двух локальных LLM-агентах (qwen3:4b и llama3.1 через Ollama) на TextNav, свипуя бюджет шагов по $\{6,12,25\}$ ($8$ эпизодов каждый; задача требует $\sim$10 низкоуровневых шагов). При щедром бюджете ($12$, $25$) оба агента успешны, и все четыре стадии срабатывают. При бюджете ниже требуемого ($6$) успех падает до $0/8$ у \emph{обеих} моделей --- и декомпозиция показывает, \emph{где}: $Q^\star{=}R^\star{=}M^\star{=}1.00$ (агенты формируют запрос, извлекают и фиксируют правильную цель в каждом эпизоде), тогда как $C^\star{=}0.00$ (ни один не завершает вовремя). Фреймворк локализует отказ точно в ограниченном бюджетом завершении ($C$), а не в формировании запроса, извлечении или материализации; то, что обе модели отказывают одинаково, показывает структурность отказа, а не артефакт модели. Это то самое диагностическое свойство, которое фреймворк заявляет, теперь продемонстрированное на \emph{реальных языковых агентах}: в остальном непрозрачный коллапс успеха ($1.00\to0.00$) разрешается в конкретную отказывающую стадию. (Замечание об инструментировании: наивный вызов chat-шаблона qwen3:4b/Ollama изначально возвращал пустые строки \texttt{response}, потому что модель тратила короткий бюджет структурированного вывода в отдельном поле размышлений; для qwen3:4b мы используем сырой structured completion и парсер первой строки, что зафиксировано для воспроизводимости.)

\paragraph{Многоагентная феноменология воспроизводится на живом LLM-коллективе.}
Помимо одноагентного исследования бюджета, тот же контракт событий
инструментирует коллектив из $N{=}3$ LLM-агентов (llama3.1), разделяющих
символьную память мест пировой рассылкой при дефиците ($M{=}2$
доступные для захвата цели, четыре раскладки, $K\in\{2,8\}$), причём все предсказания
зарегистрированы до прогонов. Каденционная феноменология из
\S\ref{p1:sec:multi-agent-results} переносится на языковой
субстрат: парное сравнение с бейзлайном без коммуникации
на идентичных раскладках показывает, что сырой обмен \emph{снижает} успех с
$2/3$ до $1/3$ на эпизод --- C-коллапс гонки за занятость, теперь
порождённый агентами, чьи запросы, извлечение тегов и коммиты
сделаны LLM, --- а столкновения целей появляются в $7/8$ эпизодов с общей
памятью. Присоединение протокола починки «резервирование и откат»
(конструктивный ответ на открытую проблему~3 из
Приложения~\ref{p1:app:open-problems}) возвращает часть потери
($0.333\to0.417$ при обеих каденциях, латентность отката снова
${\approx}K$). Два интерфейсных наблюдения переносятся на любой
дизайн LLM-поверх-символьной-памяти: LLM-формирователи запросов свободно генерируют
теги извлечения вне словаря памяти (что требует толерантных
гейтов извлечения), а малые локальные модели надёжны в семантическом
коммите, но не в координатной арифметике или систематическом исследовании
(что требует слоёного исполнения). Полный протокол и анализ провенанса ---
в сопутствующей рукописи.

Границы: это контролируемые исследования на локальных моделях; масштабирование того же контракта на полноценных агентов в стиле Voyager/MemGPT на более богатых задачах с несколькими модами отказа развивается далее (валидация на внешних фреймворках) и в Приложении~\ref{p1:app:sota-coverage}.

\begin{table}[h]
  \caption{Свип бюджета шагов LLM TextNav: Q/R/M/C локализует ограниченный завершением отказ у двух локальных Ollama-агентов. Восемь эпизодов на (модель, лимит шагов), температура 0.0; задача требует $\sim$10 низкоуровневых шагов. При бюджете ниже этого требования (лимит шагов 6) успех падает до $0.00$ у обеих моделей, однако $Q^\star{=}R^\star{=}M^\star{=}1.00$ (агенты формируют запрос, извлекают и фиксируют правильную цель в каждом эпизоде), тогда как $C^\star{=}0.00$: фреймворк локализует отказ на стадии C, а не на Q/R/M. Обе модели ведут себя идентично, так что отказ структурен (бюджет), а не специфичен для модели. Артефакты: \texttt{tmp/llm\_steplimit\_\{6,12,25\}}.}
  \label{p1:tab:llm-smoke}
  \centering
  \small
  \begin{tabular}{l c c c c c c}
    \toprule
    Модель & Лимит шагов & Успех & $Q^\star$ & $R^\star$ & $M^\star$ & $C^\star$ \\
    \midrule
    qwen3:4b & 6  & 0.00 & 1.00 & 1.00 & 1.00 & \textbf{0.00} \\
    llama3.1 & 6  & 0.00 & 1.00 & 1.00 & 1.00 & \textbf{0.00} \\
    \midrule
    qwen3:4b & 12 & 1.00 & 1.00 & 1.00 & 1.00 & 1.00 \\
    llama3.1 & 12 & 1.00 & 1.00 & 1.00 & 1.00 & 1.00 \\
    qwen3:4b & 25 & 1.00 & 1.00 & 1.00 & 1.00 & 1.00 \\
    llama3.1 & 25 & 1.00 & 1.00 & 1.00 & 1.00 & 1.00 \\
    \bottomrule
  \end{tabular}
\end{table}

\paragraph{Многоагентные оракульные интервенции замыкают цикл диагностика$\to$интервенция$\to$верификация.}
На случаях дефицита среднее $t_{\mathrm{succ}}$ при малой/промежуточной каденции ($K{\in}\{1,2,4\}$) падает с $\sim10$--$11$ в бейзлайне до $\sim7.2$ при оракульном R или оракульном M и далее до $\sim6.1$ при оракульном C, которое дополнительно устраняет путь после фиксации; при медленной каденции ($K{=}16$) бейзлайн уже в пределах $\sim1.2$ от пола R/M-оракулов (Таблица~\ref{p1:tab:multi-agent-oracle} в Приложении~\ref{p1:app:multi-oracle}). Существенно, что оракульное C даёт наибольшее сокращение времени, но наименьший прирост доли успеха (Приложение~\ref{p1:app:oracle-c-def}): при дефиците завершение стоит времени, тогда как извлечение и материализация решают успех. Индуцированное каденцией избыточное время при малых/промежуточных $K$ атрибутируется переходу M$\leftrightarrow$C; при больших $K$ бутылочное горлышко перемещается в одиночное исследование.

Многоагентный блок тем самым доставляет свою половину аргумента.
Каденция $K$ --- механически темп, с которым приватные памяти
консолидируются в коллективное представление; данные показывают, что у этого темпа есть
нетривиальный внутренний оптимум ($K\!=\!8$ на $t_{\mathrm{succ}}$), что
его канал издержек --- связывание занятости, и что оба эффекта
видимы постадийно на той же цепочке Q/R/M/C, что и для одиночных
агентов. Консолидация --- не только одноагентный предел; в
коллективе это настраиваемая, измеримая ось дизайна.

\section{Контракт на сторонних агентных стеках}
\label{p1:sec:external}

Самый сильный вопрос о внешней валидности, который читатель может задать
диагностическому протоколу, --- работает ли он на системах, которые его авторы не
строили. Поэтому мы инструментируем три сторонних агентных стека --- цикл
function-calling на OpenAI SDK, преднастроенный ReAct-агент
LangGraph и пару AssistantAgent/UserProxy из AutoGen, все на одной
локальной модели (llama3.1, температура $0$) --- идентичным
контрактом Q/R/M/C, определённым однократно над \emph{трассами инструментов}: эпизодическая
бытовая память открывается через инструменты \method{recall}/\method{goto}/%
\method{take}, и $Q^\star$ срабатывает при запросе к памяти, $R^\star$ ---
при \emph{атрибуции} (возвращённая запись называет предмет в его истинном
месте), $M^\star$ --- при фиксации уровня эпизода (\method{goto} к
истинному месту, в верности Определению~1; дисциплина первого коммита (first-commitment discipline)
$M_1$ логируется как вспомогательная диагностика качества фиксации), а
$C^\star$ --- при завершении задачи в пределах бюджета инструментов. Три варианта
инженерят по одному отказу каждый --- консолидационный разрыв (consolidation gap; факт о местоположении
никогда не был записан в память), неоднозначность с устаревшим дистрактором, бюджет
инструментов ниже требований задачи --- плюс чистый контроль; все
предсказания были зарегистрированы до прогонов, а две ревизии
семантики событий (утечка атрибуции; перетянутое $M^\star$)
зафиксированы в регистрационном файле.

\begin{table}[h]
\caption{Q/R/M/C на трёх сторонних агентных фреймворках (llama3.1,
20 эпизодов на ячейку, идентичные эпизоды/инструменты/бюджеты). Структурные
отказы локализуются везде идентично: консолидационный разрыв ---
$R{=}0.00$, а бюджетный отказ --- $C{=}0.00$ при $Q{=}R{=}1$ для
всех трёх стеков. На \emph{чистом} контроле фреймворки расходятся
от $0.50$ до $0.75$ сквозным образом --- и декомпозиция атрибутирует
этот разброс (см. текст). $M_1$: первый \method{goto} целится в истинное
место; stale$_1$: первый \method{goto} целится в устаревший дистрактор.}
\label{p1:tab:external}
\centering
\small
\begin{tabular}{llcccccc}
\toprule
Фреймворк & Вариант & $Q^\star$ & $R^\star$ & $M^\star$ & $C^\star$ & $M_1$ & stale$_1$ \\
\midrule
OpenAI SDK & контроль     & 1.00 & 1.00 & 0.85 & 0.75 & 0.25 & --- \\
           & консолидац. разрыв & 1.00 & \textbf{0.00} & 0.50 & 0.40 & 0.25 & --- \\
           & устаревшая неоднозн.   & 1.00 & 1.00 & 0.90 & 0.30 & 0.25 & 0.20 \\
           & жёсткий бюджет      & 1.00 & 1.00 & 0.25 & \textbf{0.00} & 0.25 & --- \\
\midrule
LangGraph  & контроль     & 1.00 & 1.00 & 0.70 & 0.50 & 0.50 & --- \\
           & консолидац. разрыв & 1.00 & \textbf{0.00} & 0.20 & 0.15 & 0.10 & --- \\
           & устаревшая неоднозн.   & 1.00 & 1.00 & 0.65 & 0.60 & 0.50 & 0.35 \\
           & жёсткий бюджет      & 1.00 & 1.00 & 0.10 & \textbf{0.00} & 0.10 & --- \\
\midrule
AutoGen    & контроль     & 1.00 & 1.00 & 0.85 & 0.55 & 0.45 & --- \\
           & консолидац. разрыв & 1.00 & \textbf{0.00} & 0.30 & 0.20 & 0.15 & --- \\
           & устаревшая неоднозн.   & 1.00 & 1.00 & 0.50 & 0.45 & 0.35 & 0.40 \\
           & жёсткий бюджет      & 1.00 & 1.00 & 0.15 & \textbf{0.00} & 0.15 & --- \\
\bottomrule
\end{tabular}
\end{table}

\begin{figure}[t]
\centering
\includegraphics[width=\linewidth]{figures/fig_external_stages}
\caption{Постадийные профили трёх сторонних агентных стеков на четырёх
инженерных вариантах (llama3.1, $n{=}20$ на ячейку). Запрос и извлечение
насыщаются везде идентично; два структурных отказа дают
точные нулевые столбцы на каждом стеке ($R^\star{=}0$ при
консолидационном разрыве, $C^\star{=}0$ при жёстком бюджете);
поведенческий разброс между фреймворками целиком живёт на стадиях M и
C.}
\label{p1:fig:external-stages}
\end{figure}

Три наблюдения (Таблица~\ref{p1:tab:external},
Рисунок~\ref{p1:fig:external-stages}). \textbf{(1)~Структурные
локализации точны и независимы от фреймворка.} Когда в памяти
нет факта, $R^\star{=}0.00$ на каждом стеке; когда бюджет
не вмещает завершение, $C^\star{=}0.00$ при $Q^\star{=}R^\star{=}1$ на
каждом стеке ($20/20$ эпизодов в каждом). Контракт не требует
пофреймворковой адаптации. \textbf{(2)~Чистый контроль разделяет
фреймворки, и декомпозиция говорит как.} Сквозной успех на
ОДНИХ И ТЕХ ЖЕ чистых эпизодах лежит в диапазоне от $0.50$ (LangGraph) до $0.75$ (OpenAI SDK).
Лидерборд сообщил бы «разное качество»; стадии показывают
общую патологию плюс разную экономику циклов: дисциплина первого
коммита слаба везде ($M_1{=}0.25$--$0.50$ --- первый
\method{goto} агента часто игнорирует его же только что извлечённые
свидетельства), материализация уровня эпизода восстанавливается до $0.70$--$0.85$
через циклы переспроса, а затем фреймворки различаются тем, какая часть
этого восстановления доживает до завершения ($C^\star/M^\star$ от $0.88$
у цикла OpenAI SDK до $0.65$ у AutoGen). \textbf{(3)~Устаревший
дистрактор тянет, измеримо и неравномерно.} Устаревшая запись
захватывает первый коммит в $0.20$--$0.40$ неоднозначных эпизодов
(сильнее всего у AutoGen), и цена её для завершения
специфична фреймворку: цикл OpenAI SDK платит тяжелее всех
($C^\star:0.75\to0.30$), тогда как LangGraph и AutoGen поглощают крюк
($0.50\to0.60$; $0.55\to0.45$) --- LLM-памятный аналог
гейтинга по устареванию, решающую роль которого показали коллективные
эксперименты. Честно проваленные регистрации: поведенческие предсказания,
предполагавшие чистый контроль (все стадии $\ge 0.75$; равномерное
индуцированное неоднозначностью падение завершения), сломались \emph{потому что сами фреймворки
проваливают чистую задачу на память в 25--50\% случаев и платят за
неоднозначность в разных местах} --- а это и есть наблюдения (2)--(3), не
шум; все ревизии зафиксированы как зарегистрированные провалы с
механизмами. Границы: три стека при $n{=}20$ на ячейку на одной основной
модели; Letta/CrewAI (Python $\ge$3.10) и модели API-масштаба --- естественное
расширение.

\section{Ограничения и будущая работа}
\label{p1:sec:limitations}

\textbf{Переносимость между субстратами: условия на наклоны выполняются на MiniGrid.} Условия на наклоны Эмпирического утверждения~1 выполняются и на субстрате стандартного MiniGrid. Многоагентная обёртка на основе FourRooms (построенная из \texttt{minigrid.core.grid.Grid} + примитивов Wall/Floor/Lava, идентичные операциональные определения $Q^\star,R^\star,M^\star,C^\star$; \texttt{experiments/minigrid\_multiagent\_wrapper/}) использована для свипа переносимости из 100 прогонов ($K{\in}\{1,4,8,16,64\}$, 20 сидов, $N{=}3$, $T{=}2$): Спирмен $P(M^\star|R^\star)$ от $K$ равен $-0.50$, а Спирмен $P(C^\star|M^\star)$ от $K$ равен $+0.50$ (умеренные ранговые корреляции; это малый свип из 100 прогонов, поэтому мы сообщаем размеры эффекта и не опираемся на значимость). Внутренний минимум $t_{\mathrm{succ}}$ не отделяется при этом минимальном размере свипа (5 каденций, 20 сидов, одна плотность опасностей); более плотный свип --- естественный следующий шаг. Шов между двумя субстратами тем самым закрыт на уровне утверждения о наклонах.

\textbf{Эмпирическое утверждение~1 --- только об условных частотах на уровне наклонов.} Два условия на наклоны выполняются как умеренные, кластерно-робастные размеры эффекта (Приложение~\ref{p1:app:effect-sizes}); мы не делаем формальных утверждений о форме условного произведения $\Phi$ (\S\ref{p1:sec:shift}). Будущая формулировка непосредственно в терминах поверхности временн\'{ы}х издержек или вывод при генеративной модели (пуассоновски-обновляемая память, занятость рождения-гибели) превратили бы эмпирическое утверждение в теорему.

\textbf{Оцениваемость при скрытом состоянии.} При нелогируемых внутриэпизодных переписываниях запросов или пировых побочных каналах Определения~1 и~2 --- проекции на залогированные трассы, а не полная идентификация.

\textbf{Общая система координат.} Все многоагентные результаты предполагают, что агенты
интерпретируют места в общей системе координат: межагентная идентичность
места --- это координатная близость, с семантическими профилями как
разрешителем ничьих внутри пространственного радиуса. Проблема соответствия мест ---
обозначают ли записи двух агентов вообще одно и то же место ---
здесь, следовательно, вынесена за скобки, а не решена. Сопутствующая рукопись снимает
это предположение конструктивно (идентичность мест на основе отпечатков с
консенсусным восстановлением системы координат с точностью до трансляции и поворота на $90^\circ$)
и показывает, что измерения Q/R/M/C, включая каденционные наклоны и
горизонт устаревания, неизменны поверх восстановленных систем координат; более трудные
варианты (непрерывное $SE(2)$, шум внешнего вида, выравнивание
символов) остаются открытыми.

\textbf{Обученные бейзлайны.} Бейзлайн CommNet PPO достигает успеха $0.667$ на сценарии дефицита --- конкурентоспособно с символьным пировым при $K{=}8$ ($0.65$) --- и подтверждает Q-насыщение ($Q^\star{=}1.00$) и коллапс R/M ($P(M^\star|R^\star){=}1.00$) при конкурентоспособном успехе. Структурное предсказание теперь подтверждено, а не постулировано; абляции на архитектурах с явными символьными каналами фиксации цели оставлены для будущей работы.

\textbf{Масштаб.} Многоагентный свип ограничен $N{=}12$ (Приложение~\ref{p1:app:babyai}; сигнатура заостряется, пиковый Пирсон $-0.70$ при $K{=}8$, внутренний минимум $t_{\mathrm{succ}}$ $4.77$). Более плотное масштабирование при $N{\ge}16$ --- будущая работа.
\textbf{Непрерывный субстрат.} Предварительная проверка VMAS~\citep{bettini2022vmas} воспроизводит \emph{оба} знака наклонов Эмпирического утверждения~1 в режиме связывания (Спирмен $-1.00$ и $+1.00$ по $K{\le}16$, 40 сидов; Приложение~\ref{p1:app:vmas}), переиспользуя символьную память через мост дискретизации. Это лишь направленный сигнал --- идеальная ранговая корреляция по пяти монотонным точкам, одна конфигурация, дискретизованная память, --- так что он указывает, что смещение не является артефактом клеток сетки, но полная валидация в непрерывном пространстве состояний с кластерно-робастными ДИ остаётся будущей работой.

\section{Более широкое влияние}

Статья предлагает диагностическое оценивание для агентов на основе памяти в симулированной навигации. Её положительное влияние методологично: сделать отказы на основе памяти более измеримыми между архитектурами и сделать ось каденции распределённой многоагентной памяти настраиваемой по эмпирически обоснованному критерию. Потенциальные негативные влияния косвенны и возникли бы, если бы похожая диагностика использовалась для улучшения агентов в критичных к безопасности постановках (рои дронов, автономные транспортные средства, распределённые сенсорные сети) без адекватных мер защиты. Настоящая работа --- инфраструктура уровня бенчмарка, а не готовая к развёртыванию система.

\section{Заключение}

Мы предложили Q/R/M/C --- четырёхстадийную декомпозицию навигации на основе
памяти, чьи факторы оцениваемы из стандартных логов траекторий, ---
и использовали её как диагностический инструмент на символьных, обученных,
LLM-управляемых, одноагентных и многоагентных трассах. Главный результат ---
не новая архитектура агента: это воспроизводимый протокол логирования и
оценивания, разделяющий формирование запроса, удовлетворение
извлечения, материализацию цели и завершение после извлечения.
По всем экспериментам этот протокол локализует отказы, которые
оценивание «только успех» не может, обнажает вырожденную наблюдаемость у
неявных обученных политик и делает каденцию обмена памятью измеримой
как постадийная переменная дизайна.

\begin{ack}
Рукопись подготовлена в анонимизированной форме для рецензирования.
Представленные эксперименты выполнены из единого бенчмаркового конвейера,
экспортирующего сводки по семействам задач, сравнения ассиста вкл/выкл, диагностику
методов, CSV многоагентных свипов и бутстреп-анализы из
одних и тех же эталонных бенчмарковых артефактов.
\end{ack}

{\small

}

\FloatBarrier
\begin{center}\rule{0.55\linewidth}{0.4pt}\\[2pt]\textbf{Приложения к этой части}\end{center}


\section{Формальные доказательства}
\label{p1:app:proofs}

\subsection{Аргумент для Определения~1 (одноагентная оцениваемость)}

Для каждого эпизода $i \in \{1,\dots,n\}$ пусть $Q_i^\star, R_i^\star, M_i^\star, C_i^\star$ обозначают вложенные индикаторы уровня эпизода. По построению это измеримые функции лога траектории времени исполнения, удовлетворяющие $C_i^\star \Rightarrow M_i^\star \Rightarrow R_i^\star \Rightarrow Q_i^\star$. Эмпирические оценки
\[
\hat p_Q = \tfrac{1}{n}\sum_i Q_i^\star,\quad
\hat p_{R|Q} = \tfrac{\sum_i R_i^\star}{\sum_i Q_i^\star},\quad
\hat p_{M|R} = \tfrac{\sum_i M_i^\star}{\sum_i R_i^\star},\quad
\hat p_{C|M} = \tfrac{\sum_i C_i^\star}{\sum_i M_i^\star}
\]
состоятельны в силу цепного правила плюс условного постадийно-марковского свойства, и эмпирическое произведение сходится к завершению семантического пути. \hfill$\square$

\subsection{Аргумент для Определения~2 (многоагентная оцениваемость)}

Множества обусловливания укрупняются на $\mathcal{M}_{-i}^{(K)}$ и $\mathcal{O}_{-i}$. При укрупнённом постадийно-марковском свойстве доказательство Определения~1 переносится дословно, и каждая укрупнённая условная вероятность --- бернуллиевская вероятность на измеримом событии. Оговорки, сформулированные в \S\ref{p1:sec:multi-theory} («Когда Определение~2 может нарушаться»), описывают две конкретные ситуации нарушения укрупнённого постадийно-марковского свойства. \hfill$\square$

\subsection{Аргумент для Эмпирического утверждения~1 (формулировка на уровне наклонов)}

Эмпирическое утверждение~1 --- эмпирическое утверждение, а не теорема. Аргумент ниже --- качественное обоснование, мотивирующее два знака наклонов; верифицируют их данные. Мы намеренно не называем это доказательством.

\emph{Наклон (a), при (F).} При уменьшении $K$ $\nu(\mathcal{M}_{-i}^{(K)})$ поточечно не убывает. При условии успешного извлечения вероятность фиксации на достижимой цели не убывает по $\nu$. Следовательно, $\partial P(M^\star|R^\star;K) / \partial K^{-1} > 0$.

\emph{Наклон (b), при (O).} При уменьшении $K$ пиры получают обновлённую память одновременно и сходятся на одной и той же ближайшей цели; маргинальная занятость $\Pr[\text{цель}\in\mathcal{O}_{-i}]$ растёт. Следовательно, $\partial P(C^\star|M^\star;K) / \partial K^{-1} < 0$.

\emph{Замечание об условном произведении $\Phi$.} Эмпирическое утверждение~1 не делает никаких заявлений о форме $\Phi(K) = P(M^\star|R^\star;K)\!\cdot\!P(C^\star|M^\star;K)$. При (F) и (O) наклон $\Phi$ по $K^{-1}$ есть сумма двух противонаправленных слагаемых, и меняет ли она знак --- зависит от количественных величин, не фиксируемых (F) и (O) самими по себе. В наших данных $\Phi$ монотонно по $K$ на $\{4,8,16,32,48,64\}$ и сходится к границе независимых агентов $\Phi_\infty = 0.605$. Практически значимая поверхность эффективности --- среднее время до успеха --- внутренний минимум при $K=8$ имеет; см. \S\ref{p1:sec:multi-agent-results}. \hfill$\square$

\subsection{Связь Галуа запрос--извлечение: формулировка, аргумент, границы}
\label{p1:app:galois}

\paragraph{Факт (стандартный).}
Пусть $G$ и $\Sigma$ --- множества, а $I \subseteq G \times \Sigma$ --- отношение. Для $B \subseteq \Sigma$ определим $\mathrm{ext}(B) = \{p \in G : \forall \tau \in B,\ (p,\tau) \in I\}$, а для $A \subseteq G$ определим $\mathrm{int}(A) = \{\tau \in \Sigma : \forall p \in A,\ (p,\tau) \in I\}$. Тогда для всех $A \subseteq G$, $B \subseteq \Sigma$:
\[
A \subseteq \mathrm{ext}(B) \iff B \subseteq \mathrm{int}(A);
\]
оба отображения антитонны, обе композиции $\mathrm{ext}\circ\mathrm{int}$ и $\mathrm{int}\circ\mathrm{ext}$ --- операторы замыкания, а замкнутые пары $(A, B)$ с $A = \mathrm{ext}(B)$ и $B = \mathrm{int}(A)$, упорядоченные включением объёмов, образуют полную решётку --- решётку понятий контекста $(G, \Sigma, I)$ \citep{ganter1999fca}.

\paragraph{Аргумент.}
$A \subseteq \mathrm{ext}(B)$ разворачивается в $\forall p \in A\ \forall \tau \in B:\ (p,\tau) \in I$ --- условие, симметричное по $A$ и $B$, что в точности есть $B \subseteq \mathrm{int}(A)$. Антитонность и свойства замыкания следуют стандартными однострочными аргументами \citep[Ch.~1]{ganter1999fca}. Рассматривая булеаны как тонкие категории (одну --- с обращённым порядком), выписанная эквивалентность говорит ровно то, что $\mathrm{int}$ и $\mathrm{ext}$ образуют сопряжённую пару; связь Галуа --- это сопряжение между частично упорядоченными множествами \citep[Ch.~IV]{maclane1998categories}. \hfill$\square$

\paragraph{Инстанцирование и границы.}
В нашей памяти возьмём в качестве $G$ записи мест, в качестве $\Sigma$ --- алфавит тегов, а $I_\theta = \{(p,\tau) : \text{таблица тегов места приписывает } \tau \text{ месту } p \text{ с уверенностью} \ge \theta\}$ (Приложение~\ref{p1:app:schemas}); интент запроса --- $B = q^{\mathrm{req}}$, а R-доступность --- $\mathrm{ext}(B) \ne \emptyset$. Три намеренных ограничения области. (i)~\emph{Градуированные баллы.} Развёрнутый балл извлечения градуирован (взвешенный лог-балл над обязательными, предпочтительными и штрафными тегами из Приложения~\ref{p1:app:schemas}), так что чёткая связь выше формализует пороговый скелет кандидатства по обязательным тегам; градуированный аналог --- теория нечётких связей Галуа \citep{belohlavek1999fuzzy}, и здесь мы её не используем. (ii)~\emph{Немонотонные обновления.} Обновления Дирихле--мультиномиальные не монотонны по $I_\theta$ буквально: противоречащее свидетельство может удалить инцидентность. На отчётных бенчмарках связывающий отказ --- отсутствие свидетельств, а не противоречие (91\% пустых объёмов, точность 96.9\% при непустых, \S\ref{p1:sec:refined-attribution}), так что прочтение через монотонный рост --- эмпирически релевантный режим, и мы помечаем исключение, а не отмахиваемся от него. (iii)~\emph{Нет структуры для M или C.} Материализация --- satisficing-функция выбора на непустых объёмах, а завершение --- задача управления при занятости; ни то, ни другое не несёт порядковой структуры, и мы не прикрепляем к ним категориальных утверждений. Прочтение зарабатывает своё место тем, что точно называет моду отказа (запрошенный интент с пустым объёмом лежит вне реализованной решётки понятий), даёт консолидации точный объект (рост контекста, а значит, каждого объёма) и структурно локализует ось каденции: связывание памяти $\mathcal{M}^{(K)}_{-i}$ сливает контексты пиров с темпом $K^{-1}$, тогда как занятость $\mathcal{O}_{-i}$ вне-контекстна --- что и есть структурное содержание двух знаков наклонов Эмпирического утверждения~1.

\section{Прочтение стадий через опции / иерархический MDP}
\label{p1:app:hmdp}

Это приложение развивает динамическое прочтение, на которое указано в
\S\ref{p1:sec:event-semantics}. Оно не вводит нового алгоритма: оно именует
объект, который конвейер уже вычисляет, и делает расщепление
$\Phi$-против-$t_{\mathrm{succ}}$ свойством этого объекта. Там, где
теоретико-порядковое прочтение (Приложение~\ref{p1:app:galois}) фиксирует \emph{статическое}
содержание стадий, это адресует время и издержки.

\paragraph{Q/R/M/C как аудит опции.}
Каждый запрос инстанцирует одно темпорально протяжённое действие --- опцию
$o=\langle \mathcal{I},\pi,\beta\rangle$ в смысле
\citet{sutton1999between}: множество инициации $\mathcal{I}$, внутреннюю
политику $\pi$ и условие терминации $\beta$. Стандартный
анализ уровня траекторий трактует опцию как чёрный ящик, оцениваемый её
возвратом. Декомпозиция Q/R/M/C, напротив, --- аудит \emph{внутренней}
структуры опции, и она выстраивается ребро в ребро с
факторизацией $Q^\star\!\rightarrow\!R^\star\!\rightarrow\!M^\star\!\rightarrow\!C^\star$:
\begin{description}[leftmargin=14pt,topsep=2pt,itemsep=2pt]
\item[$Q^\star\!\rightarrow\!R^\star$ (инициация).] Опция «достичь места, удовлетворяющего $q$» корректно инициируется, только когда её множество инициации непусто. В словаре \S\ref{p1:sec:event-semantics} это в точности $\mathrm{ext}(q^{\mathrm{req}})\neq\emptyset$: пустой объём --- не плохой выбор опции, а отсутствие какой-либо допустимой, так что $R$ отказывает до того, как цель вообще зафиксирована.
\item[$R^\star\!\rightarrow\!M^\star$ (политика над опциями).] Материализация --- политика верхнего уровня, коммитящаяся на единственную опцию из допустимого множества --- фиксация цели $g^\star=\chi(\mathrm{ext}(q^{\mathrm{req}}))$ при satisficing-функции выбора $\chi$ из \S\ref{p1:sec:event-semantics}.
\item[$M^\star\!\rightarrow\!C^\star$ (исполнение до терминации).] Завершение --- прокатка внутренней политики $\pi$ до срабатывания условия терминации $\beta$ (прибытие в $g^\star$ или таймаут). Это единственное ребро, чьи моды отказа внесемантичны: компетентность контроллера и, в коллективе, занятость $\mathcal{O}_{-i}$.
\end{description}
Так что Q/R/M/C при этом прочтении --- постадийно разрешающий инструмент для
внутренней структуры генератора опций иерархического MDP, а не
специализированный логгер для одной среды.

\paragraph{Почему оптимум живёт на $t_{\mathrm{succ}}$, а не на $\Phi$: прочтение через полу-MDP.}
В формализме опций / полу-MDP опция оценивается
возвратом, накопленным за \emph{случайную} длительность исполнения $\tau$,
\[
R(s,o)=\mathbb{E}\!\left[\textstyle\sum_{t=0}^{\tau-1}\gamma^{t}\,r_{t+1}\;\middle|\;s_0=s,\,o\right],
\]
где $\tau$ --- число низкоуровневых шагов до срабатывания $\beta$
\citep{sutton1999between}. Цепная вероятность
$\Phi=P(M^\star|R^\star)\!\cdot\!P(C^\star|M^\star)$ фиксирует только,
\emph{успешен ли} аудит опции; она по построению слепа к
$\tau$ --- цене его прохождения. Несущая цену метрика
$t_{\mathrm{succ}}$ --- это и есть $\tau$. В этом структурная причина --- а не
просто эмпирическое совпадение, --- что ресурсно-рациональное прочтение из
\S\ref{p1:sec:event-semantics} \citep{lieder2020resource} находит
внутренний оптимум на $t_{\mathrm{succ}}$ ($K{=}8$), тогда как $\Phi$ остаётся
монотонным: две метрики измеряют разные грани одной и той же опции ---
успех против длительности.

\paragraph{Смещение бутылочного горлышка как многоагентное столкновение опций.}
Динамическое прочтение делает компромисс M$\leftrightarrow$C из
Эмпирического утверждения~1 физически конкретным. Более быстрая рассылка (меньший $K$)
чаще сливает контексты пиров $\mathcal{M}^{(K)}_{-i}$, что может только
укрупнить объём каждого агента и потому повышает вероятность того, что
цель допустима и зафиксирована: $P(M^\star|R^\star)$ растёт. Но
поскольку агенты разделяют пространственный базис и алфавит тегов
(\S\ref{p1:sec:notation}), более свежая общая память ещё и \emph{синхронизирует}
политику верхнего уровня: несколько агентов коммитятся на одну и ту же опцию над
одной и той же дефицитной целью. При занятости $\mathcal{O}_{-i}$ это раздувает
длительность исполнения $\tau$ и подавляет $P(C^\star|M^\star)$. Каденция
$K$, таким образом, регулирует дисперсию распределения опций агентов:
редкий обмен десинхронизирует выбор опций и укорачивает $\tau$,
давая минимум $t_{\mathrm{succ}}$ при внутреннем $K$. Бутылочное горлышко,
следовательно, укоренено в конкуренции между
ростом множества инициации (семантическим, контекстным) и ценой столкновения опций
(внесемантической, занятостной) --- ровно те два ребра, на которые действует каденционный протокол
$\mathcal{C}_K$ на Рисунке~\ref{p1:fig:causal}.

\paragraph{Глосса belief--desire--intention и границы.}
Аудируемое выше состояние верхнего уровня можно назвать в
терминах belief--desire--intention --- в той форме, которую использует
компаньон о коллективной памяти: \emph{belief} --- реализованная решётка
понятий из \S\ref{p1:sec:event-semantics} (что извлекаемо), \emph{desire} ---
предикат, обусловленный состоянием, выбирающий запрос $q$ (например,
переменная жажды, переключающая обязательные теги на \method{water\_source}),
а \emph{intention} --- зафиксированная опция $g^\star$. Иерархическое
состояние тогда есть $h=(\,\text{решётка},\,\text{desire},\,g^\star,\,s\,)$, где
$s$ --- низкоуровневое физическое состояние, и три локуса отказа разделяются как
прежде: пустой объём --- отказ belief/$R$, непустой объём
без устойчивой фиксации --- отказ $M$, а фиксация, недостижимая при
занятости, --- отказ $C$. Две оговорки о границах. Мы не заявляем результата
об \emph{обучении} опций: конвейер --- фиксированный satisficing-генератор опций
\citep{simon1955behavioral}, и Q/R/M/C его аудирует. И мы прикрепляем
структуру опций только к семантическому пути --- $\mathcal{I}$ к
$Q\!\rightarrow\!R$, $(\pi,\beta)$ к $M\!\rightarrow\!C$, --- тогда как
возврат полу-MDP используется как объяснительная линза для расщепления
$\Phi$-против-$t_{\mathrm{succ}}$, а не фитируется к данным. Ни
это опционное прочтение, ни теоретико-порядковое прочтение из
Приложения~\ref{p1:app:galois} не заявляются как вклад: оба --- стандартные
конструкции, используемые как интерпретирующие линзы на измерительный протокол,
который и есть вклад статьи. Их ценность в том, что каждое даёт
проверяемое утверждение на логах --- пустой объём локализует отказ $R$,
а длительность $\tau$ полу-MDP предсказывает, что оптимум
появляется на $t_{\mathrm{succ}}$, а не на $\Phi$.

\section{Операциональные и вложенные постадийные оценки}
\label{p1:app:estimators}

\paragraph{Операциональные диагностические частоты.}
Основной бенчмарк логирует позапросные и поэпизодные диагностические величины напрямую из трасс времени исполнения (запрос сформирован, извлечение удовлетворено, цель материализована, завершение после извлечения). Это операциональные частоты, а не прямые мультипликативные факторы.

\paragraph{Вложенные постадийные оценки для факторизации.}
Для факторизации используются вложенные события уровня эпизода по первой отказавшей стадии, зафиксированной в таксономии времени исполнения: $Q^\star, R^\star, M^\star, C^\star$, вложенные по построению. Аудит согласованности в \S\ref{p1:sec:validation} сообщает $\hat P(Q^\star)$, $\hat P(R^\star|Q^\star)$, $\hat P(M^\star|R^\star)$, $\hat P(C^\star|M^\star)$, их произведение и его зазор до завершения семантического пути.

\subsection{Многоагентные операциональные определения}
\label{p1:app:multi-agent-defs}

Для многоагентного свипа каждая метрика определяется на агента и на тик, затем агрегируется. Пусть $\mathrm{Targets}_i(t)$ --- множество кандидатов, возвращённое запросом к памяти агента $i$ на тике $t$, $\mathrm{Locked}_i(t)$ --- внутрипланировочная зафиксированная цель агента, а $\mathcal{W}$ --- истинное множество целей.

\textbf{Потиковые события.} $Q_i(t) = \mathbf{1}[\mathrm{Targets}_i(t) \ne \emptyset]$; $R_i(t) = \mathbf{1}[\exists w \in \mathcal{W}: \min_{c}\|c - w\| \le \epsilon]$ ($\epsilon = 0.6$); $M_i(t) = \mathbf{1}[\mathrm{Locked}_i(t) \ne \emptyset \land \min_w\|\mathrm{Locked}_i(t) - w\| \le \epsilon]$.

\textbf{Звёздные события уровня эпизода.} $Q^\star_i = \bigvee_t Q_i(t)$; аналогично для $R^\star, M^\star$; $C^\star_i = \mathbf{1}[\text{агент } i \text{ достиг цели}]$.

\textbf{Агрегация.} Частоты по прогону усредняются по $N$ агентам; частоты по конфигурации усредняются по 40 сидам. Условные частоты $P(R^\star|Q^\star)$, $P(M^\star|R^\star)$, $P(C^\star|M^\star)$ вычисляются как отношения поэпизодных счётчиков внутри каждой конфигурации до усреднения по конфигурациям.

\section{Схема памяти и задач}
\label{p1:app:schemas}

\subsection{Полная таблица нотации}
\label{p1:app:notation}

Символы фиксированы на протяжении всей статьи.

\newcommand{\notentry}[1]{\par\smallskip\noindent\textbf{#1}\ \ }

\notentry{Стадии.}$Q$, $R$, $M$, $C$ --- четыре стадии: формирование запроса (\textbf{Q}uery), извлечение (\textbf{R}etrieval), материализация (\textbf{M}aterialization), завершение после извлечения (\textbf{C}ompletion). Они именуют и стадию, и (в выражениях вроде «звено M») соответствующее ребро цепочки.

\notentry{События.}$Q^\star, R^\star, M^\star, C^\star \in \{0, 1\}$ --- звёздные события уровня эпизода, фиксирующие успех каждой стадии (Приложение~\ref{p1:app:estimators}). Вложены по построению: $Q^\star{=}1 \Leftarrow R^\star{=}1 \Leftarrow M^\star{=}1 \Leftarrow C^\star{=}1$.

\notentry{Вероятности.}$P(\cdot)$ --- популяционная вероятность (усреднённая по распределению сидов и, в многоагентном случае, по агентам). $\hat{P}(\cdot)$ --- эмпирико-частотная оценка, вычисленная из логов.

\notentry{Один агент.}$q$ --- запрос; $\mathcal{M}$ --- состояние памяти; $Y$ --- сквозной успех задачи (может превышать $C^\star$).

\notentry{Много агентов.}$N$ --- число агентов в популяции. $T$ --- число целей в среде (ограничение $T \le N$; «дефицит» означает $N > T$). $i$ --- индекс агента. $\mathcal{M}_i$ --- приватная память агента $i$. $q_i$ --- запрос агента $i$.

\notentry{Связывание.}$\mathcal{C}_K$ --- протокол пировой широковещательной коммуникации с каденцией $K \in \mathbb{N}$ (каждые $K$ тиков каждый агент рассылает снимок памяти, \S\ref{p1:sec:archs}). $\mathcal{M}_{-i}^{(K)}$ --- последние снимки памяти пиров, доступные агенту $i$ при $\mathcal{C}_K$ (канал \emph{связывания памяти}). $\mathcal{O}_{-i}$ --- множество целей, уже заявленных другими агентами к моменту, когда действует $i$ (канал \emph{связывания занятости}).

\notentry{Переменные наклонов.}$\nu(\mathcal{M}_{-i}^{(K)})$ --- функционал свежести на памяти пиров; поточечно не убывает при уменьшении $K$ (предположение \textbf{(F)} Эмпирического утверждения~1). Наклон занятости --- предположение \textbf{(O)}.

\notentry{Сводные статистики.}$\Phi(K) := P(M^\star|R^\star;K)\cdot P(C^\star|M^\star;K)$ --- условное произведение, цепная вероятность $P(M^\star \land C^\star \mid R^\star)$. POM (произведение средних) и MOP (среднее произведений) --- две оценки $\Phi$; они различаются на $\mathrm{Cov}(P(M^\star|R^\star), P(C^\star|M^\star))$, что и есть разрыв Йенсена (\S\ref{p1:sec:multi-agent-results}). $t_{\mathrm{succ}}$ --- среднее время до успеха в тиках среды.
\par\smallskip

\subsection{Стек памяти и паттерны задач}

Стек памяти хранит символьные записи мест (уверенности тегов по местам, счётчики, посещаемость, свежесть, геометрия), записи переходов (статистики успеха, риска, издержек с быстрыми/медленными оценками), эпизодические следы (намерение + состояние при каждом посещении) и записи понятий/отношений. На каждом шаге токенизатор выпускает $(z_t, \tau_t, e_t)$ из $o_t$. Таблицы тегов мест используют Дирихле-мультиномиальное обновление; переходы --- экспоненциально сглаженные статистики; эпизодические записи прикрепляют активное намерение и состояние агента. Запрос планировщика отвечается извлечением мест-кандидатов, чьи семантические профили соответствуют текущему намерению задачи и состоянию, через взвешенный лог-балл
$s(p \mid q_t, s_t) = \sum_{\tau \in q^{\mathrm{req}}} \log P(\tau|p)
+ \alpha \sum_{\tau \in q^{\mathrm{pref}}} \log P(\tau|p)
- \beta \sum_{\tau \in q^{\mathrm{pen}}} \log P(\tau|p)
+ \gamma\, u_{\mathrm{epi}}(p, s_t)$
над обязательными, предпочтительными и штрафными тегами. Четыре семейства задач определяются семантическими паттернами мест: поиск воды (цель $\{$water$\}$), безопасный отдых ($\{$rest$\}$ с обусловливанием состоянием), возврат в целевой регион ($\{$goal-region$\}$), восстановление после опасности ($\{$post-hazard-exit$\}$).

\subsection{Одноагентные вспомогательные таблицы и рисунки}
\label{p1:app:single-results-table}

\begin{table*}[!ht]
  \caption{Лучший семантический результат по каждому семейству задач из финального 10-сидового бенчмарка, с 95\%-ми бутстреп-ДИ. Восстановление после опасности упирается в извлечение; целевой регион остаётся ограниченным нижележащими стадиями.}
  \label{p1:tab:main-results}
  \centering
  \scriptsize
  \setlength{\tabcolsep}{3.5pt}
  \begin{tabular}{p{1.5cm}p{1.7cm}c c c c c c c}
    \toprule
    Задача & Лучший метод & Ассист & Успех (95\% ДИ) & Шаги & Запрос (оп.) & Извлеч. (оп.) & Матер.\ (оп.) & Пост (оп.) \\
    \midrule
    Вода & \method{water-cp} & вкл & 0.95 [0.88, 1.00] & 17.99 & 0.71 & 0.71 & 0.71 & 0.70 \\
    Отдых & \method{rest-s2} & вкл & 0.93 [0.85, 0.99] & 23.35 & 0.61 & 0.61 & 0.61 & 0.73 \\
    Целевой регион & \method{goal-n1} & выкл & 0.54 [0.52, 0.56] & 49.34 & 0.00 & 0.00 & 0.00 & 0.00 \\
    Восст. после опасн. & \method{hazard-s7} & вкл & 0.40 [0.35, 0.45] & 10.18 & 0.11 & 0.11 & 0.11 & 0.16 \\
    \bottomrule
  \end{tabular}
\end{table*}

\begin{figure}[!ht]
  \centering
  \includegraphics[width=\linewidth]{figures/fig_oracle_waterfall.pdf}
  \caption{Оракульные постадийные интервенции на двух самых трудных одноагентных семействах задач (10 сидов $\times$ 8 эпизодов). \textbf{(a)}~Возврат к цели в FourRooms: замена любой одной стадии оракулом оставляет сквозной успех на $0.07$; бутылочное горлышко --- ниже по потоку от стадий цепочки. \textbf{(b)}~Восстановление после опасности в LavaGapS7: оракульный контроллер добавляет лишь $+0.04$, но оракульная материализация добавляет $+0.16$, а оракульное извлечение --- ещё $+0.01$ до $0.60$. Фреймворк локализует стадию бутылочного горлышка постадийно.}
  \label{p1:fig:oracle-waterfall}
\end{figure}
\FloatBarrier

\section{Переносимость и масштабирование}
\label{p1:app:babyai}

\subsection{Перенос переносимости на BabyAI (10 сидов, полный бенчмарк)}
\label{p1:app:babyai-deeper}

10-сидовый бенчмарк переносимости на \texttt{BabyAI-GoToObjMaze-v0}~\citep{chevalier2018babyai} (8 эпизодов на сид, max-steps 200) измерил все четыре одноагентных метода, используемых в основном бенчмарке (random, greedy, \texttt{babyai\_semantic\_node\_v1}, \texttt{babyai\_semantic\_plan\_v1}), с тем же логгером Q/R/M/C, применённым к трассам BabyAI. Числа в Таблице~\ref{p1:tab:babyai-deeper} подтверждают, что (a) события Q/R/M/C нетривиально эмитируются на субстрате BabyAI, (b) два семантических варианта обнажают постадийно разделимую структуру там, где несемантические бейзлайны имеют $Q^\star{=}R^\star{=}M^\star{=}0$ (семантические запросы не выпускаются), и (c) упорядочение постадийных частот совпадает с бенчмарком MiniGrid (Таблица~\ref{p1:tab:consistency}) в пределах бутстреп-ДИ. GoToObjMaze --- более трудный задачный субстрат, чем бенчмарк LavaGapS7/FourRooms из основного текста --- сквозной успех упирается в 0.15 у всех методов, --- что делает постадийную декомпозицию особенно информативной: семантический стек не выигрывает по успеху, но является единственным семейством методов, порождающим измеримые постадийные события.

\begin{table}[h]
\centering
\caption{Бенчмарк переносимости BabyAI: 10 сидов $\times$ 8 эпизодов на \texttt{BabyAI-GoToObjMaze-v0}, max-steps 200. Варианты семантического стека успешны с той же сквозной частотой, что и жадный бейзлайн (0.15), но являются единственными методами, порождающими ненулевые события Q/R/M/C --- протокол фреймворка переносится на субстрат BabyAI без модификаций. Операциональные частоты Q/R/M/C здесь оцениваются на одном и том же знаменателе (числе точек принятия решений), поэтому $Q{=}R{=}M$ у обоих семантических вариантов, поскольку GoToObjMaze эмитирует один семантический тип цели на эпизод.}
\label{p1:tab:babyai-deeper}
\small
\begin{tabular}{l c c c c c c}
\toprule
Метод & Успех & Шаги & $Q$ (оп.) & $R$ (оп.) & $M$ (оп.) & Пост (оп.) \\
\midrule
random         & 0.125 $\pm$ 0.06 & 182.7 & 0.000 & 0.000 & 0.000 & 0.000 \\
greedy         & 0.150 $\pm$ 0.05 & 171.3 & 0.000 & 0.000 & 0.000 & 0.000 \\
semantic node  & 0.150 $\pm$ 0.05 & 171.3 & 0.032 & 0.032 & 0.032 & 0.025 \\
semantic plan  & 0.150 $\pm$ 0.05 & 171.3 & 0.032 & 0.032 & 0.032 & 0.025 \\
\bottomrule
\end{tabular}
\end{table}

\textbf{Вывод.} На второй субстрат переносится измерительный протокол фреймворка --- а не его результаты для конкретных задач. Сквозной успех на BabyAI доминируется существенно более низким потолком поэпизодного успеха GoToObjMaze, а не постадийной структурой; и наоборот, постадийная структура, которую обнажает семантический стек, идентична по форме той, что тот же стек обнажает на MiniGrid. Это ровно то свойство субстратной независимости, которое методология и призвана обеспечить. Сырые данные: \texttt{tmp/babyai\_deeper\_10seed/}.

\subsection{Многоагентное масштабирование при $N=12$}
\label{p1:app:scale-N12}

Чтобы проверить, что многоагентные результаты \S\ref{p1:sec:multi-agent-results} --- не артефакт диапазона $N\!\le\!8$, был выполнен дополнительный свип из 8{,}640 прогонов при $N\!=\!12$, $T \in \{6, 8, 10\}$, тех же трёх раскладках, трёх плотностях опасностей, восьми пировых каденциях и 40 сидах. Настенное время $\approx 12$ минут на 8 ядрах. Сигнатура условных частот из \S\ref{p1:sec:multi-agent-results} сохраняется и усиливается (Таблица~\ref{p1:tab:scale-n12}).

\begin{figure}[h]
\centering
\includegraphics[width=0.92\linewidth]{figures/fig_scale_n12.pdf}
\caption{Сигнатура бутылочного горлышка при масштабировании $N\!=\!12$ ($n\!=\!8{,}640$ прогонов, 40 сидов). Сверху: покаденционные средние частоты движутся в предсказанных направлениях по $K$; внутрикаденционная корреляция Пирсона (красный пунктир, правая ось) достигает пика $-0.70$ при $K\!=\!8$ (именно этот внутрикаденционный пик и заостряется при $N{=}12$; сам M-наклон по $K$ не является кластерно-робастным при этом числе агентов, Таблица~\ref{p1:tab:scale-n12}). Снизу: среднее время до успеха сохраняет внутренний минимум при $K\!=\!8$ ($4.77$ тиков); как и на основном свипе, положение устойчиво, но запас мелок --- зазор до соседнего $K\!=\!4$ ($4.82$ тиков) составляет $\approx0.05$ тика и не разрешён при этом размере выборки. Пик смещения бутылочного горлышка мигрирует наружу с $K\!=\!4$ при $N\!\le\!8$ до $K\!=\!8$ при $N\!=\!12$, что согласуется с тем, что более сильное связывание занятости при большем числе агентов задерживает переход между режимами издержек.}
\label{p1:fig:scale-n12}
\end{figure}

\begin{table}[h]
\caption{Числовая детализация свипа масштабирования $N\!=\!12$ (пировая архитектура; $27$ конфигурационных ячеек). Размеры эффекта с кластерно-робастными 95\%-ми ДИ (блочный бутстреп по ячейкам): Спирмен $P(C^\star|M^\star)$ с $K$ равен $+0.31$ (ДИ $[+0.21,+0.41]$, исключает ноль), тогда как Спирмен $P(M^\star|R^\star)$ с $K$ --- лишь $-0.14$ (ДИ $[-0.28,+0.02]$, \emph{не} отделён от нуля при $N{=}12$). C-наклон Эмпирического утверждения~1 выполняется; M-наклон направленно согласован, но не кластерно-робастен при этом числе агентов. При $N{=}12$ заостряются внутрикаденционная корреляция и минимум $t_{\mathrm{succ}}$ (Рисунок~\ref{p1:fig:scale-n12}), а не M-наклон по $K$.}
\label{p1:tab:scale-n12}
\centering
\small
\begin{tabular}{c c c c c c}
\toprule
$K$ & $P(M^\star|R^\star)$ & $P(C^\star|M^\star)$ & MOP & среднее $t_{\mathrm{succ}}$ & корр \\
\midrule
1  & 0.758 & 0.679 & 0.475 & 5.49 & $-0.59$ \\
2  & 0.745 & 0.692 & 0.474 & 5.29 & $-0.61$ \\
4  & 0.679 & 0.768 & 0.482 & 4.82 & $-0.64$ \\
8  & 0.731 & 0.767 & 0.539 & \textbf{4.77} & $\mathbf{-0.70}$ \\
16 & 0.687 & 0.803 & 0.538 & 5.04 & $-0.62$ \\
32 & 0.676 & 0.814 & 0.538 & 5.75 & $-0.51$ \\
48 & 0.677 & 0.814 & 0.539 & 5.96 & $-0.51$ \\
64 & 0.677 & 0.816 & 0.541 & 6.12 & $-0.51$ \\
\bottomrule
\end{tabular}
\end{table}

\subsection{Q/R/M/C при семантическом, безрамочном сопоставлении мест}
\label{p1:app:semantic-matching}

Основной свип разрешает идентичность мест общими ключами $(x,y)$. Чтобы проверить, что
декомпозиция не привязана к общему координатному субстрату, мы перезапускаем
многоагентный навигационный цикл, в котором каждый агент держит \emph{приватную}
систему координат (секретную поагентную трансляцию), а места сопоставляются
\emph{по смыслу}: локальные созвездия семантических отпечатков восстанавливают
межагентную систему координат и переносят пировые свидетельства о цели в систему координат
получателя (\method{experiments/warp/semantic\_peer\_memory.py}; раннер
\method{experiments/warp/exp\_semantic\_cadence\_qrmc.py}). Агенты навигируют
тем же планировщиком на правилах к перенесённой цели; Q/R/M/C логируются
без изменений, где $M^\star$ --- фиксация кандидата в пределах $\epsilon$ от истинной
цели ($N{=}6$, $T{=}2$ дефицитных источника воды, $8$ сидов).

Стадийные события остаются корректно определёнными без общей системы координат, и фреймворк
локализует отказ координатного предположения. \emph{Семантическое} плечо фиксирует
\emph{истинную} цель при первой материализации в $100\%$ агенто-эпизодов (при
каждой каденции); \emph{координатное} плечо --- читающее приватные
координаты пира за чистую монету --- впервые фиксирует истинную цель лишь в $25$--$52\%$
случаев (доля растёт с $K$ по мере накопления агентами собственных наблюдений), т.е.\ оно
\emph{отказывает в открытую}, материализуя фантомные клетки, не являющиеся водой. Сквозной
успех ограничен дефицитом и сопоставим между плечами ($\approx0.29\!\approx\!T/N$),
как и ожидалось. Это и есть интересующее диагностическое свойство: $M^\star$ при
координатной идентичности срабатывает на не-целях, и постадийный логгер это обнажает.

Мы \emph{не} заявляем здесь каденционного сдвига M$\leftrightarrow$C: при
рассылке все-всем и горизонте в $140$ тиков семантическая система координат восстанавливается при
каждом проверенном $K$, так что $P(M^\star|R^\star)$ каденционно-инвариантно в этой постановке
($1.00$ всюду). Эксперимент изолирует вопрос \emph{субстрата идентичности мест}:
Q/R/M/C независим от архитектуры и координат и обнаруживает
материализацию с отказом в открытую, --- а не каденционный результат, который стоит на
основном координатно-ключевом свипе. Встраивание семантического сопоставления в режим
дефицита, порождающий сдвиг, --- естественная следующая интеграция.

\section{Расширенная диагностика}
\label{p1:app:diagnostics}

\subsection{Детализация условного произведения и разрыв Йенсена}
\label{p1:app:phi-detail}

\begin{table*}[h]
\caption{Поархитектурные условные частоты Q/R/M/C из свипа в 35{,}640 прогонов (40 сидов). $P(R^\star|Q^\star)$ насыщено. POM = произведение средних; MOP = среднее произведений (правильное поконфигурационное ожидаемое произведение). Две сводки систематически расходятся при малых/промежуточных $K$, потому что внутрикаденционная корреляция между $P(M^\star|R^\star)$ и $P(C^\star|M^\star)$ там сильно отрицательна (последний столбец) --- сигнатура смещения бутылочного горлышка. \textbf{Жирным}: лучший по MOP среди пировых каденций. Среднее $t_{\mathrm{succ}}$ (последний столбец) --- практически значимая мера эффективности; её внутренний минимум при $K=8$ --- результат предсказанной формы, устойчивый по положению, но мелкий по запасу (соседние быстрые каденции не разрешены; Приложение~\ref{p1:app:effect-sizes}).}
\label{p1:tab:multi-agent-cond}
\centering
\scriptsize
\setlength{\tabcolsep}{3.5pt}
\begin{tabular}{l c c c c c c c}
\toprule
Архитектура & $P(R^\star|Q^\star)$ & $P(M^\star|R^\star)$ & $P(C^\star|M^\star)$ & POM & MOP (95\% ДИ) & корр & среднее $t_{\mathrm{succ}}$ \\
\midrule
независимая          & 0.99 & 0.68 & 0.88 & 0.598 & 0.605 [0.597, 0.614] & $+0.00$ & 8.02 \\
общая шина           & 1.00 & 1.00 & 0.60 & 0.600 & 0.595 [0.586, 0.604] & $-0.10$ & 8.33 \\
центральный агрегатор   & 1.00 & 0.97 & 0.61 & 0.592 & 0.572 [0.564, 0.580] & $-0.48$ & 7.64 \\
peer ($K=1$)         & 1.00 & 0.97 & 0.60 & 0.585 & 0.573 [0.564, 0.581] & $-0.43$ & 7.80 \\
peer ($K=2$)         & 1.00 & 0.97 & 0.62 & 0.594 & 0.576 [0.568, 0.584] & $-0.54$ & 7.59 \\
peer ($K=4$)         & 1.00 & 0.92 & 0.66 & 0.599 & 0.569 [0.561, 0.577] & $-0.61$ & 8.00 \\
peer ($K=8$)         & 0.99 & 0.69 & 0.87 & 0.597 & 0.586 [0.578, 0.595] & $-0.19$ & \textbf{7.47} \\
peer ($K=16$)        & 0.99 & 0.67 & 0.89 & 0.592 & 0.590 [0.581, 0.598] & $-0.05$ & 7.87 \\
peer ($K=32$)        & 0.99 & 0.67 & 0.89 & 0.595 & 0.598 [0.589, 0.606] & $-0.04$ & 8.28 \\
peer ($K=48$)        & 0.99 & 0.68 & 0.88 & 0.598 & 0.602 [0.593, 0.611] & $-0.02$ & 8.79 \\
peer ($K=64$)        & 0.99 & 0.68 & 0.88 & 0.599 & \textbf{0.603} [0.594, 0.611] & $-0.02$ & 9.19 \\
\bottomrule
\end{tabular}
\end{table*}

Условное произведение $\Phi = P(M^\star|R^\star)\!\cdot\!P(C^\star|M^\star)$ на свипе в 35{,}640 прогонов монотонно растёт по $K$ на $K\in\{4,8,16,32,48,64\}$ и сходится к границе независимых агентов $\Phi(K\!\to\!\infty)=0.605$. Парный бутстреп: $\Delta(K{=}32{-}K{=}16) = +0.008$ (ДИ $[+0.006, +0.010]$), $\Delta(K{=}48{-}K{=}16) = +0.011$, $\Delta(K{=}64{-}K{=}16) = +0.012$; уровень независимых агентов лежит выше каждого пирового $K$, кроме $K\!=\!64$ (где ДИ зазора включает ноль). Это утверждение об \emph{уровне} $\Phi$ на проверенных каденциях, а не о \emph{форме} $\Phi(K)$: согласуясь с монотонным трендом, мы не заявляем внутреннего экстремума (\S\ref{p1:sec:shift}).

Отрицательная внутрикаденционная ковариация, задокументированная в \S\ref{p1:sec:multi-agent-results}, создаёт разрыв неравенства Йенсена между $\Phi$ и маргинальным произведением средних: $E[X]\!\cdot\!E[Y] > E[X\!\cdot\!Y]$ при $\mathrm{Cov}(X,Y) < 0$. Среднее произведений сообщается как первичная сводка, потому что оно корректно впитывает ковариацию. Рисунок~\ref{p1:fig:bottleneck-shift} показывает K-свиповый вид обеих условных частот и среднего $t_{\mathrm{succ}}$; Рисунок~\ref{p1:fig:jensen-scatter} показывает поконфигурационный скаттер, который суммирует отрицательная ковариация.

\begin{figure}[t]
  \centering
  \includegraphics[width=0.85\linewidth]{figures/fig_bottleneck_shift.pdf}
  \caption{Смещение бутылочного горлышка по восьми пировым каденциям (40-сидовый свип). Сверху: $P(M^\star|R^\star)$ падает, а $P(C^\star|M^\star)$ растёт по $K$ (умеренные, кластерно-робастные спирменовские размеры эффекта $-0.53$ и $+0.43$; ДИ в Приложении~\ref{p1:app:effect-sizes}). Снизу: среднее время до успеха имеет внутренний минимум при $K=8$ ($7.47$ тиков) ниже бейзлайна независимых агентов ($8.02$, пунктир); положение устойчиво, но его запас над соседними быстрыми каденциями ($K{=}1,4$) не разрешён (Приложение~\ref{p1:app:effect-sizes}). Планки погрешностей и затенённая полоса --- 95\%-е бутстреп-ДИ.}
  \label{p1:fig:bottleneck-shift}
\end{figure}

\begin{figure}[t]
  \centering
  \includegraphics[width=0.95\linewidth]{figures/fig_jensen_scatter.pdf}
  \caption{Поконфигурационный скаттер $P(M^\star|R^\star)$ против $P(C^\star|M^\star)$ при трёх каденциях. Сильная отрицательная корреляция при $K=4$ (Пирсон $-0.61$) --- сигнатура смещения бутылочного горлышка: те же конфигурации, что успешны на $M$, склонны отказывать на $C$, и наоборот. Корреляция затухает до шума к $K=64$.}
  \label{p1:fig:jensen-scatter}
\end{figure}

\subsection{Декомпозиция дизайна и сверка числа прогонов}
\label{p1:app:design}

Число $35{,}640$ раскладывается в точности как в Таблице~\ref{p1:tab:design}. Базовый дизайн --- девять допустимых дефицитных пар $(N,T)$ (для $N{=}3$: $T{\in}\{2,3\}$; $N{=}5$: $T{\in}\{2,3,5\}$; $N{=}8$: $T{\in}\{2,3,5,8\}$), скрещённых с тремя раскладками, тремя плотностями опасностей и 40 сидами, что даёт $3{,}240$ базовых конфигураций. Пировая архитектура прогоняется на всех восьми каденциях ($3{,}240\times 8 = 25{,}920$); три бескаденционные архитектуры (независимая, общая шина, центральный агрегатор) вносят по $3{,}240$ каждая. Число $25{,}920$ прогонов, используемое для пирового кластерно-робастного анализа (Приложение~\ref{p1:app:effect-sizes}), --- это в точности пировый блок; первоначальный разведочный свип из $12{,}960$ прогонов (\S\ref{p1:sec:shift}, «Хронология») и свипы расширенных каденций и $N{=}12$ (Приложение~\ref{p1:app:babyai}) --- отдельные наборы прогонов и не смешиваются с основным свипом.

\begin{table}[h]
\caption{Сверка числа прогонов для основного многоагентного свипа ($35{,}640$ прогонов). Счётчики пересчитаны напрямую из опубликованного лога прогонов.}
\label{p1:tab:design}
\centering
\small
\begin{tabular}{l l r}
\toprule
Фактор & Уровни & Число \\
\midrule
Дефицитные пары $(N,T)$ & $(3,2)(3,3)(5,2)(5,3)(5,5)(8,2)(8,3)(8,5)(8,8)$ & 9 \\
Раскладки & симметричная, асимметричная, случайная & 3 \\
Плотности опасностей & $0.0,\ 0.05,\ 0.10$ & 3 \\
Сиды & & 40 \\
\midrule
Базовые конфигурации & $9\times3\times3\times40$ & 3{,}240 \\
Peer (все 8 каденций $K$) & $3{,}240\times8$ & 25{,}920 \\
Независимая / шина / центральная & $3\times3{,}240$ & 9{,}720 \\
\midrule
\textbf{Итого} & & \textbf{35{,}640} \\
\bottomrule
\end{tabular}
\end{table}

\subsection{Размеры эффекта и кластерно-робастный вывод}
\label{p1:app:effect-sizes}

Величины $p<10^{-4}$, приведённые в \S\ref{p1:sec:multi-agent-results}, сообщаются для полноты, но не являются основанием утверждения: на десятках тысяч прогонов почти любой ненулевой наклон «значим», а прогоны не независимы --- они кластеризуются внутри конфигурации $(N,T,\text{раскладка},\text{опасность})$. Полный свип содержит $25{,}920$ пировых прогонов, но лишь \textbf{81 независимую ячейку дизайна}. Поэтому мы сообщаем размеры эффекта с кластерно-робастными доверительными интервалами из блочного бутстрепа, ресемплирующего \emph{конфигурационные ячейки} (а не строки), $B{=}4000$ ресемплов.

Два наклона Эмпирического утверждения~1 выживают как подлинные, но умеренные эффекты: спирменовская корреляция $P(M^\star|R^\star)$ с $K$ равна $\rho=-0.53$ (95\%-й кластерный ДИ $[-0.59,-0.47]$, сильная), а $P(C^\star|M^\star)$ с $K$ --- $\rho=+0.43$ (95\%-й кластерный ДИ $[+0.37,+0.50]$, умеренная). Оба интервала исключают ноль при ресемплировании на уровне ячеек, так что смещение бутылочного горлышка --- не артефакт коррелированного числа прогонов; это умеренный монотонный тренд, и мы описываем его именно так, а не как почти детерминированный закон.

Внутренний оптимум на среднем времени до успеха устойчив по \emph{положению}, но мелок по \emph{запасу}. Парный блочный бутстреп на ячеечных разностях $t_{\mathrm{succ}}$ даёт $t_{\mathrm{succ}}(K{=}8)$ ниже $t_{\mathrm{succ}}(K{=}16)$ на $0.40$ тика (ДИ $[0.08,0.83]$) и ниже $K{=}64$ на $1.67$ тика (ДИ $[0.75,2.83]$), оба разрешены; но его преимущество над $K{=}4$ ($0.49$, ДИ $[-0.47,1.51]$) и $K{=}1$ ($0.27$, ДИ $[-0.43,0.97]$) --- в пределах шума.\footnote{Этот парный бутстреп и спирменовские корреляции размеров эффекта выше вычислены на том же исполнении свипа, что и Таблица~\ref{p1:tab:multi-agent-cond} ($25{,}920$ пировых прогонов); статистика $t_{\mathrm{succ}}$ определена только на успешных эпизодах, так что парная ячеечная разность ($1.67$ для $K{=}8$ против $K{=}64$) относится к этой субпопуляции и совпадает с разностью маргинальных средних в Таблице~\ref{p1:tab:multi-agent-cond} ($9.19{-}7.47{=}1.72$) с точностью до округления. Размеры эффекта идентичны на независимом исполнении \texttt{paper1\_clean} ($-0.531$/$+0.434$), что подтверждает их непривязанность к конкретному прогону.} Взвешенная успехом метрика эффективности ($t_{\mathrm{succ}}$, делённое на долю успеха) имеет тот же внутренний $\arg\min$ при $K{=}8$, так что положение оптимума устойчиво к определению эффективности, даже если его запас над соседними быстрыми каденциями не разрешён при этом размере выборки. Мы читаем это как поддержку \emph{формы}, предсказанной компромиссом M$\leftrightarrow$C, а не остро настроенной оптимальной каденции.

\subsection{Конструктная валидность: успех задачи $Y$ против завершения семантического пути $C^\star$}
\label{p1:app:y-vs-c}

Успех $Y$ в принципе может обойти семантический путь, так что постадийное бутылочное горлышко в $Q/R/M/C$ не обязано объяснять вариацию $Y$. Два факта ограничивают это опасение. Во-первых, в многоагентном свипе, о котором Эмпирическое утверждение~1, успех \emph{определяется} как достижение цели \texttt{water\_source}, что и есть в точности событие $C^\star$; эмпирически $Y$ и $C^\star$ совпадают прогон в прогон (корреляция $1.00$, средний зазор $0.00$). Центральное многоагентное утверждение тем самым не уязвимо для возражения о параллельном канале. Во-вторых, зазор $Y>C^\star$ --- одноагентный феномен на лёгких семействах задач (вода, отдых), где запасной контроллер может достичь цели без семантической фиксации цели; он наибольший ровно там, где успех уже высок и диагноз наименее нужен, и исчезает на трудных семействах (восстановление после опасности, целевой регион), где фреймворк и делает свою работу. Поэпизодная, постадийно разрешённая медиация $Y$ (а не $C^\star$) требует логов, помечающих каждый успех тем, прошёл ли он семантический путь, и является планируемым уточнением инструментирования.

\subsection{Робастность диагноза пустого множества кандидатов к порогу ($\theta$)}
\label{p1:app:theta-robustness}

Одноагентная атрибуция «отказывает генерация кандидатов, не ранжирование»
(\S\ref{p1:sec:refined-attribution}) опирается на высокую долю пустых множеств кандидатов.
Рецензент спрашивает, не артефакт ли это порога обязательных тегов
$\theta$ (развёрнутое значение $0.5$; гейт кандидатства ---
$\mathrm{profile}[\text{tag}]\ge\theta$ в
\method{symbolic\_memory.\_required\_matches\_profile}). Мы свипуем $\theta$ через
переопределение в построителе запросов (\method{experiments/big\_experiment/theta\_ablation.py});
крайние точки подтверждают, что гейт работает, а рабочий диапазон показывает, что
диагноз --- не пороговый эффект.

\begin{table}[h]
\caption{Доля пустых множеств кандидатов против порога обязательных тегов $\theta$ (семантический метод, ассист выкл., восстановление после опасности / целевой регион). Доля \emph{инвариантна} по всему реалистичному диапазону $\theta\in[0.05,0.95]$; лишь вырожденные крайние точки её сдвигают, и лишь до пола отсутствия кандидатов. Пустоты, следовательно, --- отсутствие сохранённых свидетельств, а не подпороговая фильтрация.}
\label{p1:tab:theta-robustness}
\centering
\small
\begin{tabular}{c c c}
\toprule
$\theta$ & доля пустых (восст. после опасности) & доля пустых (целевой регион) \\
\midrule
$0.00$ (принять всё) & 0.50 & --- \\
$0.05$ & 0.533 & 0.267 \\
$0.15$ & 0.533 & 0.267 \\
$0.25$ & 0.533 & 0.267 \\
$0.50$ (развёрнутый) & 0.533 & 0.267 \\
$0.75$ & 0.533 & 0.267 \\
$0.95$ & 0.533 & 0.267 \\
$1.00$ (не принимать ничего) & 0.75 & --- \\
\bottomrule
\end{tabular}
\end{table}

Прочтение решающее (Таблица~\ref{p1:tab:theta-robustness}). При $\theta{=}1.0$ ничего не проходит, и доля пустых растёт; при $\theta{=}0.0$ гейт обязательных тегов принимает каждый узел ($\mathrm{conf}\ge 0$ выполняется всегда), однако доля пустых падает лишь до $0.50$ --- своего пола, --- то есть в этих точках принятия решений \emph{нет узла-кандидата какой бы то ни было уверенности} для извлечения. По всему разумному диапазону $\theta\in[0.05,0.95]$ доля плоская на $0.533$: ослабление или ужесточение порога не возвращает кандидата, потому что требуемое семантическое место никогда не было сохранено. Это в точности режим «отсутствие свидетельств, а не противоречие», названный в \S\ref{p1:sec:event-semantics}, теперь показанный порогово-инвариантным, так что атрибуция накопления памяти --- не артефакт $\theta$. (Абсолютные доли здесь --- на малом свипе восстановления после опасности / целевого региона и не являются заголовочными $91\%$ из \S\ref{p1:sec:refined-attribution}, которые относятся к полному бенчмарку; суть в $\theta$-инвариантности, а не в уровне.)

\subsection{Робастность к нарушениям Предположения~1}
\label{p1:app:assumption1-stress}

Определение~1 покоится на условном постадийно-марковском свойстве (Предположение~1). Два практически значимых нарушения --- адаптивные переписывания запроса, не всплывающие в логгер, и скрытые межстадийные обновления памяти --- ровно те случаи, когда \emph{нелогируемая латентная} переменная связывает несмежные стадии. Мы количественно оцениваем результирующее смещение оценки контролируемой симуляцией ($4\!\times\!10^5$ эпизодов на настройку; \method{experiments/big\_experiment/assumption1\_stress\_test.py}). Эпизоды генерируются с известными механизменными частотами $P(M^\star|R^\star){=}0.70$, $P(C^\star|M^\star){=}0.65$ и скрытой бинарной латентной $U$, сдвигающей \emph{обе} частоты M и C на $\pm\delta$; наивная оценка обусловливается лишь предыдущей звёздной стадией, не $U$.

\begin{table}[h]
\caption{Смещение наивных условных оценок при нелогируемой латентной силы $\delta$, связывающей стадии M и C (нарушая Предположение~1). $P(M^\star|R^\star)$ остаётся несмещённым; $P(C^\star|M^\star)$ смещено вверх, существенно лишь при сильном связывании.}
\label{p1:tab:assumption1-stress}
\centering
\small
\begin{tabular}{c c c}
\toprule
$\delta$ (сила связывания) & смещение $\hat P(M^\star|R^\star)$ & смещение $\hat P(C^\star|M^\star)$ \\
\midrule
0.00 & $-0.001$ & $-0.002$ \\
0.05 & $+0.001$ & $+0.004$ \\
0.10 & $+0.001$ & $+0.015$ \\
0.15 & $-0.000$ & $+0.031$ \\
0.20 & $+0.000$ & $+0.056$ \\
0.25 & $-0.001$ & $+0.089$ \\
\bottomrule
\end{tabular}
\end{table}

Два вывода (Таблица~\ref{p1:tab:assumption1-stress}). Условная вероятность R$\to$M несмещена при каждом $\delta$: латентная переменная перевзвешивает эпизоды, но не конфаундит это ребро. Условная вероятность C$\to$M приобретает однознаковое смещение \emph{вверх} --- латентная делает эпизоды с высоким M также эпизодами с высоким C, так что обусловливание на M пересчитывает завершение, --- но смещение составляет лишь $+1.5$ процентного пункта при мягком нарушении ($\delta\!\le\!0.10$) и достигает $+5.6$ пункта лишь при сильном нелогируемом связывании ($\delta{=}0.20$). Уязвимость фреймворка к нарушениям Предположения~1, следовательно, ограничена, и направление его отказа известно --- что и есть релевантное утверждение о робастности для применения логгера к обученным или LLM-агентам, где такие связывания более обычны.

\subsection{Предварительная проверка на непрерывном субстрате (VMAS)}
\label{p1:app:vmas}

Самая запрашиваемая проверка робастности --- специфично ли смещение бутылочного горлышка
M$\leftrightarrow$C для сеточных субстратов. Мы сообщаем \emph{предварительный}
положительный результат на субстрате с непрерывным пространством состояний (VMAS~\citep{bettini2022vmas}).
Символьная пировая память и операциональные определения Q/R/M/C переиспользуются
без изменений; тонкий мост дискретизации отображает непрерывные $(x,y)$ в грубые клетки,
чтобы память была применима, а контроллер действует в непрерывном пространстве к
материализованной путевой точке (\method{experiments/vmas\_portability/}). Материализация
использует ту же чувствительную к занятости фиксацию, что и сеточный планировщик (кандидат, чья
цель уже заявлена, не фиксируется).

\begin{table}[h]
\caption{Предварительная переносимость VMAS: оба знака наклонов Эмпирического утверждения~1 воспроизводятся на непрерывном субстрате в режиме связывания ($N{=}8$, $T{=}3$, 40 сидов, горизонт $25$). $\mathrm{Spearman}(P(M^\star|R^\star),K){=}-1.00$ и $\mathrm{Spearman}(P(C^\star|M^\star),K){=}+1.00$. За пределами $K\!\approx\!32$ система входит в независимый режим (редкая рассылка, нет соперничества), и $P(M^\star|R^\star)$ отскакивает --- та же граница режимов, что видна на сетке.}
\label{p1:tab:vmas}
\centering
\small
\begin{tabular}{c c c}
\toprule
$K$ & $P(M^\star|R^\star)$ & $P(C^\star|M^\star)$ \\
\midrule
1  & 1.000 & 0.182 \\
2  & 0.583 & 0.305 \\
4  & 0.448 & 0.405 \\
8  & 0.305 & 0.608 \\
16 & 0.278 & 0.639 \\
\bottomrule
\end{tabular}
\end{table}

Мы помечаем это как предварительное и не вливаем в заголовочные утверждения: это
одна конфигурация сценария, горизонт короток, кластерно-робастных ДИ пока нет,
память питается через мост дискретизации (непрерывная
\emph{динамика} с дискретизованной символьной \emph{памятью}, а не полностью непрерывная
память), а сообщаемый Спирмен $\pm1.00$ --- \emph{идеальная} ранговая
корреляция всего лишь по пяти монотонным точкам ($K\le16$) --- слабое свидетельство само по
себе, указывающее направление, а не разрешённый размер эффекта. Что оно устанавливает --- так это что два знака наклонов не являются артефактом
клеток сетки: тот же механизм «свежесть против соперничества» действует, когда позиции
и движение непрерывны. Повышение этого до заголовочного результата требует масштабирования
до нескольких ячеек дефицита и раскладок с кластерно-робастным бутстрепом из
Приложения~\ref{p1:app:effect-sizes}.

\subsection{Каденционная робастность: сдвиг отслеживает дефицит, а не устройство сетки}
\label{p1:app:cadence-robustness}

Рецензент может возразить, что смещение бутылочного горлышка встроено в сетку с дефицитными
целями и фиксациями. Если же это сигнатура соперничества за \emph{дефицитные
общие} цели, её условия на наклоны должны ломаться при изобилии целей. Мы
перезапускаем пировый каденционный свип в трёх режимах на той же среде ($N{=}8$, случайная
раскладка, опасность $0.05$, $20$ сидов; \method{experiments/big\_experiment/exp\_cadence\_robustness.py}):
дефицит ($T{=}3$), пограничный ($T{=}N{=}8$) и изобилие ($T{=}14{>}N$).

\begin{table}[h]
\caption{Два условия на наклоны Эмпирического утверждения~1 (\S\ref{p1:sec:shift}) выполняются при дефиците и выполаживаются при изобилии. $P(M^\star|R^\star)$ должно падать, а $P(C^\star|M^\star)$ расти по $K$; «корр» --- внутрикаденционная корреляция Пирсона двух звеньев при $K{=}64$ (сигнатура сдвига). $P(C^\star|M^\star)$ --- операциональное отношение (ограничено $1$; изобилие завершает без формальной фиксации).}
\label{p1:tab:cadence-robustness}
\centering
\small
\begin{tabular}{l c c c c}
\toprule
Режим & $P(M^\star|R^\star)$: $K{=}1\!\to\!64$ & $P(C^\star|M^\star)$: $K{=}1\!\to\!64$ & корр$(M,C)$ @ $K{=}64$ & argmin $t_{\mathrm{succ}}$ \\
\midrule
Дефицит ($T{=}3$)      & $0.86 \to 0.49$ (падает) & $0.51 \to 0.79$ (растёт) & $-0.86$ & $K{=}4$ (внутр.) \\
Пограничный ($T{=}N{=}8$) & $0.79 \to 0.72$        & $0.77 \to 0.95$ (растёт) & $-0.79$ & $K{=}8$ (внутр.) \\
Изобилие ($T{=}14$)    & $0.74 \to 0.77$ (плоско)  & $1.00 \to 0.99$ (плоско)  & $-0.29$ & $K{=}4$ \\
\bottomrule
\end{tabular}
\end{table}

Оба наклона присутствуют при дефиците, ослабевают на границе и \emph{выполаживаются}
при изобилии: при большем числе целей, чем агентов, соперничества нет, так что быстрый
обмен ни перенасыщает материализацию ($P(M^\star|R^\star)$ остаётся плоским
на $\sim0.75$ вместо всплеска с последующим падением), ни коллапсирует завершение
($P(C^\star|M^\star)$ остаётся на $\sim1$ даже при $K{=}1$ вместо падения до
$0.51$). Внутрикаденционная антикорреляция M$\leftrightarrow$C монотонно
ослабевает с изобилием ($-0.86\to-0.79\to-0.29$). Сдвиг, следовательно,
отслеживает \emph{механизм соперничества} и исчезает при его устранении --- он
не является фиксированным артефактом дизайна «сетка плюс фиксации». Честно отметим, что
внутренний оптимум $t_{\mathrm{succ}}$ более общ (присутствует во всех трёх
режимах): он отражает издержки исследования/пути очень быстрого \emph{или} очень
медленного обмена, где соперничество за цели --- один из вкладчиков, но не единственный.

\subsection{Открытые проблемы, подсказанные диагностикой}
\label{p1:app:open-problems}

Четыре конкретных вопроса следуют напрямую из измерений \S\ref{p1:sec:refined-attribution} и \S\ref{p1:sec:multi-agent-results}.
\textbf{(1)~Адаптивная каденция.} Эмпирическое утверждение~1 предполагает фиксированный $K$; сигнатура смещения бутылочного горлышка подсказывает онлайн-контроллер $K_t$, рассылающий, когда (оцененный) выигрыш на стороне M превышает занятостную цену на стороне C.
\textbf{(2)~Селективная консолидация.} Обмен снимками сливает всё; правила слияния, приоритизирующие свидетельства о недопредставленных семантических тегах, атакуют голодание генерации кандидатов из \S\ref{p1:sec:refined-attribution} напрямую.
\textbf{(3)~Слияние, взвешенное доверием, при занятости.} Центральный агрегатор уже выполняет слияние, взвешенное доверием, но платит цену C-коллапса; децентрализованные правила доверия, учитывающие \emph{намерение} пира (а не только его свидетельства), могли бы расцепить свежесть и гонку за цели. \emph{Со времени подачи первого драфта эта проблема решена конструктивно}: децентрализованный протокол, в котором фиксация цели агента подсаживает лёгкое \emph{резервирование} на его следующую запланированную рассылку (без координатора, без дополнительного канала), а фиксации, доминируемые более ранним резервированием, отзываются анти-$M^\star$ откатом с экспоненциальным бэкоффом. На ячейках дефицита это возвращает потерю завершения быстрых каденций на \emph{безусловных} метриках --- доля успеха растёт при каждом $K\in\{1,2,4\}$ с непересекающимися бутстреп-ДИ, а быстрый обмен с протоколом превышает бейзлайн оптимума $K{=}8$ ($0.614$--$0.620$ против $0.583$) --- с латентностью отката, ограниченной ${\approx}K$ тиками. Детали --- в сопутствующей рукописи; открытая проблема пересмотрена в сторону её более трудных вариантов (адаптивная область резервирования, гетерогенное доверие).
\textbf{(4)~Общность субстрата.} Декомпозиция переносится между MiniGrid, BabyAI, собственной сеткой и многоагентной обёрткой стандартного MiniGrid с идентичными операциональными определениями и подтверждёнными знаками наклонов Эмпирического утверждения~1 ($\S$~\ref{p1:sec:limitations}), а предварительная проверка воспроизводит оба наклона на VMAS-субстрате с непрерывным пространством состояний (Приложение~\ref{p1:app:vmas}). Декомпозиция также координатно-независима: при \emph{семантическом}, безрамочном сопоставлении мест (агенты в приватных системах координат, места сопоставляются по смыслу) стадийные события остаются корректно определёнными и обнажают материализацию координатного предположения с отказом в открытую (Приложение~\ref{p1:app:semantic-matching}). Полная валидация в непрерывном пространстве состояний с кластерно-робастными ДИ и более богатые 3D-субстраты (Habitat~\citep{savva2019habitat}) остаются следующими шагами переносимости.

\subsection{Расширенный обзор литературы}
\label{p1:app:related-work-extended}

Классические перспективы когнитивных карт мотивируют явные абстракции мест \citep{tolman1948cognitive,okeefe1978hippocampus,kuipers2000spatial}; современный воплощённый ИИ сочетает картирование и планирование \citep{gupta2017cognitive} или открывает топологическую память (SPTM, \citealp{savinov2018semi}); целеобусловленный RL обусловливается явными целями или эмбеддингами \citep{andrychowicz2017hindsight,ghosh2021learning}, а работы о репрезентации преемника дают предиктивную структуру \citep{dayan1993improving,stachenfeld2017hippocampus,momennejad2017successor}. Наша постановка отличается тем, что цель --- семантический паттерн, извлекаемый из памяти, а не координата. В многоагентном RL доминирует CTDE \citep{lowe2017maddpg,foerster2018ctde}, а обученная коммуникация добавляет пошаговую переменную обмена \citep{sukhbaatar2016commnet,foerster2016learning}; протоколы gossip и консенсуса \citep{boyd2006gossip,xiao2007distributed} --- ближайший сетевой аналог нашей оси каденции, поднятой здесь на уровень когнитивной архитектуры. Системы LLM-агентов сочетают рассуждение/действие, рефлексию, явную память и многоагентную беседу \citep{yao2022react,shinn2023reflexion,wang2023voyager,park2023generative,wu2023autogen,packer2023memgpt,chhikara2025mem0}. Наш фреймворк ортогонален: вместо обучения коммуникационной политики сквозным образом он \emph{диагностирует}, где в цепочке Q/R/M/C данная архитектура успешна или отказывает, и даёт предсказание уровня наклонов о каденции, проверяемое на логах.

\subsection{SOTA-бейзлайн и матрица покрытия}
\label{p1:app:sota-coverage}

Таблица~\ref{p1:tab:sota-coverage} разделяет прямые экспериментальные бейзлайны и более широкие SOTA-референсы. Локальные эксперименты намеренно скромны и полностью инструментированы: random / graph-only / SPTM-like / обученный BC в MiniGrid, CommNet PPO в многоагентной сетке и исследование бюджета шагов LLM на TextNav на двух локальных моделях Ollama (в котором Q/R/M/C локализует отказ, ограниченный завершением, \S\ref{p1:sec:multi-agent-results}). ReAct, Reflexion, Voyager, Generative Agents, AutoGen, MemGPT и Mem0 здесь не реимплементируются как полные системы; они определяют текущее пространство дизайна агентов с памятью, относительно которого Q/R/M/C позиционируется как диагностический протокол.

\paragraph{Позиционирование относительно SOTA-дизайнов памяти.}
Пространство дизайна агентов с памятью можно организовать по двум осям: один против многих агентов, и как память представлена и разделяется (Таблица~\ref{p1:tab:sota-quadrants}). Векторно-эмбеддинговая память (MemGPT, Voyager --- одноагентные; Generative Agents --- многоагентная) извлекает непрозрачным поиском ближайших соседей, так что отказ не локализуем по стадии. Дизайны с общим контекстом (MetaGPT, AutoGen) сваливают всё в одно окно, растворяя поагентное состояние. MARL с обученной коммуникацией (CommNet и более сильные CTDE-методы QMIX \citep{rashid2018qmix}, MAPPO \citep{yu2022mappo}, MADDPG \citep{lowe2017maddpg}) оптимизирует канал сообщений сквозным образом и, как мы верифицируем для CommNet PPO (\S\ref{p1:sec:multi-agent-results}), не эмитирует разделимых событий запроса/фиксации/завершения. Ячейка, которую занимает наша система --- поагентная \emph{символьная} память с явным интерфейсом запроса/извлечения/материализации, обмениваемая с настраиваемой каденцией, --- в этой матрице пуста, и именно этот интерфейс делает четыре стадии Q/R/M/C наблюдаемыми. Наше утверждение против SOTA, следовательно, --- не более высокая награда, а более высокая \emph{диагностируемость при сопоставимом успехе задачи}: тот самый интерфейс, который сильный обученный метод выторговывает в пользу сквозной оптимизации, и позволяет фреймворку локализовать, где память помогает.

\begin{table}[h]
\caption{Где фреймворк сидит в пространстве дизайна агентов с памятью. Постадийная наблюдаемость Q/R/M/C требует явного интерфейса запроса/фиксации, который и предоставляет ячейка «символьная поагентная».}
\label{p1:tab:sota-quadrants}
\centering
\small
\begin{tabular}{p{3.1cm} p{4.3cm} p{4.3cm}}
\toprule
Представление памяти & Один агент & Много агентов \\
\midrule
Векторно-эмбеддинговая память & MemGPT, Voyager \citep{packer2023memgpt,wang2023voyager} & Generative Agents \citep{park2023generative} \\
Общее контекстное окно & --- & AutoGen \citep{wu2023autogen} \\
MARL с обуч.\ коммуникацией & --- & CommNet, QMIX, MAPPO, MADDPG \citep{sukhbaatar2016commnet,rashid2018qmix,yu2022mappo,lowe2017maddpg} \\
Символьная поагентная + явный интерфейс & Q/R/M/C один агент (\S\ref{p1:sec:single}) & \textbf{эта работа} (пировая память, каденция $K$) \\
\bottomrule
\end{tabular}
\end{table}

\begin{table*}[h]
\caption{SOTA-покрытие для Статьи 1. «Инструментировано» означает, что прогон эмитирует операциональные события $Q^\star,R^\star,M^\star,C^\star$, используемые Определением~1.}
\label{p1:tab:sota-coverage}
\centering
\small
\begin{tabular}{p{2.45cm}p{2.55cm}p{2.7cm}p{4.0cm}}
\toprule
Семейство & Репрезентативный метод & Статус в этой статье & Точка сравнения Q/R/M/C \\
\midrule
Топологическая навигационная память & SPTM-подобный графовый бейзлайн \citep{savinov2018semi} & Прямой бейзлайн MiniGrid & Проверяет, обнажает ли топологическая память без семантического запроса/фиксации нетривиальные стадийные события. \\
Клонирование поведения & Обученный BC на сеточных наблюдениях & Прямой бейзлайн MiniGrid & Сильный успех на лёгких задачах, но без явного интерфейса запроса/фиксации, так что диагностика структурно ограничена. \\
Обученная коммуникация & CommNet PPO \citep{sukhbaatar2016commnet,schulman2017proximal} & Прямой многоагентный бейзлайн & Конкурентоспособный успех; $Q^\star$ насыщается и R/M схлопываются, верифицируя потерю наблюдаемости. \\
Кооперативный MARL (CTDE / декомпозиция ценности) & QMIX, MAPPO, MADDPG \citep{rashid2018qmix,yu2022mappo,lowe2017maddpg} & Референс первоисточников & Более сильные сквозные бейзлайны, чем CommNet; они оптимизируют совместную политику/канал сообщений без явного интерфейса запроса/фиксации, так что применим тот же аргумент о коллапсе наблюдаемости --- фреймворк квантифицирует диагностическую цену неявной координации, а не соревнуется по награде. \\
LLM-агенты рассуждения/действия & ReAct \citep{yao2022react} & Референс первоисточников & Q/R/M/C может инструментировать результирующие трассы действий, если логируются запросы, извлечённые свидетельства, фиксации целей и завершения. \\
Рефлексивные LLM-агенты & Reflexion \citep{shinn2023reflexion} & Референс первоисточников & Рефлексию можно трактовать как вышестоящую ревизию Q/R; фреймворк спрашивает, улучшает ли она R, M или C. \\
Воплощённые LLM-агенты & Voyager \citep{wang2023voyager} & Референс первоисточников & Извлечение навыков/памяти естественно отображается на R и M, но полное воспроизведение вне рамок этой статьи. \\
Социальные/памятные симуляции & Generative Agents \citep{park2023generative} & Референс первоисточников & Их цикл наблюдение/память/рефлексия мотивирует логирование стадийно-специфичного использования памяти, а не только правдоподобия. \\
Многоагентные LLM-фреймворки & AutoGen \citep{wu2023autogen} & Референс первоисточников & Трассы бесед дают кандидатные события Q/R/M/C; эта статья поставляет измерительный контракт. \\
Системы долговременной LLM-памяти & MemGPT / Mem0 \citep{packer2023memgpt,chhikara2025mem0} & Референс первоисточников & Персистентные системы памяти можно оценивать по тому, где они отказывают: извлечение, материализация или завершение. \\
Локальный LLM-субстрат & qwen3:4b, llama3.1 & Прямое исследование бюджета шагов TextNav & Q/R/M/C локализует отказ, ограниченный завершением ($C^\star{\to}0$ при насыщенных $Q/R/M$), в реальных LLM-агентах; см.\ Таблицу~\ref{p1:tab:llm-smoke}. \\
\bottomrule
\end{tabular}
\end{table*}

\subsection{Числа многоагентных оракульных интервенций}
\label{p1:app:multi-oracle}

\begin{table}[h]
\caption{Многоагентные оракульные интервенции на случаях дефицита при трёх каденциях. Среднее время до успеха (и доля успеха), усреднённые по $(N,T)\in\{(3,2),(5,3),(8,5)\}$, двум раскладкам (асимметричная, случайная), 15 сидам. Оракульное R возвращает истинные позиции воды; оракульное M фиксирует ближайшую реальную воду; \textbf{оракульное C} гарантирует завершение реальной материализованной фиксации (Приложение~\ref{p1:app:oracle-c-def}). Оракульное C даёт наибольшее сокращение \emph{времени} (устраняет путь), но наименьший прирост \emph{успеха} --- завершение есть временная издержка, а не бутылочное горлышко успеха, которое сидит выше по потоку на R/M.}
\label{p1:tab:multi-agent-oracle}
\centering
\small
\begin{tabular}{c c c c c}
\toprule
$K$ & Бейзлайн & Оракул R & Оракул M & Оракул C \\
\midrule
\multicolumn{5}{l}{\emph{среднее $t_{\mathrm{succ}}$}}\\
1  & 10.33 & 7.22 & 7.21 & \textbf{6.19} \\
4  & 10.98 & 7.22 & 7.21 & \textbf{6.07} \\
16 &  8.45 & 7.22 & 7.21 & \textbf{6.41} \\
\midrule
\multicolumn{5}{l}{\emph{доля успеха}}\\
1  & 0.591 & \textbf{0.614} & 0.602 & 0.591 \\
4  & 0.575 & \textbf{0.614} & 0.602 & 0.600 \\
16 & 0.578 & \textbf{0.614} & 0.602 & 0.598 \\
\bottomrule
\end{tabular}
\end{table}

\subsection{Оракульное C: определение и прочтение}
\label{p1:app:oracle-c-def}

Оракульное C --- аналог оракульных R и M для стадии завершения, добавленный, чтобы симметрично замкнуть тройку R/M/C. Реальный конвейер Q/R/M работает без изменений (реальный запрос к памяти, реальная материализация), и как только агент зафиксировал цель, являющуюся подлинным \texttt{water\_source}, оракул гарантирует прибытие, помещая агента на эту клетку; дефицит сохраняется, потому что уже заявленная или занятая клетка недоступна, так что оракульное C не может сделать успешными более $T$ агентов. Раннер --- \method{experiments/big\_experiment/exp\_oracle\_interventions.py} (режим \method{oracle\_C}).

Из Таблицы~\ref{p1:tab:multi-agent-oracle} следуют два прочтения. (i)~Оракульное C даёт \emph{наименьшее} среднее $t_{\mathrm{succ}}$ из трёх интервенций ($\sim6.1$--$6.4$ против $\sim7.2$ у R и M), потому что оно устраняет всё время пути после фиксации; это подтверждает, что реальная часть избыточного времени --- путь контроллера к корректно материализованной цели. (ii)~Оракульное C даёт \emph{наименьший} прирост доли успеха ($+0.01$--$0.02$ над бейзлайном, ниже и оракульного R с $0.614$, и оракульного M с $0.602$), потому что при дефиците ограничение на то, \emph{успешен ли} агент, находится выше по потоку --- извлечь и зафиксировать ещё доступную цель, --- а не достичь её. Два факта вместе --- это в точности тезис фреймворка в интервенционной форме: завершение --- где тратится время, извлечение/материализация --- где выигрывается успех.

\subsection{Безразличные к стадиям каденционные гипотезы и что меняет фреймворк}
\label{p1:app:standalone-hypotheses}

Первая версия многоагентного каденционного свипа (12{,}960 прогонов при 20 сидах; поглощена свипом из 35{,}640 прогонов в этой статье) изначально планировалась с четырьмя операциональными гипотезами, сформулированными \emph{до} Эмпирического утверждения~1. Они сохранены здесь, потому что мотивировали итоговый фреймворк. Фреймворк был сформулирован после завершения свипа; не заявляется, что он предсказал числа свипа априори (см. «Хронологию» в \S\ref{p1:sec:shift}).

\textbf{Исходные гипотезы (до фреймворка).}
\begin{description}[leftmargin=10pt,topsep=2pt,itemsep=1pt]
\item[H1.] Внутренний $K^\star \in \{2,4,8\}$, минимизирующий среднее время до успеха, существует для большинства ячеек $(N, T, \text{раскладка})$.
\item[H2.] $K^\star$ монотонно растёт с дефицитом $\rho = N/T$.
\item[H3.] При $\rho \le 1$ $K^\star = 1$.
\item[H4.] Пирово-широковещательные архитектуры Парето-доминируют централизованную агрегацию по (среднему времени до успеха, доле успеха).
\end{description}

\textbf{Автономный результат на исходном 20-сидовом свипе.}
Практическое утверждение было направленно поддержано: peer при эмпирически лучшем $K$ был быстрее централизованного в режиме дефицита (скромный эффект; Манн--Уитни $p{=}0.044$ на этом малом автономном свипе, сообщается как разведочный, а не как первичное свидетельство, в согласии с линией «сначала размер эффекта» Приложения~\ref{p1:app:effect-sizes}). Из четырёх структурных гипотез только H4 была явно поддержана на подмножестве случайных раскладок. H1 выполнялась в 7 из 18 ячеек; H2 имела незначимую отрицательную корреляцию; H3 поддержана в 4 из 9 допустимых ячеек. Расширенный свип из 35{,}640 прогонов, используемый в остальной части статьи, даёт совпадающие автономные числа в пределах бутстреп-неопределённости.

\textbf{Что меняет фреймворк.}
Гипотезы H1--H4 безразличны к стадиям: каждая привязывает $K^\star$ к дефициту $\rho$, не специфицируя механизм. Эмпирическое утверждение~1, сформулированное после свипа, привязывает наклоны $P(M^\star|R^\star)$ и $P(C^\star|M^\star)$ к каденции; условия на наклоны проверяемы на логах. Данные верифицируют оба знака как умеренные, кластерно-робастные размеры эффекта (Приложение~\ref{p1:app:effect-sizes}). Мы не делаем формальных утверждений о форме условного произведения $\Phi$, по которому данные монотонны. Внутренний оптимум выполняется на метрике эффективности --- среднем времени до успеха.

\section{Анализ ассиста вкл/выкл}
\label{p1:app:assist}

По всем четырём одноагентным семействам задач итоговый успех практически неизменен при ассисте вкл/выкл ($\Delta \le 0.10$). Единственный нетривиальный эффект --- на отдыхе, где ассисты улучшают завершение после извлечения на~$+0.05$, не меняя частот запроса/извлечения, что согласуется с действием ассистов на материализацию и пост-извлеченческое исполнение.

\section{Переобученный ретривер и бейзлайн CommNet}
\label{p1:app:learned-baselines}

\subsection{Переобученный MLP генерации кандидатов (\S\ref{p1:sec:refined-attribution})}

\textbf{Извлечение датасета.} Из 10-сидового бенчмарка восстановления после опасности
(MiniGrid-LavaGapS7-v0, метод \texttt{milestone\_state\_conditioned\_hazard\_recovery\_v7})
извлечён 18-мерный вектор признаков для каждого шага каждого эпизода каждого сида --- 3 числовых скаляра состояния агента (energy, thirst, risk\_budget),
9 уверенностей семантических тегов, 6 one-hot намерения --- и бинарная метка: \emph{был ли
этот символьный узел когда-либо выбран кандидатом в течение прогона?} Итоги датасета:
814 примеров, 113 положительных (13.9\%).

\textbf{Разбиение.} По сидам: train $\{2,3,4,5,6,7\}$ (505 примеров, 9.7\% положительных),
val $\{8,9\}$ (170, 22.4\% положительных), test $\{10,11\}$ (139, 18.7\% положительных).

\textbf{Архитектура.} $18 \to 32 \to 32 \to 1$, активации ReLU, BCE с
$\text{pos\_weight}\!\approx\!9.3$ для дисбаланса классов, Adam $\text{lr}=2{\cdot}10^{-3}$
с $\text{weight\_decay}=10^{-4}$, батч 64, 200 эпох, лучшая по val-лоссу.

\textbf{Результаты.}

\begin{table}[h]
\centering
\small
\begin{tabular}{l c c c}
\toprule
Метрика & Train & Val & Test \\
\midrule
Точность                          & 0.584 & 0.553 & 0.590 \\
Точность мажоритарного бейзлайна  & 0.903 & 0.776 & 0.813 \\
Precision @ recall $=0.70$        & 0.171 & 0.278 & 0.275 \\
\bottomrule
\end{tabular}
\end{table}

\begin{figure}[h]
\centering
\includegraphics[width=0.95\linewidth]{figures/fig_refined_attribution.pdf}
\caption{Диагностика уточнённой атрибуции. \emph{Слева:} 349 вызовов извлечения, агрегированных по 10 сидам; 317 (90.8\%) вернули пустые множества кандидатов, 31 (8.9\%) вернул хотя бы одного удовлетворяющего кандидата, лишь 1 (0.3\%) вернул кандидатов, отранжированных, но неудовлетворяющих. Бутылочное горлышко --- генерация кандидатов, а не их ранжирование. \emph{Справа:} обученный MLP уступает мажоритарному бейзлайну на всех трёх сид-дизъюнктных разбиениях, подтверждая, что одних лишь мгновенных признаков состояния недостаточно для предсказания пригодности кандидата.}
\label{p1:fig:refined-attribution}
\end{figure}

MLP уступает мажоритарному классу по точности на всех разбиениях и
достигает лишь умеренной precision при фиксированном recall. Мгновенных признаков
состояния недостаточно, чтобы предсказать пригодность кандидата ниже по потоку. Это
то самое негативное-но-информативное наблюдение, на которое ссылается \S\ref{p1:sec:refined-attribution}:
атрибуция фреймворка «ограничено извлечением» уточняется до
«ограничено накоплением памяти», а не «ограничено ранжированием».

\subsection{Бейзлайн обученной коммуникации в стиле CommNet (\S\ref{p1:sec:multi-agent-results})}

\textbf{Архитектура политики.} Наблюдение агента --- 4-мерный one-hot направления
$+$ локальное окно типов клеток $5\!\times\!5\!\times\!5$ (125) $+$ 2-мерная нормализованная
позиция $=$ 131. Энкодер с общими весами $131{\to}64{\to}64$ (ReLU), 16-мерная
проекция коммуникации, mean-pool по пирам (исключая себя), объединение $80{\to}64$ (ReLU),
голова актора $64\!\to\!4$ (TURN\_LEFT, TURN\_RIGHT, FORWARD, NOOP), голова критика
$64\!\to\!1$. Это архитектура обученной коммуникации в стиле CommNet~\citep{sukhbaatar2016commnet}.

\textbf{Обучение.} REINFORCE~\citep{williams1992simple} с ценностным бейзлайном; лосс $=\!-\,\mathrm{mean}(\log\pi(a)\cdot(R\!-\!V))\,+\,0.5\cdot\mathrm{MSE}(V,R)\,-\,0.01\cdot H[\pi]$; Adam $\text{lr}=3{\cdot}10^{-4}$, клиппинг градиента $1.0$. Шейпинг награды: $+1.0$ за первый успех на агента, $-0.01$ шаговая цена, $-0.1$ за контакт с опасностью. 2000 эпизодов с циклом по 200 сидам среды, настенное время $100$\,с на CPU. Оценка на отложенных 50 сидах среды 250--299.

\textbf{Кривая обучения.} Поэпизодный успех (окно последних 100) остаётся в
$[0.13, 0.23]$ на протяжении обучения; монотонного улучшения после $\sim\!200$
эпизодов нет. Политика сходится к локальному оптимуму, где один из трёх агентов
находит воду случайно.

\begin{figure}[h]
\centering
\includegraphics[width=0.95\linewidth]{figures/fig_commnet_baseline.pdf}
\caption{Диагностика бейзлайна обученной коммуникации. \emph{Слева:} кривая обучения CommNet (сглаженная доля успеха по окну в 100 эпизодов) на 2000 обучающих эпизодов; успех остаётся ниже 0.25 на всём протяжении. Пунктирная линия --- символьная пировая архитектура при $K\!=\!8$ на том же сценарии для сравнения. \emph{Справа:} профиль Q/R/M/C обученного CommNet на 50 отложенных сидах против символьного peer при $K\!=\!8$. $Q^\star$ насыщен на $1.0$ у CommNet (вырожден по построению); $R/M$ схлопываются вместе; только у символьной архитектуры условные частоты разделены.}
\label{p1:fig:commnet}
\end{figure}

\textbf{Оценка Q/R/M/C (50 отложенных сидов, стохастическое сэмплирование действий).}

\begin{table}[h]
\centering
\small
\begin{tabular}{l c}
\toprule
Величина & Значение \\
\midrule
Частота $Q^\star$                              & \textbf{1.00} (насыщена) \\
Частота $R^\star$                              & 0.50 \\
Частота $M^\star$                              & 0.50 (схлопывается на $R$) \\
Частота $C^\star$                              & 0.19 \\
$P(R \mid Q)$                               & 0.50 \\
$P(M \mid R)$                               & 1.00 \\
$P(C \mid M)$                               & 0.42 \\
Доля успеха                                & 0.19 \\
Средние шаги                                  & 80 (бюджет исчерпан) \\
\bottomrule
\end{tabular}
\end{table}

Сравните с символьной пировой архитектурой при $K{=}8$: $P(M^\star|R^\star){=}0.69$, $P(C^\star|M^\star){=}0.87$, успех $\sim 0.65$.

\paragraph{Вариант PPO: структурная вырожденность верифицирована при конкурентоспособном успехе.}
Обученный PPO~\citep{schulman2017proximal} вариант той же архитектуры CommNet (клип $\epsilon{=}0.2$, обобщённая оценка преимущества~\citep{schulman2016gae} с $\lambda{=}0.95$, $\gamma{=}0.99$, 4 PPO-эпохи на роллаут, 64 роллаута на обновление, 200 обновлений, $\approx 96$ минут на одном ядре CPU; полная настройка в \texttt{experiments/commnet\_ppo\_baseline/train\_ppo\_commnet.py}) достигает успеха $0.667$ на 100 отложенных сидах --- конкурентоспособно с символьной пировой архитектурой при $K{=}8$ ($0.65$).

\begin{table}[h]
\centering
\small
\caption{REINFORCE против PPO CommNet против символьного peer ($K{=}8$) на сценарии дефицита ($N{=}3$, $T{=}2$, асимметричная, опасность $0.05$). PPO конкурентоспособен по доле успеха, но всё равно демонстрирует Q-насыщение и коллапс R/M --- диагностическое утверждение верифицировано, а не постулировано.}
\label{p1:tab:ppo-commnet}
\begin{tabular}{l c c c}
\toprule
Метрика & CommNet REINFORCE & \textbf{CommNet PPO} & Симв. peer ($K{=}8$) \\
\midrule
Доля успеха                  & 0.19   & \textbf{0.667} & 0.65 \\
Частота $Q^\star$                & 1.00   & \textbf{1.00}  & 0.99 \\
Частота $R^\star$                & 0.50   & \textbf{0.997} & 0.69 \\
Частота $M^\star$                & 0.50   & \textbf{0.997} & 0.69 \\
$P(M^\star|R^\star)$          & 1.00   & \textbf{1.00}  & 0.69 \\
$P(C^\star|M^\star)$          & 0.42   & \textbf{0.67}  & 0.87 \\
R/M разделены?                & нет     & \textbf{нет}    & \textbf{да} \\
\bottomrule
\end{tabular}
\end{table}

Результат PPO превращает структурный аргумент из утверждения («бюджет обучения этого бы не изменил») в измерение: $Q^\star$ остаётся насыщенным на 1.00, $P(M^\star|R^\star)$ остаётся 1.00 (коллапс R/M) даже при конкурентоспособном успехе. Стадия Q вырождена по построению у любой политики, чей коммуникационный канал --- вещественный вектор с обученной линейной проекцией --- канал никогда не бывает информативно пустым, независимо от бюджета обучения, --- а без явного символьного интерфейса фиксации цели событие M не может отделиться от R. Только архитектуры с явным интерфейсом \emph{запроса/фиксации} (символьный стек из \S\ref{p1:sec:archs}) обнажают постадийно разделимые события Q/R/M/C.

\section{Команды воспроизводимости}
\label{p1:app:repro}

Подачу сопровождает анонимизированный публичный репозиторий. Артефакты основной статьи сгенерированы следующими скриптами; всюду предполагаются \texttt{PYTHONPATH=.} и виртуальное окружение Python 3 с \texttt{numpy}, \texttt{scipy}, \texttt{pandas}, \texttt{matplotlib}, \texttt{imageio}, \texttt{gymnasium} и \texttt{minigrid}.

\paragraph{Вычислительная инфраструктура.} Все эксперименты выполнялись на одной
потребительской машине: Apple Silicon (10 ядер CPU для свипов через
\texttt{ProcessPoolExecutor}), macOS (Darwin 25.4), Python 3.9.6.
Ключевые версии библиотек: \texttt{numpy} 1.x, \texttt{scipy},
\texttt{torch} (бейзлайн CommNet-PPO и VMAS), \texttt{vmas} 1.5.2,
\texttt{minigrid}/\texttt{gymnasium}. LLM-эксперименты используют локальные
модели через Ollama (\texttt{llama3.1:latest} 8B, \texttt{qwen3:4b}) при
температуре $0$ с детерминированным дисковым кэшированием ответов;
исследование внешних фреймворков дополнительно использует \texttt{openai} 2.44,
\texttt{langgraph} 0.6, \texttt{langchain-ollama} и
\texttt{pyautogen} 0.9 (классический AutoGen). GPU-кластер или платный API
не использовались; полный многоагентный блок воспроизводится за ${\approx}22$
минуты на 8 ядрах, а каждое LLM-исследование --- за минуты-часы локального
инференса.

\subsection*{Воспроизведение одной командой (многоагентный блок, $\approx 22$ минуты на 8 ядрах)}
\path{bash scripts/reproduce_paper.sh}

\subsection*{Чистая сюита Статьи 1}
\begin{itemize}[leftmargin=16pt]
\item smoke-сюита (все локальные не-LLM блоки): \path{python scripts/run_paper1_clean_experiments.py --mode smoke --skip_llm --out_dir tmp/paper1_clean_experiments_smoke}
\item полная локальная сюита (одноагентный, оракульный, многоагентный, CommNet, переносимость/шум): \path{python scripts/run_paper1_clean_experiments.py --mode full --skip_llm --workers 8 --out_dir tmp/paper1_clean_experiments_full}
\item свип бюджета шагов LLM TextNav (требует локальный Ollama), запуск один раз на бюджет $K'\in\{6,12,25\}$: \path{python -m experiments.llm_collective.run_model_sweep --models qwen3:4b llama3.1:latest --episodes 8 --step_limit <K'> --out_dir tmp/llm_steplimit_<K'>}
\end{itemize}

\subsection*{Один агент (\S\ref{p1:sec:validation}--\S\ref{p1:sec:oracle})}
\begin{itemize}[leftmargin=16pt]
\item статейный бенчмарк и сюита обученных бейзлайнов: \path{scripts/benchmark_symbolic_memory_article.py}
\item статейный анализ и генерация рисунков: \path{scripts/analyze_symbolic_memory_article.py}
\item синтетическая проверка корректности факторизации: \path{scripts/validate_qrmc_factorization.py}
\item бенчмарк оракульных постадийных интервенций: \path{scripts/benchmark_oracle_stage_interventions.py}
\item бенчмарк переноса на BabyAI: \path{scripts/compare_semnav_babyai.py}
\end{itemize}

\subsection*{Много агентов (\S\ref{p1:sec:multi-agent-results})}
\begin{itemize}[leftmargin=16pt]
\item Базовый батч --- peer при $K\in\{1,2,4,8,16\}$ \emph{и} три бескаденционные архитектуры (независимая, общая шина, центральный агрегатор), 40 сидов ($25{,}920$ прогонов $=16{,}200$ peer $+\,9{,}720$ бескаденционных), $\approx 18$ мин:
\path{python experiments/big_experiment/exp_cadence_phase.py --mode full --workers 8 --out_dir tmp/big_experiment_qrmc_40}
\item Батч расширенных каденций --- только peer при $K\in\{32,48,64\}$ ($9{,}720$ прогонов), $\approx 4$ мин:
\path{python experiments/big_experiment/exp_extra_K.py --workers 8 --out_dir tmp/big_experiment_extraK}
\item[] \emph{Замечание.} Эти два батча разбивают итог $35{,}640$ по \emph{диапазону каденций}; поархитектурное разбиение из Приложения~\ref{p1:app:design} ($25{,}920$ peer по всем восьми каденциям против $9{,}720$ бескаденционных) переиспользует те же два итога для \emph{разных} подмножеств (совпадение: $3$ расширенные каденции $\times\,3{,}240 = 3$ бескаденционные архитектуры $\times\,3{,}240$). Объединённый лог --- \path{tmp/paper1_clean_experiments_full/multiagent_cadence_full/runs.csv}.
\item Постадийный анализ Q/R/M/C и бутстреп-ДИ:
\path{python experiments/big_experiment/analyze_qrmc.py --runs_csv tmp/big_experiment_qrmc_40/runs.csv --out_dir tmp/big_experiment_qrmc_40}
\item Многоагентные оракульные интервенции:
\path{python experiments/big_experiment/exp_oracle_interventions.py --workers 8 --out_dir tmp/big_experiment_oracle}
\item Smoke-тесты фаз 1--4:
\path{python scripts/run_all_smokes.py --out_dir tmp/all_smokes}
\item Обученный PPO бейзлайн CommNet (Таблица~\ref{p1:tab:ppo-commnet}, $\approx 96$ мин на одном ядре CPU):
\path{python -m experiments.commnet_ppo_baseline.train_ppo_commnet --total_updates 200 --rollouts_per_update 64}
\item Оценка Q/R/M/C политики PPO на 100 отложенных сидах:
\path{python experiments/commnet_baseline/eval_with_qrmc.py --policy_path experiments/commnet_ppo_baseline/ppo_policy.pt --out_dir experiments/commnet_ppo_baseline --n_episodes 100}
\item Свип переносимости на субстрате MiniGrid (\S\ref{p1:sec:limitations}, 100 прогонов, $\approx 2$ мин):
\path{python -m experiments.minigrid_multiagent_wrapper.run_portability_sweep --out_dir tmp/minigrid_portability}
\end{itemize}

\section{Замечание об ассетах и лицензиях}
\label{p1:app:assets}

Эксперименты используют MiniGrid~\citep{MinigridMiniworld23} и BabyAI~\citep{chevalier2018babyai}, оба выпущены под лицензией MIT. Локально реализованные бейзлайны и скрипты запуска бенчмарков, введённые в этой статье, будут выпущены под лицензией MIT в составе пакета артефактов camera-ready.

\clearpage
\FloatBarrier
\FloatBarrier

\clearpage
\part{Семантический варп: завершение, обусловленное провенансом}
\noindent\textit{Навигация по чужой памяти как измеряемое событие: аннотация чужеродности, закон варп-дистанции и координационный протокол Warp Drive.}

\paragraph{Аннотация части.} \textbf{Задача.} В коллективной памяти для многоагентной навигации агент может принять обязательство по цели --- и достичь её, --- ни разу её не наблюдав: все подтверждающие свидетельства пришли от соседей. Народная культура называет это навигацией по чужой песне; мы называем это \emph{семантическим варпом} (semantic warp). Метрики, учитывающие только успех, и даже постадийные метрики этого не видят, поскольку ни одна существующая методика оценки не обусловливает завершение (completion) \emph{происхождением} (provenance) свидетельств, стоящих за принятым обязательством.
\textbf{Подход.} Мы определяем событие варпа $\wstar_\theta := \mstar \wedge (\varphi \ge \theta)$, где чужеродность (foreignness) $\varphi$ --- это доля соседей в массе свидетельств, \emph{взятая в тех же весах, которыми слияние отбирало кандидата}, и зафиксированная в момент лока (lock). Варп --- не пятая стадия Q/R/M/C, а аннотация, расщепляющая $P(\cstar|\mstar)$ на страты по происхождению. Поскольку символьные слияния разлагаются по источникам, мы получаем (i) контрфактическую атрибуцию ценности варпа через детерминированный повтор с маскированием чужих свидетельств, (ii) \emph{закон варп-дистанции} в замкнутой форме, выведенный из правил доверие$\times$устаревание (trust$\times$staleness), и (iii) децентрализованный протокол \emph{Warp Drive} (оптимистичный лок + резервация + анти-$\mstar$-откат) для патологии варп-коллизий.
\textbf{Результаты.} На развёртке дефицита из 2{,}400 эпизодов сырые варп-локи почти никогда не завершаются напрямую ($P(\cstar|\wstar)=0.004$--$0.020$ против $0.22$--$0.41$ для собственных локов; доверительные интервалы не пересекаются), однако чужие свидетельства всё же несут ценность как транспорт: до $51\%$ завершений собственных локов --- это варп-ассистированные цепочки. Закон дистанции предсказывает горизонты формирования намерения с точностью до такта, включая отложенную проверку $6/6$ с предсказаниями, зарегистрированными до исполнения. Warp Drive разворачивает компромисс: доля успеха растёт при каждой каденции (cadence) с непересекающимися интервалами, и быстрый обмен с протоколом обходит лучший базовый вариант с фиксированной каденцией ($0.614$--$0.620$ против $0.583$). Строгий межсессионный и межагентный варп воспроизведён на LLM-субстрате ($10/10$ планировок против $0/10$ без чужой памяти).
\textbf{Утверждение.} Мы утверждаем, что коллективную память следует оценивать не только по тому, что агент извлекает, но и по тому, \emph{чьё свидетельство вызвало его обязательство} и \emph{при каком временном и координационном контракте} это свидетельство остаётся применимым.
\textbf{Тезис.} Ценность коллективной памяти состоит не в самой информации, а в \emph{провенансе, горизонте допустимости и координационном контракте}, прикреплённых к ней.
\textbf{Границы.} Сеточные субстраты и одна небольшая текстовая LLM-демонстрация; вклад --- измерительный слой, фальсифицируемый закон и минимальный протокол, а не общая теория.

\FloatBarrier
\section{Введение}
\label{p2:sec:intro}

В культурах песенных троп путник, знающий песню, способен пересечь
местность, которой никогда не видел: навигационная компетентность
передаётся без физического посещения. Многоагентные системы с общей
памятью совершают тот же акт постоянно --- агент принимает
обязательство по источнику воды, целевой области или расположению
инструмента, известным лишь из широковещательной рассылки соседа, ---
однако ни одна стандартная методика оценки не отличает этот акт от
обычной навигации по памяти. Метрики успеха агрегируют поверх него;
постадийная диагностика вроде Q/R/M/C считает обязательство
($\mstar$) и завершение ($\cstar$), но не то, \emph{чьи свидетельства}
их подкрепили.

Эта статья делает такой акт полноправным измеряемым событием.
\emph{Семантический варп} $\wstar$ --- это материализация цели, масса
подтверждающих свидетельств которой (почти) целиком чужеродна, а
\emph{завершение варпа} $P(\cstar|\wstar)$ спрашивает, доводит ли
чужая песня действительно до места. Три свойства символьной
коллективной памяти делают конструкцию строгой, и ни одно из них не
выполняется для непрозрачных векторных памятей:

\begin{enumerate}[leftmargin=20pt,itemsep=2pt,topsep=2pt]
\item \textbf{Слияния разлагаются по источникам.} Чужеродность
  $\varphi\in[0,1]$ вычисляется из тех же по-источниковых весов
  (доверие, устаревание, поддержка), которыми само слияние
  ранжировало кандидата --- это прочтение решения, а не новая
  эвристика.
\item \textbf{Повторы детерминированы.} Повторный прогон эпизода с
  маскированием чужих вкладов на этапе слияния даёт
  \emph{контрфактический} выигрыш варпа для каждого эпизода, а не
  корреляцию.
\item \textbf{Правила обмена имеют замкнутую форму.} Затухание
  доверия и устаревания влечёт \emph{закон варп-дистанции}: явный
  горизонт, за которым чужие свидетельства уже не могут формировать
  намерение --- фальсифицируемый и фальсифицируемо различный у разных
  архитектур.
\end{enumerate}

Наши измерения затем вынуждают переосмысление. Сырой варп почти
бесполезен в точке завершения: при дефиците варп-локи завершаются на
порядок реже, чем локи на собственных свидетельствах
(\S\ref{p2:sec:raw}), потому что одновременные получатели одной и той же
рассылки соревнуются за одну цель, и большинство обречено проиграть.
Однако те же чужие свидетельства доказуемо ценны как
\emph{транспорт}: до половины всех завершений на собственных
свидетельствах --- это цепочки, начавшиеся с провалившегося
варп-лока, который довёз агента до зоны самоподтверждения
(\S\ref{p2:sec:taxi}). А лёгкий децентрализованный контракт ---
резервация, подсаженная на уже существующую рассылку, плюс
анти-$\mstar$-откат в духе симуляций Time
Warp~\cite{jefferson1985}
--- возвращает то, что теряет сырой варп: с протоколом \emph{быстрый}
обмен становится лучшим режимом по безусловной доле успеха
(\S\ref{p2:sec:drive}). Отсюда тезис в одном предложении:
\textbf{ценность коллективной памяти --- не информация, а
координационный контракт, прикреплённый к информации.}

\paragraph{Вклад работы.}
\begin{enumerate}[leftmargin=20pt,itemsep=2pt,topsep=2pt]
\item Первая метрика завершения, обусловленная происхождением
  свидетельств, для коллективной памяти: $\wstar_\theta$, $\varphi$,
  замороженная в момент лока, разложение
  $P(\cstar|\mstar)=P(\cstar|\mstar,\varphi{<}\theta)\,P(\varphi{<}\theta|\mstar)
  + P(\cstar|\wstar)\,P(\wstar|\mstar)$ и контрфактический выигрыш
  варпа через детерминированный повтор с маскированием источников
  (\S\ref{p2:sec:warp-def}--\S\ref{p2:sec:raw}).
\item \emph{Закон варп-дистанции}
  $\mathrm{age}_{\max}(\tau_i)=\alpha^{-1}\ln(\tau_i c/\tau)$,
  выведенный из собственных правил доверие$\times$устаревание самой
  памяти, подтверждённый с точностью до такта, включая отложенную
  проверку с зарегистрированными предсказаниями ($6/6$ точных, в том
  числе ячейка, где горизонт отсечки обязан сработать раньше
  экспоненциального гейта --- и срабатывает), и различие формы на
  уровне архитектур: у пиринговой памяти с фиксированной каденцией
  гейта \emph{нет} (\S\ref{p2:sec:law}).
\item Протокол \emph{Warp Drive}: оптимистичный варп-лок, резервации,
  переносимые на существующей каденции рассылки (без координатора), и
  анти-$\mstar$-откат с экспоненциальной выдержкой. Он превентивно
  разрешает известную патологию введённого в заблуждение агента,
  ограничивает задержку отката величиной ${\approx}K$ тактов и
  поднимает безусловную долю успеха выше лучшего базового варианта с
  фиксированной каденцией (\S\ref{p2:sec:drive}).
\item Строгий межсессионный, межагентный варп на LLM-субстрате:
  агент, никогда не наблюдавший комнату, решает там задачу
  исключительно по памяти другого агента из предыдущей сессии;
  воспроизведено на $10/10$ планировках против контроля $0/10$
  (\S\ref{p2:sec:llm}).
\end{enumerate}

\paragraph{Связь со статьями-компаньонами.}
Статья о фреймворке Q/R/M/C рассматривает постадийные показатели как
измерительный инструмент; статья о коллективной памяти устанавливает
архитектуры и компромисс каденции. Настоящая статья добавляет поверх
обеих ось происхождения: те же логи, одна новая аннотация и
механизмы, которые эта аннотация обнажает. % TODO: cite companions per venue policy.

\section{Постановка}
\label{p2:sec:setting}

\paragraph{Субстрат и задача.} $N$ агентов исследуют сетку
$12\times10$ с $M$ семантическими целями-водоёмами ($N>M$: дефицит) и
опасностями. Каждый агент поддерживает символьную память мест (теги,
уверенности, поддержка, свежесть). Четыре архитектуры обмена поверх
идентичного субстрата и планировщика: \method{independent} (без
коммуникации), \method{shared} (единая шина событий),
\method{centralized} (консенсусный агрегатор), \method{peer}
(рассылка каждые $K$ тактов, асимметричное доверие), плюс
\method{CSM} (peer $+$ явные правила слияния/доверия/устаревания со
скоростью затухания $\alpha{=}0.05$ и порогом включения
$\tau{=}0.30$). % TODO: compress substrate details; cite companion §2.

\paragraph{Инструментирование Q/R/M/C.} На каждом такте для каждого
агента логируются формирование запроса (Q), удовлетворение выборки
(R), материализация цели как лок планировщика ($\mstar$) и завершение
($\cstar$). Условные показатели уровня эпизода следуют определению из
статьи-компаньона. Все измерения варпа ниже --- аннотации этих
событий $\mstar$; ни один инвариант планировщика или памяти не
изменяется.

\section{Событие варпа}
\label{p2:sec:warp-def}

\paragraph{Происхождение и чужеродность.} Каждая единица
свидетельства, входящая в слияние, несёт свой источник. Для
кандидата $m$, на котором планировщик ставит лок,
\[
\varphi(m) \;=\;
\frac{\sum_{e\in\mathrm{supp}(m),\,\pi(e)\neq\mathrm{self}} w(e)}
     {\sum_{e\in\mathrm{supp}(m)} w(e)},
\]
где $w(e)$ --- \emph{собственные веса слияния}: $\text{trust}\cdot
\log(1{+}\text{support})$ для консенсуса peer/centralized,
$\text{trust}\cdot e^{-\alpha\,\text{age}}\cdot\text{conf}$ для CSM,
массы наблюдений для общей шины. Для архитектур peer и centralized
$\varphi$ восстанавливается точно из повкладовых полей, которые
слияние уже экспортирует; измерительный слой читает состояние и
никогда его не пишет.

\paragraph{Правило заморозки.} $\varphi$ фиксируется в момент лока.
Если агент позже подтверждает место собственными глазами, $\varphi$
\emph{не} пересчитывается --- иначе каждый успешный варп размывал бы
сам себя на подходе, и $P(\cstar|\wstar)$ была бы неизмерима по
построению.

\paragraph{События варпа и страты.} $\wstar_\theta := \mstar \wedge
(\varphi\ge\theta)$ при $\theta{=}1$ (строгий: цель ни разу не
наблюдалась самим агентом, что перекрёстно проверяется по журналу
наблюдений) и $\theta{=}0.8$ (мягкий). Доля варпа $P(\wstar|\mstar)$
--- структурный дескриптор (тождественно $0$ для
\method{independent}); завершение варпа $P(\cstar|\wstar)$ --- новая
величина. Разложение по правилу цепочки из \S\ref{p2:sec:intro} делает
варп точным уточнением $P(\cstar|\mstar)$, а не новой стадией:
идентифицируемость Q/R/M/C не затронута.

\paragraph{Регрессионный якорь.} Известный эпизод введённого в
заблуждение агента из сравнения архитектур (агент A ставит лок на
$(3,8)$, известную только от C; C занимает её первым; A застревает)
аннотируется инструментированием как строгий $\wstar$ с
$\varphi=1.0$ и коллизией --- патология, которую фреймворк обязан
видеть, а протокол из \S\ref{p2:sec:drive} обязан чинить.

\section{Сырой варп редко завершается --- но перевозит}
\label{p2:sec:raw}

\paragraph{Дизайн.} $2{,}400$ эпизодов: $(N,M)\in\{(3,2),(5,3),(8,5),
(4,4),(6,2)\}$ $\times$ три планировки $\times$
$\{\text{peer-}K2,\text{peer-}K8,\text{CSM-}K8,\text{shared}\}$
$\times$ 20 сидов $\times$ \{полный, с маскированием чужого\} при
идентичных сидах. Единица анализа: событие лока; доверительные
интервалы --- бутстреп с кластеризацией по эпизодам.

\begin{figure}[t]
\centering
\includegraphics[width=0.72\linewidth]{figures/fig_warp_strata}
\caption{Завершение, стратифицированное по $\varphi$, в зависимости
от каденции рассылки $K$ (выделенная развёртка peer: ячейки дефицита,
случайные планировки, 20 сидов, бутстреп-интервалы с кластеризацией
по эпизодам). По мере уменьшения $K$ доля варпа $P(\wstar|\mstar)$
растёт с $0.11$ до $0.55$, тогда как $P(\cstar|\wstar)$ остаётся ниже
$0.05$, а показатель собственной страты не меняется: коллапс
завершения при быстром обмене сосредоточен в страте $\wstar$, что
даёт механистическое прочтение сдвигу узкого места из
статьи-компаньона. Подскок доли при $K{=}16$ (рассылки редки, но и
сами агенты мало что открывают) приводится как наблюдённый.}
\label{p2:fig:strata}
\end{figure}

Рисунок~\ref{p2:fig:strata} показывает каденционный срез на выделенной
развёртке peer ($K\in\{1,2,4,8,16\}$): доля варпа --- это то, что
покупает быстрый обмен, и это ровно та страта, которая не
завершается. Таблица~\ref{p2:tab:hw1} даёт срез по архитектурам.

\begin{table}[h]
\centering
\caption{Страты завершения по происхождению при дефиците ($N>M$),
полные прогоны. Варп-локи завершаются на порядок реже, чем локи на
собственных свидетельствах; все пары доверительных интервалов не
пересекаются. (H-W1 подтверждена.)}
\label{p2:tab:hw1}
\small
\begin{tabular}{lcccc}
\toprule
Архитектура & $P(\cstar|\wstar_{0.8})$ & $n_{\wstar}$ & $P(\cstar|\mstar,\text{own})$ & $n_{\text{own}}$ \\
\midrule
peer-$K{=}2$ & 0.004 [0.001, 0.008] & 1{,}437 & 0.223 [0.186, 0.275] & 2{,}738 \\
peer-$K{=}8$ & 0.020 [0.000, 0.038] & 149    & 0.389 [0.329, 0.446] & 1{,}734 \\
CSM-$K{=}8$  & 0.017 [0.006, 0.044] & 637    & 0.414 [0.368, 0.453] & 1{,}684 \\
shared       & 0.004 [0.002, 0.008] & 2{,}099 & 0.279 [0.237, 0.332] & 2{,}250 \\
\bottomrule
\end{tabular}
\end{table}

\subsection{Варп как транспорт: цепочки локов}
\label{p2:sec:taxi}

Прямое завершение недооценивает ценность чужих свидетельств.
\emph{Варп-ассистированное собственное завершение} --- это
завершённый лок на собственных свидетельствах, которому (у того же
агента, в том же эпизоде) предшествовал сброшенный мягкий
$\wstar$-лок: варп провалился как пункт назначения, но удался как
поездка.

\begin{table}[h]
\centering
\caption{Анализ цепочек локов (случайные планировки, полные прогоны).
Первый $\wstar$-лок в цепочке ставится с расстояния
${\approx}8$--$10$ клеток; завершающий собственный лок --- с $2.0$
клеток, то есть ровно с радиуса наблюдения, точки
самоподтверждения.}
\label{p2:tab:taxi}
\small
\begin{tabular}{lccc}
\toprule
Ветвь & собств. $\cstar$ & варп-ассистированные (доля) & $r(\text{first }\wstar) \to r(\text{own})$ \\
\midrule
shared          & 248 & 127 (\textbf{51.2\%}) & $9.5 \to 2.0$ \\
peer-$K{=}2$    & 223 & 99 (\textbf{44.4\%})  & $9.7 \to 2.0$ \\
CSM-$K{=}8$     & 279 & 13 (4.7\%)            & $5.8 \to 2.0$ \\
peer-$K{=}8$    & 246 & 7 (2.8\%)             & $8.4 \to 2.0$ \\
\bottomrule
\end{tabular}
\end{table}

Это разрешает кажущийся парадокс контрфактических данных: общая шина
при умеренном дефиците выигрывает от чужих свидетельств $+10.8$ такта
$[+2.2,+20.9]$, хотя её прямое $P(\cstar|\wstar)$ равно $0.004$.
Выигрыш течёт через цепочки; правило заморозки честно записывает
завершение в собственную страту, а анализ цепочек восстанавливает
вклад варпа как отдельную измеренную величину.

\subsection{Контрфактический выигрыш варпа}
\label{p2:sec:gain}

Маскирование чужих свидетельств на этапе слияния при том же сиде даёт
$\mathrm{WG} = t_{\text{cens}}(\text{masked}) -
t_{\text{cens}}(\text{full})$ на эпизод. Карта знаков по (планировка,
дефицит $\rho{=}N/M$) зависит от режима: положительна для общей шины
на непредсказуемых планировках при умеренном дефиците ($+10.8$ при
$\rho{=}1.5$), отрицательна для быстрого пирингового обмена
(peer-$K2$, случайные планировки: $-7.1$ $[-13.5,-1.3]$),
${\approx}0$ при жёстком дефиците. Поагентная маска (один агент
маскирован, остальные обмениваются как обычно) подтверждает знак на
индивидуальном уровне: $\mathrm{WG}_i = -10.25$ $[-20.91,-0.16]$
такта на ячейке быстрого обмена при дефиците, при этом побочный
эффект на немаскированных агентах ${\approx}0$ ($-0.62$
$[-3.06,+1.91]$). Сырой варп \emph{вредит} своему получателю ровно
там, где обмен самый быстрый --- цена коллизий превышает ценность
информации. % TODO: full sign map to appendix.

\paragraph{Разведочный результат: давление коллизий.} Давление
со-локов (другие агенты, залоченные на ту же цель) коррелирует с
провалом варпа с предсказанным знаком, но пренебрежимой величиной
(Спирмен $+0.033$, $p{=}0.018$, $n{=}5{,}117$); мы приводим это как
разведочный результат, а не как подтверждённый предиктор.
% honest downgrade of pre-registered H-W5

\section{Закон варп-дистанции}
\label{p2:sec:law}

CSM включает чужое свидетельство с весом $\tau_i\, e^{-\alpha\,
\text{age}}\, c$ (доверие $\tau_i$, уверенность $c{=}0.95$) и порогом
включения $\tau{=}0.30$. Следовательно, чужое свидетельство перестаёт
формировать намерение при
\[
\mathrm{age}_{\max}(\tau_i) \;=\; \tfrac{1}{\alpha}\,
\ln\!\big(\tau_i\, c / \tau\big), \qquad \text{определено при } \tau_i c > \tau .
\]

\paragraph{Протокол «свидетель--путник».} Свидетель один раз
наблюдает воду и замолкает; путник, никогда её не видевший, получает
снимок памяти с параметризованным возрастом. Каждый лок путника ---
строгий $\wstar$ по построению. Рисунок~\ref{p2:fig:law} суммирует оба
эксперимента.

\begin{figure}[t]
\centering
\includegraphics[width=\linewidth]{figures/fig_warp_law}
\caption{Закон варп-дистанции. \textbf{(a)} Вес чужого свидетельства
в слиянии при затухании доверие$\times$устаревание; $\times$ отмечает
эмпирически измеренный последний возраст, при котором запрос путника
ещё обнажает цель --- каждый лежит на пересечении своей кривой с
порогом включения $\tau$. У пиринговой памяти с фиксированной
каденцией члена устаревания нет (пунктир): её гейт никогда не
закрывается; архитектуры различаются \emph{формой} кривой.
\textbf{(b)} Предсказанные и эмпирические точки излома завершения в
полных навигационных эпизодах: исходная сетка (кружки; возраст
дискретизован с шагом 2, поэтому точки лежат не более чем на один шаг
сетки ниже диагонали) и зарегистрированная отложенная проверка
(звёздочки; шаг возраста 1, предсказания записаны на диск до
исполнения, $6/6$ точных, включая never-ячейку в $(-1,-1)$ и ячейку,
где горизонт отсечки срабатывает раньше экспоненциального гейта).}
\label{p2:fig:law}
\end{figure}

\begin{table}[h]
\centering
\caption{Кривая гейта: наибольший возраст свидетельства, при котором
запрос путника ещё обнажает цель. У пиринговой памяти с фиксированной
каденцией члена устаревания нет: гейт никогда не закрывается
(проверено до возраста 120 при всех уровнях доверия) --- архитектуры
различаются \emph{формой} кривой, а не только её уровнем.}
\label{p2:tab:gate}
\small
\begin{tabular}{lcc}
\toprule
доверие $\tau_i$ & эмпирический $\mathrm{age}_{\max}$ & предсказание \\
\midrule
1.00 & 23 & 23.05 \\
0.80 & 18 & 18.59 \\
0.60 & 12 & 12.84 \\
0.40 & 4  & 4.73 \\
0.25 & гейт закрыт & 0 (trust-flip) \\
\bottomrule
\end{tabular}
\end{table}

\paragraph{Точка излома завершения.} В полных навигационных эпизодах
неравенство осуществимости получает одно дискретное уточнение: путник
самоподтверждает цель, как только она оказывается в радиусе
наблюдения, поэтому чужому гейту достаточно продержаться до такта
$t_{\text{gate}} = d - r_{\text{obs}} - 1$, что даёт точку излома
$\mathrm{bp} = \mathrm{age}_{\max} - t_{\text{gate}}$; отсечка
снимков исполняется на тактах рассылки и добавляет границу $10K -
t_{\text{bcast}}$. Все константы закона неизменны. Все 11 ячеек CSM
исходной сетки совпадают в пределах разрешения дискретизации; кривые
завершения --- монотонные ступени; trust-flip дискретно обнуляет
события $\wstar$.

\begin{table}[h]
\centering
\caption{Отложенная проверка с зарегистрированными предсказаниями:
шесть ячеек, отсутствующих в исходной сетке, предсказания записаны на
диск \emph{до} любого эпизода, шаг сетки возраста 1. Все шесть ---
точные целочисленные совпадения, включая never-ячейку и ячейку, где
горизонт отсечки обязан сработать раньше экспоненциального гейта.}
\label{p2:tab:holdout}
\small
\begin{tabular}{lccc}
\toprule
Ячейка & связывающая граница & предсказание & эмпирика \\
\midrule
$K{=}8$, $\tau_i{=}0.9$, $d{=}9$  & гейт    & 14 & \textbf{14} \\
$K{=}8$, $\tau_i{=}0.7$, $d{=}9$  & гейт    & 9  & \textbf{9}  \\
$K{=}8$, $\tau_i{=}0.5$, $d{=}15$ & гейт    & никогда & \textbf{никогда} \\
$K{=}4$, $\tau_i{=}1.0$, $d{=}6$  & гейт    & 20 & \textbf{20} \\
$K{=}2$, $\tau_i{=}1.0$, $d{=}12$ & \emph{отсечка} & 12 & \textbf{12} \\
$K{=}2$, $\tau_i{=}0.8$, $d{=}8$  & гейт    & 13 & \textbf{13} \\
\bottomrule
\end{tabular}
\end{table}

Закон, таким образом, предсказателен вне выборки с точностью до
такта: правила CSM действительно являются механизмом, который
формирует (и расформировывает) намерение. Он также переинтерпретирует
результат статьи-компаньона о том, что CSM доминирует над пиринговой
памятью с фиксированной каденцией: доверие$\times$устаревание --- это
\emph{варп-гейт} --- CSM не обменивается лучше, она \emph{отказывается
варпать по устаревшей или недоверенной песне}.

\paragraph{Переносимость на непрерывную динамику.} Все три
зарегистрированных предсказания выживают на непрерывном субстрате
VMAS через тот же мост «непрерывное$\to$сетка», что и в тесте
переносимости статьи-компаньона: (i) страты по происхождению остаются
разнесёнными на порядок ($P(\cstar|\wstar)=0.012$ $[0.000,0.042]$
против $0.562$ $[0.457,0.709]$ при $K{=}1$); (ii) доля варпа
монотонно падает с каденцией ($0.723\to0.366\to0.000$ по $K{=}1/4/16$);
(iii) закон становится \emph{метрическим}: при однократной калибровке
времени «путь до видимости» на эпизоде со свежим свидетельством
предсказанные точки излома совпадают в пределах сетки дискретизации,
а все ячейки за метрическим радиусом варпа
($v\cdot\mathrm{age}_{\max}$ короче расстояния) предсказаны как
\emph{никогда} и наблюдены как \emph{никогда}. Один по-настоящему
непрерывный вывод: инерция --- физический канал утечки варпа --- агент
может по инерции вкатиться в зону видимости после закрытия гейта;
устраняется контроллером «нет лока --- тормози», восстанавливающим
сеточную семантику.

\section{Протокол Warp Drive}
\label{p2:sec:drive}

Варп-коллизия --- лок на чужих свидетельствах по цели, которую займёт
кто-то другой, --- доминирующий отказ быстрого обмена. Warp Drive
добавляет три примитива, все децентрализованные:

\begin{enumerate}[leftmargin=20pt,itemsep=2pt,topsep=2pt]
\item \textbf{Оптимистичный варп-лок} --- локи на чужих
  свидетельствах остаются разрешены немедленно (планировщик не
  изменяется).
\item \textbf{Варп-резервация} --- на ближайшем такте рассылки
  владельца резервация едет на обычном снимке (та же каденция $K$,
  без координатора, без дополнительного канала); получатели подавляют
  зарезервированные цели в своих запросах. Приоритет: (такт лока,
  расстояние в момент лока, идентификатор агента).
\item \textbf{Анти-$\mstar$-откат} (семантика Time
  Warp~\cite{jefferson1985}) --- агент, чей лок доминирован чужой
  резервацией, отзывает его: планировщик перезапрашивает, цель уходит
  в экспоненциальную выдержку ($4\to8\to\cdots\to32$ тактов) с
  разрешением ничьих по идентификатору агента против осцилляций.
  Журнал событий остаётся только дополняемым (append-only).
\end{enumerate}

\paragraph{Проверка децентрализации.} Прогон протокола прямыми
поагентными вызовами --- без рантайма, без окружения --- даёт те же
решения об отзыве по одному лишь собственному входящему ящику каждого
агента.

\paragraph{Якорь «введённого в заблуждение A».} С Warp Drive
патологический лок разрешается \emph{превентивно}: резервация C
приходит до того, как A вообще залочится на $(3,8)$; A завершает на
другой цели. В развёртке задержка отката доминированных локов конечна
и составляет ${\approx}K$: $1.0/1.5/2.7$ такта при $K{=}1/2/4$
(базовый вариант: до конца эпизода).

\begin{figure}[t]
\centering
\includegraphics[width=\linewidth]{figures/fig_warp_drive}
\caption{Warp Drive на безусловных метриках (840 эпизодов;
бутстреп-интервалы). \textbf{(a)} Доля успеха растёт при каждой
каденции с непересекающимися интервалами, и каждая быстрая ветвь с
протоколом превосходит базовую «сладкую точку» $K{=}8$ (пунктир).
\textbf{(b)} Цензурированное время до успеха соответственно падает.
Утверждение о том, что протокол снимает главный аргумент против
быстрого обмена, опирается на эти безусловные величины, а не на
(чувствительные к знаменателю) условные показатели.}
\label{p2:fig:drive}
\end{figure}

\begin{table}[h]
\centering
\caption{Безусловные результаты, 840 эпизодов (ячейки дефицита
$\times$ случайные/асимметричные планировки $\times$ 20 сидов). Доля
успеха под Warp Drive растёт при каждой каденции с непересекающимися
интервалами, цензурированное время падает, и быстрый обмен с
протоколом обходит базовую «сладкую точку» $K{=}8$
(рисунок~\ref{p2:fig:drive}). Перепись локов показывает сжатие
знаменателя условных метрик --- поэтому утверждение сформулировано на
безусловных величинах.}
\label{p2:tab:wd}
\small
\begin{tabular}{lccc}
\toprule
Ветвь & доля успеха & $t_{\text{cens}}$ (провал $=120$) & перепись локов \\
\midrule
$K{=}1$ базовый & 0.592 [0.574, 0.608] & 55.3 [53.7, 57.2] & 0.986 \\
$K{=}1$ +WD  & \textbf{0.620} [0.609, 0.629] & 53.1 [52.0, 54.4] & 0.985 \\
$K{=}2$ базовый & 0.572 [0.549, 0.592] & 57.2 [55.1, 59.7] & 0.973 \\
$K{=}2$ +WD  & \textbf{0.614} [0.601, 0.624] & \textbf{53.2} [52.0, 54.6] & 0.926 \\
$K{=}4$ базовый & 0.577 [0.554, 0.598] & 57.3 [55.1, 59.8] & 0.944 \\
$K{=}4$ +WD  & \textbf{0.620} [0.609, 0.630] & \textbf{52.2} [50.9, 53.7] & 0.739 \\
$K{=}8$ базовый & 0.583 [0.558, 0.604] & 55.2 [52.9, 57.9] & 0.682 \\
\bottomrule
\end{tabular}
\end{table}

На условной оси восстановление измеряется относительно потолка
дефицита $\min(m,M)/m$ (никакой протокол не наколдует лишней воды):
$0.90 / 2.15 / 6.6$ восстановимого коллапса при $K{=}1/2/4$; значения
выше 1 возникают потому, что подавление по резервациям убирает
обречённые локи из популяции $\mstar$. Вместе с
таблицей~\ref{p2:tab:taxi} это замыкает дугу: сырой варп проваливается
на завершении (\S\ref{p2:sec:raw}); с прикреплённым координационным
контрактом варп на быстрой каденции становится лучшим режимом.
«Сладкая точка» каденции из статьи-компаньона переинтерпретируется:
умеренный $K$ никогда не был про свежесть информации --- это был
неявный ограничитель частоты варп-коллизий, и явный контракт
справляется с этой работой лучше.

\section{Межсессионный варп на LLM-субстрате}
\label{p2:sec:llm}

Наконец, строгий варп пересекает одновременно границу агента и
границу сессии на языковом субстрате. В сессии 1 агент A (локальный
LLM-экстрактор тегов над наблюдениями на естественном языке; Ollama
\method{llama3.1}, детерминированное кеширование, нулевая стоимость
API) обходит трёхкомнатный текстовый дом, и его символьная память
выгружается на диск. В сессии 2 агент B --- свежий агент в свежем
экземпляре окружения, из которого удалён штатный направляющий оракул,
так что дальнее знание может прийти только из памяти, --- загружает
дамп A как чужое происхождение и решает задачу «принеси яблоко на
стол». Каждый лок B на яблоке или столе --- строгий $\wstar$ с
$\varphi{=}1.0$; исполнение слоистое (LLM принимает семантические
решения --- какое запомненное место преследовать, когда взять/положить,
--- а перемещение к залоченной путевой точке выполняет
детерминированный моторный примитив, в соответствии со стеком
планировщика субстрата).

\begin{table}[h]
\centering
\caption{Межсессионный варп на 10 рандомизированных планировках.
Завершение засчитывается на радиусе аффорданса окружения (действие
взять/положить в пределах манхэттенского расстояния 1, и свидетель
помечает якоря на достижимых клетках).}
\label{p2:tab:llm}
\small
\begin{tabular}{lccc}
\toprule
 & успех & среднее $t_{\text{succ}}$ & планировки с завершённым строгим $\wstar$ \\
\midrule
варп (память A) & \textbf{10/10} & 8.3 & \textbf{10/10} \\
контроль (пустая память) & 0/10 & --- & --- \\
\bottomrule
\end{tabular}
\end{table}

Насколько нам известно, это первый задокументированный межсессионный,
межагентный строгий семантический варп на LLM-субстрате. Контраст
$10/10$ против $0/10$ показывает, что варп несущий, а не
декоративный.

\paragraph{Масштабирование до живого коллектива.} При $N{=}3$
LLM-агентах, обменивающихся пиринговой рассылкой в условиях дефицита
($M{=}2$ яблока, четыре планировки, $K\in\{2,8\}$), вся
феноменология варпа воспроизводится на языковом субстрате (все
предсказания зарегистрированы до прогонов): строгие $\wstar$-локи
происходят; доля варпа падает с каденцией ($0.417$ при $K{=}2$ против
$0.083$ при $K{=}8$ --- та же форма, что и на сетке
$0.554\!\to\!0.105$ и в VMAS $0.723\!\to\!0.000$); варп-коллизии
появляются в $7/8$ эпизодов. Парное сравнение с базовым вариантом без
коммуникации на идентичных планировках показывает, что сырой обмен
\emph{снижает} успех с $2/3$ до $1/3$ на эпизод --- цена коллизий
превышает ценность информации на LLM-субстрате ровно как в
\S\ref{p2:sec:gain}. Прикрепление неизменённого протокола Warp Drive к
LLM-коллективу замыкает дугу и на языке: суммарный успех растёт с
$0.333$ до $0.417$ на обеих каденциях, а задержки отката снова
${\approx}K$ ($1$ такт при $K{=}2$, $7$ при $K{=}8$). Восстановление
частично (${\approx}25\%$ коллапса против $90$--$100\%$ на сетке): на
языковом субстрате патология коллизий усугубляется потерями
семантического уровня (галлюцинированные якорные теги, шумные локи),
которые никакая резервация не лечит --- измеримое разложение мы
оставляем будущей работе. Закон дистанции также выдерживает
свидетельства, произведённые LLM: при константе гейта $c$, равной
уверенности, которую экстрактор реально выдал
($c_{\text{meas}}{=}0.900$), все четыре зарегистрированные точки
излома ложатся в пределах одного такта от предсказания.

Два вывода об интерфейсе переносятся на любой дизайн
«LLM поверх символьной памяти»: формирователь запросов на LLM свободно
генерирует поисковые теги вне словаря памяти (что требует гейта
any-of), а малые локальные модели надёжны в семантическом
обязательстве, но не в координатной арифметике и не в систематическом
исследовании (что требует описанного выше слоистого исполнения и
явного яруса исследования --- без него открытие не происходит, и варпу
нечем питаться).

\begin{figure}[t]
\centering
\includegraphics[width=0.62\linewidth]{figures/fig_warp_universality}
\caption{Одна величина, три субстрата: доля варпа $P(\wstar|\mstar)$
в зависимости от каденции рассылки $K$ на символьной сетке, в
непрерывном сценарии VMAS и в текстовом мире с LLM. Кривые почти
совпадают --- то, какая часть материализаций коллектива работает на
чужих свидетельствах, определяется каденцией обмена, а не субстратом.
(Ветвь LLM: 8 эпизодов, только точечные оценки.)}
\label{p2:fig:universality}
\end{figure}

Рисунок~\ref{p2:fig:universality} накладывает кривую «доля
варпа--каденция» по всем трём субстратам: величина ведёт себя как
свойство архитектуры обмена, свободное от субстрата.

\section{Связанные работы}
\label{p2:sec:related}

\paragraph{Общая память в коллективах агентов.} Generative
Agents~\cite{park2023generative} и многоагентные LLM-фреймворки
разделяют потоки памяти или контекст, но не обнажают ни интерфейса
лока, ни по-источникового разложения слитого убеждения, так что
завершение, обусловленное происхождением, там неизмеримо. Векторные
памяти (например, MemGPT~\cite{packer2023memgpt}) сливают
арифметикой вложений; такое слияние не разлагается по источникам, что
блокирует и $\varphi$, и контрфактический повтор с маскированием.
\paragraph{Эмерджентная коммуникация.} CommNet~\cite{sukhbaatar2016commnet}
и TarMAC~\cite{das2019tarmac} обучают непрерывные каналы;
статья-компаньон показывает, что такие каналы схлопывают
наблюдаемость стадий. Варп требует явного события материализации,
которого эти архитектуры не обнажают.
\paragraph{Стигмергия.} Координация в стиле феромонов --- это
\emph{анонимный} варп: свидетельства передаются без происхождения.
Наши результаты количественно показывают, чего стоит анонимность
(доля варпа на общей шине максимальна при минимальном завершении
варпа) и что покупает происхождение (гейты доверия, резервации,
откат).
\paragraph{Оптимистичная конкурентность.} Примитив отката --- это
Time Warp~\cite{jefferson1985}, пересаженный из параллельной
симуляции в навигацию по памяти: локи --- оптимистичные вычисления,
свидетельство занятости --- запоздавшее сообщение (straggler),
анти-$\mstar$ --- анти-сообщение. Насколько нам известно, эта
пересадка нова.
% TODO: add ALFWorld/TextWorld refs for the LLM substrate; MiniGrid ref.

\section{Ограничения и честные негативные результаты}
\label{p2:sec:limits}

(0) Все эксперименты предполагают ОБЩУЮ систему координат: тождество
мест (place identity) --- это близость координат, а семантика служит
лишь разрешению ничьих внутри пространственного радиуса; проблема
межагентного соответствия мест вынесена за скобки, а не решена.
(Расширение без общей системы координат --- семантическое тождество
по отпечаткам с восстановлением системы координат через консенсус по
ориентирам, отказывающее закрытым образом (fail closed) при
перцептивном алиасинге там, где координатное допущение отказывает
открытым образом (fail open), --- реализовано и провалидировано
отдельно и выходит за рамки настоящей работы.) (1) Сеточные миры,
один непрерывный сценарий VMAS и один небольшой текстовый мир; без
фотореалистичного субстрата. (2) Предиктор давления со-локов выжил
только как разведочный (\S\ref{p2:sec:gain}). (3) Карта знаков выигрыша
варпа зависит от режима, а не имеет чистой внутренней структуры,
которую ожидала предрегистрированная гипотеза; детерминированные
симметричные/асимметричные планировки несут малую дисперсию по сидам,
так что статистический вес лежит на случайных планировках. (4)
Уточнение закона дистанции (горизонт самоподтверждения, расписание
отсечки) было сформулировано после просмотра исходной сетки ---
смягчено, но не стёрто зарегистрированной отложенной проверкой. (5)
LLM-демонстрация использует одно семейство моделей и намеренно малое
окружение; это демонстрация измеримости, а не бенчмарк. (6) Все
отклонения от предрегистрированного дизайна перечислены в
приложении~\ref{p2:app:deviations}.

\section{Заключение}

Варп превращает «агенты разделяют память» из свойства архитектуры в
измеренное, каузально атрибутируемое событие с фальсифицируемым
законом дистанции и протоколом, который на него опирается. Измерения
растворяют информационно-центричное прочтение коллективной памяти:
сырые чужие свидетельства чаще всего подводят получателя на
завершении и окупаются как транспорт, тогда как двухстрочный
координационный контракт --- резервируй при локе, откатывайся при
доминировании --- делает самый быстрый режим обмена одновременно и
лучшим. Управляющая поверхность, которая имеет значение, --- это
происхождение, а не объём.

\paragraph{Взгляд вперёд: контракт как одна функция из нескольких.}
Контрактная перспектива подсказывает прочтение, которое мы формулируем
как фальсифицируемую программу, а не как вывод: память --- не один
механизм, а несколько функций с различимыми сигнатурами отказов.
Изученный здесь гейт доверие$\times$давность реализует
\emph{память-для-доверия} (отключите его --- и завершение рушится
через устаревшие и отравленные коммиты; переброс доверия и негативный
результат сырого обмена --- ровно эта сигнатура); она отлична от
\emph{памяти-для-возврата} (собственная консолидация; её абляция
обескровливает генерацию кандидатов) и от \emph{памяти-для-рассказа}
(что вообще стоит транслировать; её абляция оставляет retrieval
получателей неулучшенным на недосэмплированных ими тегах). При таком
прочтении правильный критерий валидности разделяемой памяти --- не
прямая верификация содержимого, а \emph{экологическая пригодность}:
способно ли обязательство, сформированное на чужих свидетельствах,
быть доведённым до завершения под контрактом --- а это ровно та
величина, которую делают измеримой $P(\cstar \mid \wstar)$ и её закон.

\FloatBarrier
\begin{center}\rule{0.55\linewidth}{0.4pt}\\[2pt]\textbf{Приложения к этой части}\end{center}


\section{Отклонения от предрегистрированного дизайна}
\label{p2:app:deviations}
\begin{enumerate}[leftmargin=20pt,itemsep=2pt,topsep=2pt]
\item Разрешение ничьих в резервациях расширено до (такт лока,
  \emph{расстояние в момент лока}, идентификатор агента): чистое
  разрешение по идентификатору отдаёт одновременные локи дальнему
  агенту вместо ближнего первооткрывателя цели.
\item Восстановление протокола измеряется относительно потолка
  дефицита $\min(m,M)/m$, а не эталона $K{=}8$, что отделяет
  структурный коллапс от координационного; безусловные метрики
  приводятся рядом (таблица~\ref{p2:tab:wd}).
\item Точки излома завершения используют дискретно-временное
  уточнение из \S\ref{p2:sec:law} (константы закона неизменны;
  отложенная проверка зарегистрирована).
\item Давление коллизий операционализировано как число со-локов, а не
  число со-получателей (структурно постоянное $N{-}2$ при
  широковещании).
\item LLM-демонстрация: слоистое исполнение (семантика --- LLM,
  перемещение детерминированное); завершение засчитывается на радиусе
  аффорданса.
\end{enumerate}

\section{Воспроизводимость}
\label{p2:app:repro}
Все эксперименты выполняются на CPU за минуты--десятки минут;
LLM-демонстрация использует локальную модель Ollama с
детерминированным кешированием.
\begin{verbatim}
PYTHONPATH=. python experiments/warp/exp_warp_anchor.py
PYTHONPATH=. python experiments/warp/exp_warp_gain.py --seeds 20
PYTHONPATH=. python experiments/warp/analyze_warp_chains.py
PYTHONPATH=. python experiments/warp/exp_warp_gain_individual.py
PYTHONPATH=. python experiments/warp/exp_warp_age_law.py
PYTHONPATH=. python experiments/warp/exp_warp_age_law_holdout.py
PYTHONPATH=. python experiments/warp/exp_warp_drive.py --seeds 20
PYTHONPATH=. python experiments/warp/exp_warp_cross_session.py --layouts 10
\end{verbatim}
% TODO: anonymised artifact link for submission.

\section{Полная карта знаков контрфактического выигрыша}
\label{p2:app:signmap}
% TODO: import the 60-cell sign map from tmp/warp/w1_gain/w1_summary.json
% as a table; include the H-W2 discussion of regime boundaries.
\FloatBarrier

\clearpage
\part{Песня, а не точка: маршруты и смысловая идентичность}
\noindent\textit{Между агентами переносится песня: маршруты с провенансом рёбер, идентичность мест без общих систем координат, абляция словаря и анатомия песни.}

\paragraph{Аннотация части.} \textbf{Проблема.} Сопутствующая статья о варпе измеряет, способен ли
агент достичь цели, известной лишь из памяти напарника. В этой
конструкции --- и, по существу, во всех мультиагентных системах с общей
памятью --- скрыты два допущения. Во-первых, переносится \emph{точка}:
широковещательные снимки несут сводки мест, но не рёбра, поэтому путник
(traveler) узнаёт ГДЕ и вынужден заново выводить КАК. Во-вторых,
идентичность места задаётся \emph{общей системой координат}: <<то же
место>> означает <<близкие $(x,y)$ в сетке, которую среда выдаёт каждому
агенту>>; семантика лишь разрешает конфликты внутри пространственного
радиуса. В исходной культуре сонглайн устроен иначе: песня И ЕСТЬ
маршрут, а её места опознаются по тому, как они выглядят, а не по
координатам, которых никто не разделяет.
\textbf{Подход.} Мы снимаем оба допущения и измеряем, что каждое из них
скрывало. (i) \emph{Маршрутный варп (route warp)}: широковещания несут
пройденные рёбра с провенансом; $\phiroute$ --- чужая доля массы рёбер
вдоль зафиксированного пути, а строгий $\rwstar$ --- связный чужой
под-путь, по которому путник никогда не ходил. (ii) \emph{Смысловая
идентичность мест}: места сопоставляются по локальным семантическим
констелляциям (трансляционно-инвариантным отпечаткам мест,
fingerprints); межагентная система отсчёта --- сначала сдвиг, затем
дискретная подгруппа $\mathbb{Z}^2\rtimes C_4$ группы $SE(2)$ (сдвиги и повороты на четверть оборота) --- восстанавливается по взаимно-уникальным
совпадениям совместно посещённых ориентиров, и чужие свидетельства,
включая никогда не виденные цели, переносятся через неё.
\textbf{Результаты.} То, что даёт знание маршрута, --- ступенька, а не
наклон: перенос мест (place transfer) рушится на ПЕРВОЙ же стене (успех
$1.0$ при факторе обхода (detour factor) $D{=}1.00$ против $0.0$ при
$D{=}1.38$ --- хуже слепого исследования), тогда как перенос маршрута
завершает задачу всюду. Старение разрывает перенесённый маршрут в точке,
предсказуемой с точностью до конкретного ребра ($110/110$ режимов; все
$40$ разрывов (ruptures) в середине маршрута в двух классах миров --- на
предсказанном индексе, ошибка $0$), а безопасность едет вместе с
маршрутом: $0$ попаданий в опасности на пути свидетеля (witness) против
$3.25$ $[2.45,4.10]$ у путника с одними местами, причём риск
ВОЗОБНОВЛЯЕТСЯ в предсказанной точке разрыва. Смысловая идентичность
восстанавливает секретные системы отсчёта точно ($100\%$ для сдвигов;
$80/80$ при $SE(2)$), возвращает $88\%$ производительности общей системы
координат в полном peer-стеке и при перцептивном алиасинге отказывает
\emph{закрытым образом} (fail closed) --- там, где координатное
допущение отказывает \emph{открытым} (fail open), --- а сопутствующий
закон расстояния воспроизводится с точностью до тика (точки излома
$20/14/9$) без общего начала координат и без общего севера. Абляция
словаря ориентиров показывает, что механизм идентичности
\emph{безопасен в своём алфавите признаков}: обеднение словаря вплоть
до одного класса тегов стоит покрытия (точное восстановление $1.00 \to
0.83$), но никогда --- корректности (ноль fail-open переносов в $252$
ячейках). Наконец, эксперимент по анатомии песни разделяет роли внутри
самой песни: чисто реляционная песня (сигнатуры ориентиров $+$
такт-векторы, \emph{вовсе без глобальных координат}) переносит никогда
не виденную цель точно за $6.4\%$ битовой стоимости снапшота, тогда
как одни такты отказывают \emph{открытым образом} ($11/11$ неверных
захватов при рассогласовании якорей), а одни сигнатуры вовсе не могут
перенести невиданное.
\textbf{Утверждение.} Мы утверждаем, что переносимая единица навигационной памяти --- не координата цели и не изолированная запись места, а реляционный маршрут с провенансом: ориентиры задают идентичность, рёбра дают перенос, а их композиция позволяет агенту исполнить чужой опыт без общей системы координат.
\textbf{Тезис.} Описание маршрута без идентичности небезопасно, а идентичность места без переноса инертна; исполнимая память требует их композиции. Коллективную навигацию делает работающей не общая карта.
Переносится песня --- последовательность мест, заякоренных смыслом, ---
и каждое обеднение песни (точка вместо маршрута, координаты вместо
смысла) отказывает специфическим, измеримым и теперь предсказанным
способом.
\textbf{Границы.} Сеточные миры; сдвиги и повороты на $90^\circ$;
измерение и минимальные механизмы, а не общая теория.

\FloatBarrier
\section{Введение}
\label{p3:sec:intro}

Сонглайн кодирует путешествие: последовательность мест, каждое из
которых опознаётся по тому, что там находится, связанных в путь
порядком песни. Когда мультиагентные системы <<делят память>>, они
обычно делят нечто куда более бедное --- точку (pin) на общей карте:
координаты цели в системе отсчёта, в которой, как предполагается,
живут все агенты. Сопутствующая статья о варпе\footnote{Сопутствующие
материалы: измерительный фреймворк Q/R/M/C, архитектурное исследование
коллективной семантической памяти и статья об измерении семантического
варпа. Ссылки отложены согласно политике одновременной подачи.} сделала
навигацию по чужой памяти измеряемым событием ($\wstar$: материализация,
чьи свидетельства целиком чужие) и показала, что её ценность управляется
координационным контрактом, приложенным к информации. Но её конструкция
молча унаследовала оба обеднения:

\begin{enumerate}[leftmargin=20pt,itemsep=2pt,topsep=2pt]
\item \textbf{Точка, а не песня.} Широковещательные снимки несут
  концепты мест --- центроиды, теги, уверенности, --- но не рёбра.
  Получатель узнаёт, ГДЕ вода; маршрут выводится заново с нуля. В
  открытых сетках это незаметно: манхэттенский путь и есть маршрут.
  Но как только пути становятся дорогими, это решает всё.
\item \textbf{Координаты, а не смысл.} <<То же место>> смешивает два
  вопроса, которые серия до сих пор не разделяла. \emph{Семантический
  предикат} --- тот же ли это КЛАСС желаемого объекта (вода, GOOD)? ---
  едет в каждом снимке, и именно над ним определено удовлетворение
  задаче. \emph{Соответствие сущностей} (entity correspondence) --- тот
  же ли это референт в мире? --- операционализировано как координатная
  близость, а семантике оставлено лишь разрешение конфликтов внутри
  радиуса. Межагентная задача соответствия мест --- причина, по которой
  слияние карт трудно в робототехнике, --- вынесена за скобки указом.
\end{enumerate}

Эта статья снимает оба допущения, по одному за раз, с измерительной
дисциплиной всей серии: каждое утверждение регистрируется до запусков,
отказы относительно регистрации сообщаются вместе с их механизмами, а
сильнейшие результаты --- точные совпадения с предсказаниями в
замкнутой форме.

\paragraph{Вклад.}
\begin{enumerate}[leftmargin=20pt,itemsep=2pt,topsep=2pt]
\item \textbf{Маршрутный варп} ($\rwstar$, $\phiroute$): провенанс на
  уровне рёбер и строгое событие маршрутного варпа --- связный чужой
  под-путь, по которому путник никогда не ходил (\S\ref{p3:sec:route}).
  Преимущество переноса маршрута над переносом мест --- \emph{обрыв},
  а не наклон: перенос одних мест умирает на первой же стене (успех
  $1.0 \to 0.0$ между измеренными факторами обхода $1.00$ и $1.38$) и
  там ХУЖЕ полного отсутствия переноса, тогда как перенос маршрута
  завершает задачу всюду (\S\ref{p3:sec:cliff}).
\item \textbf{Закон разрыва}: свидетель штампует рёбра последовательно,
  поэтому его след стареет неравномерно, и запоздавший путник гонится
  за угасающей песней. Проигрывание гейта доверие$\times$давность серии
  по калиброванной кинематике предсказывает, \emph{где рвётся песня},
  с точностью до конкретного ребра: $110/110$ совпадений режимов, все
  $21$ разрыв в середине маршрута --- на предсказанном индексе (ошибка
  $0$), плюс дискретный переброс доверия, полностью обнуляющий
  фиксации маршрутов (\S\ref{p3:sec:rupture}).
\item \textbf{Стратификация по опасностям}: оси стен и опасностей
  независимы --- стены ломают \emph{завершение} переноса мест,
  опасности ломают его \emph{безопасность} ($0$ против $3.25$
  попаданий при равном завершении и $+0.7$ тика), а когда песня
  рвётся, риск возобновляется на предсказанном ребре
  (\S\ref{p3:sec:hazard}).
\item \textbf{Смысловая идентичность мест}: трансляционно-инвариантные
  семантические отпечатки, взаимно-уникальное сопоставление ориентиров
  и консенсусное восстановление системы отсчёта --- сначала сдвиг,
  затем дискретный случай $\mathbb{Z}^2\rtimes C_4$ с секретными поворотами на четверть оборота
  (\S\ref{p3:sec:identity}). Точное восстановление системы отсчёта
  ($100\%$; $80/80$ при $SE(2)$); онлайн-выравнивание в полном
  peer-стеке, возвращающее $88\%$ разрыва с общей системой координат;
  доказанная асимметрия отказов (смысловая идентичность отказывает
  \emph{закрытым образом} при алиасинге и при принципиально
  невосстановимой симметрии $180^\circ$; координатная --- отказывает
  \emph{открытым}); и сопутствующий закон расстояния, воспроизведённый
  с точностью до тика ($20/14/9$) вовсе без общей системы отсчёта
  (\S\ref{p3:sec:stack}). Зарегистрированная абляция словаря ориентиров
  показывает, что механизм \emph{безопасен в своём алфавите
  признаков}: вплоть до одного содержательного класса обеднение
  превращает восстановление в отказ --- и никогда в неверный захват
  ($0$ fail-open переносов в $252$ ячейках; \S\ref{p3:sec:ablation}).
  Эксперимент по анатомии песни затем разделяет роли внутри песни:
  \emph{узлы несут идентичность, рёбра несут перенос} --- одни такты
  отказывают открытым образом, одни сигнатуры не достигают
  невиданного, а их композиция воспроизводит перенос полного снапшота
  за $6.4\%$ бит (\S\ref{p3:sec:anatomy}).
\end{enumerate}

\section{Постановка}
\label{p3:sec:setting}

Сеточные миры (от $12{\times}10$ до $14{\times}12$) с семантическими
целями-водой, опасностями и стенами; агенты с символической памятью
мест, peer-широковещанием с каденцией $K$ и стандартным планировщиком
серии (захват цели $=$ событие $\mstar$ фреймворка Q/R/M/C; завершение
--- $\cstar$). Всюду используется правило включения
доверие$\times$давность из сопутствующей работы о коллективной памяти
с неизменными константами: масса свидетельства
$\text{trust}\cdot e^{-\alpha\,\text{age}}\cdot m$, порог
$\tau{=}0.30$, $\alpha{=}0.05$. События варпа следуют сопутствующей
статье о варпе: провенанс $\varphi$ замораживается в момент захвата;
завершение атрибутируется непрерывному захвату от чужой фиксации
($\varphi{=}1$ при первом захвате) до успеха.
% TODO: compress; cross-reference companions.

\section{Маршрутный варп: перенос КАК}
\label{p3:sec:route}

\paragraph{Провенанс рёбер.} Широковещательные снимки дополнительно
несут СОБСТВЕННЫЕ пройденные рёбра отправителя
$(a,b,\text{count},\text{last tick})$; чужие рёбра никогда не
ретранслируются (тот же контракт приватности, что и у снимков мест: ни
один агент не читает чужой граф, только транспорт). Для пути $P$,
который фиксирует планировщик,
\[
\phiroute(P) \;=\;
\frac{\sum_{e\in P,\;\pi(e)\neq\mathrm{self}} w(e)}
     {\sum_{e\in P} w(e)},
\qquad w(e) = \text{trust}\cdot e^{-\alpha\,\text{age}(e)}\cdot
\log(1{+}\text{count}(e)),
\]
--- те же веса, которые использует маршрутный гейт. \textbf{Строгий
$\rwstar$} --- это зафиксированный, пройденный, связный под-путь из
$\ge 5$ рёбер с $\phiroute = 1$: весь куплет --- чужой. Варп мест ---
вырожденный случай (длина пути $0$: чужая лишь цель), что даёт
двумерную плоскость провенанса $(\varphi_{\mathrm{place}},
\phiroute)$.

\paragraph{Первый строгий $\rwstar$.} На лабиринте-гребёнке с
измеренным фактором обхода $D{=}2.65$ путник фиксирует и проходит
полностью чужой маршрут из $45$ рёбер, завершая задачу за $54$ тика; с
теми же чужими свидетельствами о МЕСТАХ, но без рёбер стандартный
планировщик так и не приходит (бюджет $300$ тиков), а слепое
исследование занимает $126$.

\subsection{Обрыв: что даёт знание маршрута}
\label{p3:sec:cliff}

\begin{figure}[t]
\centering
\includegraphics[width=\linewidth]{figures/fig_route_cliff}
\caption{Обрыв (парные плечи на идентичных лабиринтах, 20 сидов на
корзину измеренного обхода). \textbf{(a)} Перенос одних мест не
деградирует с ростом обхода --- он рушится на ПЕРВОЙ же стене,
локализованной пробной корзиной между $D{=}1.00$ и $D{=}1.38$; перенос
маршрута завершает задачу всюду. \textbf{(b)} Ниже обрыва знать ГДЕ
без КАК --- \emph{хуже, чем не знать ничего}: манхэттенская тяга к
известной цели прижимает планировщик к стенам (от $-150$ до $-168$
тиков против слепого исследования), тогда как в открытом мире то же
знание стоит $+57$.}
\label{p3:fig:cliff}
\end{figure}

\begin{table}[h]
\centering
\caption{Три плеча на идентичных лабиринтах (спарены по сиду),
разбитые по ИЗМЕРЕННОМУ фактору обхода $D$; 20 сидов на корзину,
цензурированные выигрыши (отказ $=300$). Перенос одних мест не
деградирует с ростом $D$ --- он рушится на первой стене и там хуже
слепого исследования (манхэттенская тяга к известной цели прижимает
планировщик к стенам). Перенос маршрута завершает задачу всюду; каждое
завершение --- строгий $\rwstar$ с $\phiroute=1$.}
\label{p3:tab:cliff}
\small
\begin{tabular}{lcccc}
\toprule
$D$ & выигрыш маршрут--место [CI] & выигрыш маршрут--слепой [CI] & успех r/p/b \\
\midrule
1.00 & $+4.7$ $[+3.4,+6.0]$ & $+62$ $[+44,+80]$ & 20/\textbf{20}/20 \\
1.38 & $+281$ $[+279,+283]$ & $+113$ $[+77,+155]$ & 20/\textbf{0}/16 \\
2.03 & $+273$ $[+271,+275]$ & $+123$ $[+91,+160]$ & 20/\textbf{0}/17 \\
3.32 & $+262$ $[+259,+264]$ & $+110$ $[+72,+152]$ & 20/\textbf{0}/15 \\
\bottomrule
\end{tabular}
\end{table}

Рисунок~\ref{p3:fig:cliff} показывает обе стороны результата.
Два предзарегистрированных предсказания здесь провалились ---
поучительным образом. Отрицательный контроль при $D{=}1$ был
зарегистрирован в тиках ($|\text{gain}|\le 3$) и дал $+4.7$: зазор ---
это \emph{экономия поворотов} (маршруты BFS идут прямыми отрезками;
планировщик доминирующей оси чередует оси), а не знание маршрута ---
по ДЛИНЕ пути зазор в точности $0.0$ $[0.0,0.0]$ во всех 20 ячейках
открытой сетки. А утверждение <<выигрыш растёт с $D$>> неизмеримо: при
нулевом успехе переноса мест выигрыш насыщается на потолке
цензурирования. Правильная форма закона --- ступенька, локализованная
пробной корзиной в $D\in(1.00,1.38)$, то есть на первой стене.
% re-registration documented in the artifact.

\section{Закон разрыва: где рвётся песня}
\label{p3:sec:rupture}

Свидетель штампует ребро $i$ на тике $i$: его след стареет
неравномерно --- начало маршрута всегда старше конца, --- и путник,
запоздавший на $a_0$, \emph{гонится за угасающей песней}. Под
стандартным гейтом связывающим является ТЕКУЩЕЕ ребро (старейшее из
оставшихся), что даёт три режима по оси $a_0$: завершение; разрыв на
предсказуемом индексе пути; мёртв по прибытии. Предиктор полностью
детерминирован: калибровочный прогон без затухания записывает
кинематику путника (индекс пути на каждый тик, с учётом поворотов), а
зарегистрированное предсказание проигрывает гейт в замкнутой форме по
этой траектории --- тот же приём, который дал $6/6$ на отложенной
выборке сопутствующей статьи.

По 5 лабиринтам $\times$ 11 задержкам $\times$ 2 уровням доверия (110
ячеек, зарегистрированы до исполнения): \textbf{110/110 совпадений
режимов; все 21 разрыв в середине маршрута --- ровно на предсказанном
индексе ребра (ошибка $0$)}; и переброс доверия для маршрутов --- при
доверии $0.7$ масса ребра с однократным проходом даёт
$\mathrm{age}_{\max} = 9.6 < L$, так что \emph{ни один маршрут не
фиксируется вовсе} ($0$ наблюдённых фиксаций; дискретный выключатель
сопутствующего закона, теперь для путей). Разрыв перенесённого
маршрута --- не шум, а закон: те же константы
$(\alpha,\tau,\text{trust})$, которые задают, КОГДА закрывается гейт
варпа мест, задают и то, ГДЕ рвётся маршрутный варп.

\section{Стратификация по опасностям: песня защищает}
\label{p3:sec:hazard}

\begin{table}[h]
\centering
\caption{Открытые поля, густо усеянные проходимыми, но штрафуемыми
опасностями (плотность $0.20$); маршрут свидетеля избегает опасностей
по опыту; 20 парных миров, затухание выключено (ось безопасности
изолирована от оси старения). Оба режима переноса ЗАВЕРШАЮТ задачу ---
ось стен (табл.~\ref{p3:tab:cliff}) и ось опасностей независимы;
различается цена.}
\label{p3:tab:hazard}
\small
\begin{tabular}{lccc}
\toprule
Плечо & попадания в опасности [CI] & завершение & среднее $t$ \\
\midrule
маршрут (чужой безопасный путь) & \textbf{0} (все 20 ячеек) & 20/20 & 16.6 \\
место (только ГДЕ) & 3.25 $[2.45, 4.10]$ & 20/20 & 15.9 \\
слепой & 28.7 & --- & --- \\
\bottomrule
\end{tabular}
\end{table}

\begin{figure}[t]
\centering
\includegraphics[width=0.92\linewidth]{figures/fig_route_hazard}
\caption{Песня защищает. \textbf{(a)} Попадания в опасности за эпизод:
путник на маршруте свидетеля сохраняет измеренный профиль риска свидетеля в пределах валидированного префикса маршрута ($0$
во всех 20 мирах); путник с одними местами платит $3.25$ попадания за
то же завершение; слепое исследование платит $28.7$. \textbf{(b)}
Композиция с законом разрыва: в режиме разрыва путник не получает ни
одного попадания, пока длится песня, а после разрыва риск
возобновляется --- и разрыв происходит на предсказанном экспериментом
R2 индексе ребра во всех $19/19$ ячейках.}
\label{p3:fig:hazard}
\end{figure}

Путник, следующий по пути свидетеля, сохраняет измеренный профиль
риска свидетеля в пределах валидированного префикса маршрута --- за надбавку времени в $+0.7$ тика
(рис.~\ref{p3:fig:hazard}a). Композиция с \S\ref{p3:sec:rupture}: в режиме
разрыва (19/20 миров) путник получает \textbf{0 попаданий до разрыва}
и $0.58$ $[0.11,1.32]$ после него, когда откатывается на знание мест и
заново входит в поле, --- причём все 19 разрывов снова на
предсказанном R2 индексе (ошибка $0$; предиктор с точностью до ребра
воспроизведён на новом классе миров без какой-либо подстройки). Песня
защищает, пока длится, и перестаёт защищать ровно там, где по закону
кончается куплет.

\section{Смысловая идентичность мест}
\label{p3:sec:identity}

Всё сказанное до сих пор --- и всё в сопутствующих статьях ---
держится на общей системе координат. Теперь мы убираем её: приватная
система отсчёта каждого агента --- это истинная сетка под секретным
сдвигом (позже: плюс секретный поворот, кратный $90^\circ$), а
широковещания несут приватные координаты, бессмысленные для
получателя, пока он не решит, что ОЗНАЧАЮТ места отправителя.

\paragraph{Два уровня <<того же самого>> и что такое ориентир.}
У межагентной идентичности мест два уровня, которые нельзя смешивать.
\emph{Семантический предикат} говорит, что две клетки принадлежат
одному желаемому классу (\texttt{water\_source}, GOOD): он едет в
каждом снимке, и именно над ним определено удовлетворение задаче. Но
никакой предикат не может решить \emph{соответствие сущностей} --- что
две записи памяти обозначают один и тот же референт в мире, --- потому
что предикат и есть то, что алиасирует: любая вода выглядит как вода.
Маршрутно-относительное описание (<<$n$ шагов после последнего
GOOD>>) его тоже не решает: это провенанс пути, а не идентичность его
концов (\S\ref{p3:sec:route}--\S\ref{p3:sec:rupture} измеряют, что оно даёт
и где рвётся). Соответствие в этой серии решает \emph{ориентир}
(landmark): место, чья локальная семантическая констелляция
поддерживает взаимно-уникальное сопоставление между агентами. Ориентир
--- это \emph{роль}, а не тип клетки: какие классы тегов способны её
играть, ниже само становится предметом измерения
(\S\ref{p3:sec:ablation}).

\paragraph{Отпечатки и рукопожатие.} Отпечаток места --- это локальная
семантическая констелляция вокруг посещённой клетки: какие
содержательные теги (опасности, стены, вода) стоят на каких
относительных смещениях; он трансляционно-инвариантен и свободен от
координат. Сопоставление \emph{взаимно уникально} (пара принимается,
только если каждая сторона --- единственный кандидат другой стороны
выше порога точного равенства; отправитель знает куда больше мест, чем
получатель, и односторонняя уникальность алиасирует далёкие
<<двойники>> в окрестность получателя), отпечаткам нужно $\ge 2$
значимых ключей, а сдвиг системы отсчёта --- модальная дельта $\ge 3$
согласующихся пар. Ориентиры обязаны нести хотя бы один
СОДЕРЖАТЕЛЬНЫЙ тег: узоры из одной лишь границы (<<пустота>>)
описывают форму мира, а не место, и исключаются --- урок угловой
ловушки ниже.

\paragraph{Результаты свидетель--путник.} По 4 смещениям $\times$ 2
плотностям опасностей $\times$ 10 сидам: семантический выравниватель
восстанавливает секретное смещение \emph{точно в 100\%} прогонов в
мирах с ориентирами (в среднем $15.1$ взаимной пары) и всякий раз
завершает задачу со строгим $\wstar$ ($\varphi{=}1$); координатный
бейзлайн при любом ненулевом смещении в $100\%$ прогонов захватывает
клетку без воды (его редкие <<успехи>>, $6.7\%$, --- это смещённый
маршрут, случайно пересекающий настоящую воду). В алиасированном мире
(безликая внутренняя полоса) два режима идентичности ломаются в
ПРОТИВОПОЛОЖНЫХ направлениях: семантическая идентичность не делает
\emph{ни одного} захвата --- отказывает закрытым образом, --- тогда
как координатная продолжает уверенно захватывать неверные клетки ---
отказывает открытым. Именно эта асимметрия, а не сырая точность, ---
самый глубокий аргумент в пользу смысловой идентичности: когда мир не
поддерживает опознание, семантическая система это знает.

\subsection{Полный стек и снятие последней оси}
\label{p3:sec:stack}

\paragraph{Онлайн-выравнивание в peer-стеке.} При $N{=}4$ агентах в
условиях дефицита, каждый в своей секретной системе отсчёта,
выравнивание складывается ПО ХОДУ эпизода: агенты исследуют, совместно
посещённые ориентиры накапливаются, и системы отсчёта по каждому
отправителю защёлкиваются посреди прогона (наступление в $100\%$
эпизодов, средний тик ${\approx}15$). Объединённый успех: оракул с
общей системой отсчёта $0.456$, координаты при рассогласовании
$0.350$, семантика $0.444$ --- \textbf{возвращено 88\% зазора}, со
строгим $\wstar$, работающим через восстановленные системы отсчёта.

\paragraph{Закон без системы отсчёта.} Перезапуск сопутствующего
закона расстояния с гейтом доверие$\times$давность ПОВЕРХ
семантической идентичности при смещении системы отсчёта даёт точки
излома $20/14/9$ --- \emph{в точности те же целые числа, что и
замкнутая форма при общей системе отсчёта}. Гейт и восстановление
системы отсчёта --- ортогональные слои: закон не замечает, что общая
сетка исчезла.

\paragraph{$SE(2)$: без общего севера.} С добавлением секретных
поворотов одна гипотеза поворота на каждые $90^\circ$ переиспользует
трансляционный выравниватель, а победитель обязан \emph{доминировать}
($\ge 2\times$ консенсуса ближайшего конкурента). Результаты: точное
восстановление $(r^\ast,\delta^\ast)$ в \textbf{80/80} комбинациях
поворот$\times$смещение$\times$сид со $100\%$ завершением со строгим
$\wstar$; в $180^\circ$-симметричном мире --- где система отсчёта
невосстановима \emph{в принципе} --- выравниватель отказывается
($0/10$ захватов), тогда как координаты отказывают открытым образом
($10/10$ неверных захватов); и закон расстояния снова ничего не
замечает ($\mathrm{bp}=20$ при свидетеле, повернутом на $90^\circ$).
Чисто трансляционный выравниватель, для контраста, никогда не
восстанавливает повернутую систему отсчёта ($0/60$), но не всегда
отказывает закрытым образом: в $2/60$ ячеек $180^\circ$-симметричная
подструктура опасностей позволяет ему отказать ОТКРЫТЫМ ---
эмпирическая мотивация правила доминирования.

\begin{figure}[t]
\centering
\includegraphics[width=0.92\linewidth]{figures/fig_route_identity}
\caption{Смысловая идентичность при $SE(2)$-системах отсчёта.
\textbf{(a)} Каждый строгий $\wstar$-захват семантического стека
указывает на настоящую воду ($8/8$); координатный стек делает $992$
захвата, и все --- фантомы ($0/992$); само восстановление системы
отсчёта точно в $80/80$ комбинаций свидетель--путник. \textbf{(b)}
Асимметрия отказов в мирах, где опознание невозможно (безликая полоса;
$180^\circ$-симметричная планировка, где система отсчёта
невосстановима в принципе): семантическая идентичность отказывается
захватывать --- отказывает закрытым образом, --- тогда как
координатная продолжает уверенно захватывать неверные клетки ---
отказывает открытым.}
\label{p3:fig:identity}
\end{figure}

\paragraph{Угловая ловушка.} В полном стеке сырой успех
недискриминативен при поворотах (повернутый координатный <<яд>>
большей частью отображается за пределы сетки и вырождается в
псевдоисследование), поэтому зарегистрированный дискриминатор ---
\emph{точность варп-захватов}: доля строгих $\wstar$-захватов,
указывающих на настоящую воду. До фильтра пустоты ранние граничные
отпечатки (агенты стартуют в углах; границы прямоугольника
$180^\circ$-симметричны) питали гипотезы неверной системы отсчёта, и
семантическая точность была $0.159$; с исключением границ как
не-ориентиров она равна \textbf{$1.000$ ($8/8$) против $0.000$
($0/992$ фантомных захватов) у координат}, с наступлением
выравнивания в $100\%$ эпизодов и чистыми регрессиями на всех более
ранних экспериментах (рис.~\ref{p3:fig:identity}).

\subsection{Что делает ориентир ориентиром: абляция словаря}
\label{p3:sec:ablation}

Весь механизм идентичности держится на фиксированном словаре
ориентиров $\Sigma = \{\texttt{стена}, \texttt{опасность},
\texttt{вода}, \texttt{пустота}\}$: какие классы тегов допускаются в
отпечаток. Все утверждения выше, стало быть, условны на выбранном
вручную алфавите. Последний эксперимент делает алфавит самим
воздействием: восстановление системы отсчёта свидетель--путник (миры
с рассеянными стенами И рассеянными опасностями; $3$ смещения
$\times$ $12$ сидов на плечо) при семи словарях --- от полного
$\Sigma$ до одиночных классов.

\begin{table}[h]
\centering
\caption{Абляция словаря ориентиров ($36$ ячеек на плечо, всего $252$;
зарегистрированные предсказания A1--A4, все PASS). Обеднение алфавита
стоит \emph{покрытия} --- восстановление превращается в отказ, --- но
никогда \emph{корректности}: каждый выполненный перенос, при любом
словаре, идёт через точно восстановленную систему отсчёта. Атрибуция
= доля массы ключей консенсусных пар при полном $\Sigma$.}
\label{p3:tab:ablation}
\small
\begin{tabular}{lcccc}
\toprule
Словарь & точное восст. & средн. пар & точность захвата & fail-open \\
\midrule
полный $\Sigma$ & $1.00$ & $33.2$ & $1.00$ & $0$ \\
$-$ пустота & $1.00$ & $32.0$ & $1.00$ & $0$ \\
$-$ опасности & $0.92$ & $11.8$ & $1.00$ & $0$ \\
$-$ стены & $1.00$ & $15.5$ & $1.00$ & $0$ \\
только опасности & $1.00$ & $13.7$ & $1.00$ & $0$ \\
только стены & $0.83$ & $10.8$ & $1.00$ & $0$ \\
только пустота & $0.00$ & $0.0$ & --- & $0$ \\
\bottomrule
\end{tabular}
\end{table}

Четыре зарегистрированных предсказания --- четыре попадания
(табл.~\ref{p3:tab:ablation}). \textbf{(A1) Абляция безопасна
(fail-safe):} в $252$ ячейках нет ни одного fail-open переноса ---
асимметрия закрытых отказов из \S\ref{p3:sec:identity} выдерживает не
только неоднозначные \emph{миры}, но и обеднённые \emph{алфавиты}.
\textbf{(A2) Пустота --- леса, а не идентичность:} граничные узоры
несут $0.000$ массы консенсусных ключей, и их удаление ничего не
меняет --- согласуется с угловой ловушкой, где они активно отравляли
гипотезы поворота. \textbf{(A3) Одного содержательного класса
достаточно:} одни опасности восстанавливают каждую систему отсчёта;
одни стены --- $0.83$, отказываясь в остальных. Соответствию нужно
лишь $\ge 3$ взаимно-уникальных констелляций над \emph{каким-нибудь}
содержательным классом --- бюджет ограниченной рациональности, а не
богатая онтология. \textbf{(A4) Атрибуция измерима:} констелляции
опасностей несут $0.56$ консенсусной массы, стены --- $0.44$. Какие
признаки работают ориентирами --- это, таким образом, не удачно
угаданное модельное допущение, а измеримое свойство пары
мир$\times$словарь, и при неверном словаре механизм деградирует вдоль
безопасной оси.

\subsection{Анатомия песни: узлы несут идентичность, рёбра --- перенос}
\label{p3:sec:anatomy}

\S\ref{p3:sec:identity} открывался утверждением, которое до сих пор мы
лишь аргументировали: маршрутно-относительное описание (<<$n$ шагов
после последнего GOOD>>) --- провенанс пути, а не идентичность его
концов. Последний эксперимент это измеряет. Свидетель сжимает свой
маршрут к воде в чисто \emph{реляционную} песню --- последовательность
куплетов $(\sigma, e)$, где $\sigma$ --- констелляционная сигнатура
ориентира, а $e$ --- точный такт-вектор от предыдущего куплета, ---
не несущую \emph{никаких глобальных координат}. Путник, ограниченный
внутренней полосой и никогда не видевший воду, её получает. Четыре
зарегистрированных плеча (кодек бит зафиксирован до прогонов; $12$
парных миров):

\begin{itemize}[leftmargin=16pt,itemsep=2pt,topsep=2pt]
\item \textbf{Одни такты} (чистое <<$n$ шагов от последнего GOOD>>):
  применимы, только \emph{предположив} якорь. При общем якоре ---
  точны ($11/11$, отрицательный контроль); при рассогласовании якорей
  захватывают всякий раз и не правы \emph{никогда} ($11/11$ fail-open
  промахов, смещённых ровно на разность якорей). Такты не несут
  идентичность --- они паразитируют на идентичности якоря.
\item \textbf{Одни сигнатуры}: опознают совместно посещённые места,
  но им не по чему идти --- \emph{ноль} переносов невиданной цели.
  Одни узлы не достигают невиданного.
\item \textbf{Сигнатуры $+$ такты}: куплеты сопоставлены
  взаимно-уникально, цепочка тактов проверена на метрическую
  согласованность между сопоставленными куплетами (loop-closure
  тест), цель дорассчитана от последнего сопоставленного ориентира.
  Покрытие $1.0$, каждый перенос точен, ноль fail-open --- за $366$
  бит на песню.
\item \textbf{Полные снапшоты} (машинерия \S\ref{p3:sec:identity}):
  информационная верхняя граница --- та же корректность за $5{,}690$
  бит.
\end{itemize}

Все четыре зарегистрированных предсказания сбылись. Песня --- не
облегчённый снапшот, а другой \emph{род} объекта: идентичность живёт
в узлах (ориентирах), перенос --- в рёбрах (тактах), и их композиция
даёт перенос уровня полного снапшота за $6.4\%$ битовой стоимости,
наследуя fail-safe свойство \S\ref{p3:sec:ablation}: всё, что ломается,
ломается отказом.

\section{Связанные работы}
\label{p3:sec:related}

\paragraph{Слияние карт и распознавание мест.} Мультироботный SLAM
трактует межкартовую ассоциацию данных как задачу оценивания над
метрическими признаками~\cite{cadena2016slam}; визуальное
распознавание мест~\cite{lowry2016vpr} сопоставляет внешний вид. Наш
вклад --- не конкурирующий сопоставитель, а измерительный слой поверх:
завершение, обусловленное провенансом, через восстановленные системы
отсчёта; асимметрия <<отказ закрытым/открытым образом>> как
зарегистрированное предсказание; инвариантность закона при удалении
системы отсчёта.
\paragraph{Обмен маршрутами.} Стигмергия переносит пути анонимно
(градиенты феромонов); обучаемая коммуникация (например,
CommNet~\cite{sukhbaatar2016commnet}) переносит политики непрозрачно.
Ни та, ни другая не выставляет разложение по провенансу рёбер, поэтому
ни та, ни другая не может сказать --- тем более предсказать, --- где
рвётся перенесённый маршрут. Предиктор разрыва --- это идея Time
Warp~\cite{jefferson1985}, применённая к затуханию памяти:
детерминированное проигрывание гейта по калиброванной кинематике.
\paragraph{Сонглайны.} Нарративная рамка теперь несущая в обоих
направлениях: песня --- это последовательность (маршрутный варп), её
места опознаются по смыслу (семантическая идентичность), она угасает
куплет за куплетом (закон разрыва) и защищает идущего, пока длится
(стратификация по опасностям).
% TODO: add refs: ProcTHOR/ALFWorld for future substrates.

\section{Ограничения и честные отрицательные результаты}
\label{p3:sec:limits}

(1) Сеточные миры; только сдвиги и повороты на $90^\circ$ ---
непрерывное $SE(2)$ и шум внешнего вида остаются будущей работой
(порог отпечатков --- точное равенство, поскольку субстрат бесшумен;
шумным субстратам нужна машинерия зазоров из области распознавания
мест). (2) Четыре зарегистрированных предсказания провалились и
сообщаются вместе с механизмами: отрицательный контроль $D{=}1$ в
тиках (экономия поворотов; точен по длине пути), монотонность выигрыша
по $D$ (цензурирование; истина --- обрыв), сравнение успеха полного
стека при поворотах (недискриминативно; заменено точностью захватов) и
<<трансляция отказывает закрытым образом>> ($2/60$ отказывают открытым
через симметричные подструктуры). (3) Следование по маршруту --- новый
ярус планировщика; все три плеча делят его, но ценность маршрутного
варпа измерена относительно именно этого потребителя. (4) Безопасность
по опасностям переносится по построению, когда маршрут свидетеля
безопасен; измеренное содержание --- в том, что перенос мест НЕ МОЖЕТ
её нести и что защита кончается на предсказанном разрыве. (5)
Завершение варп-захватов на уровне дефицита по-прежнему управляется
феноменологией столкновений из сопутствующей статьи. (6) Словарь
ориентиров $\Sigma$ зафиксирован вручную; \S\ref{p3:sec:ablation}
измеряет его абляцию, но никакой механизм его не \emph{отбирает}. Что
даёт признаку статус ориентира здесь --- наш выбор алфавита плюс
готовность мира сотрудничать; механизм, в котором признаки
зарабатывают этот статус участием в успешных выравниваниях --- словарь
формируется координацией, а не проектируется, --- естественный
следующий шаг серии, и здесь он не предпринимается.

\section{Заключение}

Сведите песню к точке на общей карте --- и незаметно ломаются две
вещи: КАК (фатально на первой же стене, дорого в любом поле
опасностей) и КТО-ЧТО-ИМЕЕТ-В-ВИДУ (фатально при любом рассогласовании
систем отсчёта --- и опасно именно потому, что отказывает открытым
образом). Верните песню --- рёбра с провенансом, места со смыслом ---
и оба отказа становятся измеримыми, закономерными и в значительной
мере исправленными: маршруты переносят завершение и безопасность и
рвутся только там, где велит закон старения; смысл восстанавливает
системы отсчёта точно, честно отказывается, когда мир неоднозначен, и
переносит весь измерительный стек --- $\varphi$, гейты, закон
расстояния --- без изменений в мир, где агенты не делят никаких
координат. Даже алфавит, которым записана песня, оказался измеримым,
а не священным: идентичность переживает обеднение словаря безопасно,
теряя покрытие, но не истину. Точка никогда не была памятью. Память
--- это песня; и открытый вопрос, который заостряет абляция, --- не
какие признаки ЯВЛЯЮТСЯ ориентирами, а какие признаки координация
самого коллектива СДЕЛАЕТ ориентирами.

\FloatBarrier
\begin{center}\rule{0.55\linewidth}{0.4pt}\\[2pt]\textbf{Приложения к этой части}\end{center}


\section{Зарегистрированные предсказания и ревизии}
\label{p3:app:registered}
Все предсказания были записаны на диск до исполнения; артефакт
сохраняет обе версии всюду, где операционализация пересматривалась.
Ревизии: отрицательный контроль R1 перемерен по длине пути;
монотонность выигрыша R1 переформулирована как детекция обрыва;
дискриминатор полного стека W9 сменён с успеха на точность
варп-захватов; функционирование строгого $\wstar$ в W8 расширено на
варп-ассистированные собственные завершения (такси-цепочки). Абляция
словаря (A1--A4, \S\ref{p3:sec:ablation}) и эксперимент по анатомии песни
(S0.1--S0.4, \S\ref{p3:sec:anatomy}, кодек бит зарегистрирован вместе с
предсказаниями) прошли как зарегистрировано, без ревизий.
% TODO: inline the registration JSONs.

\section{Воспроизводимость}
\label{p3:app:repro}
\begin{verbatim}
PYTHONPATH=. python experiments/warp/exp_route_warp_r0.py
PYTHONPATH=. python experiments/warp/exp_route_warp_r1.py --seeds 20
PYTHONPATH=. python experiments/warp/exp_route_warp_r2.py --seeds 5
PYTHONPATH=. python experiments/warp/exp_route_warp_r3.py --seeds 20
PYTHONPATH=. python experiments/warp/exp_warp_semantic_identity.py
PYTHONPATH=. python experiments/warp/exp_warp_semantic_stack.py --seeds 20
PYTHONPATH=. python experiments/warp/exp_warp_rotation.py --seeds 10
PYTHONPATH=. python experiments/warp/exp_warp_landmark_ablation.py
PYTHONPATH=. python experiments/song_grammar/exp_s0_song_smoke.py
\end{verbatim}
\FloatBarrier

\clearpage
\part{Коллективная семантическая память: слияние, доверие, устаревание}
\noindent\textit{Механика peer-памяти — три явных правила, третий режим на плоскости M×C — и категорная рамка, в которой общность интерпретации — параметр, а не предположение.}

\paragraph{Аннотация части.} \textbf{Проблема.} Коллективную память часто считают монотонной: чем больше общей информации, тем лучше. В кооперативной навигации с дефицитными семантическими целями эта интуиция не работает. Быстрое глобальное распространение информации улучшает материализацию целей, но заставляет агентов состязаться за одни и те же цели; медленный обмен избегает конкуренции, но оставляет агентов с устаревшей или неполной памятью.
\textbf{Подход.} Мы реализуем четыре топологии памяти над одним и тем же символьно-семантическим субстратом: независимые приватные памяти, общую шину событий, центральный консенсусный агрегатор и пиринговую широковещательную рассылку с каденцией (cadence) $K$. Затем мы вводим минимальную \emph{коллективную семантическую память} (CSM): пиринговую рассылку плюс явные правила \emph{слияния} (merge), \emph{доверия} (trust) и \emph{устаревания} (staleness). Показатели стадий Q/R/M/C используются только как измерительная инструментовка, а не как вклад данной статьи.
\textbf{Результаты.} На многоагентной развёртке из 35\,640 запусков семейство пиринговых протоколов с фиксированной каденцией демонстрирует сдвиг узкого места: быстрая рассылка насыщает материализацию, но обрушивает завершение, тогда как медленная рассылка имеет противоположный профиль. Среднее время до успеха имеет внутренний оптимум при $K{=}8$. На бенчмарке дефицита из 3\,240 запусков минимальная CSM строго доминирует над пятью пиринговыми базовыми вариантами с фиксированной каденцией по доле успеха: CSM достигает $0.617$ при 95\%-м ДИ $[0.610,0.624]$, выше верхней границы ДИ каждого peer-$K$.
\textbf{Границы применимости.} Эксперименты проведены на грид-мирах. Утверждение состоит не в том, что это полная теория коллективной памяти; это воспроизводимое архитектурное разбиение, показывающее, что слияние/доверие/устаревание --- правильные явные поверхности управления после того, как одна лишь каденция исчерпана.

\FloatBarrier
\section{Введение}

Многие многоагентные системы обмениваются информацией либо записывая
в единую глобальную память, либо посылая обучаемые сообщения, либо
периодически обмениваясь локальным состоянием. Эти механизмы обычно
оцениваются сквозным образом. Если успех растёт, коммуникация
признаётся полезной; если нет --- причина остаётся непрозрачной. В этой
статье изучается более узкий архитектурный вопрос: когда несколько
агентов поддерживают семантическую память о местах, чем именно следует
делиться, с какой частотой и с каким весом?

Ответ не сводится к <<делиться всем и как можно скорее>>. В условиях
дефицита, когда $N$ агентов конкурируют за $T<N$ семантических целей,
свежая глобальная память может привести к тому, что все агенты выберут
одну и ту же цель. Память при этом точна, но бесполезна. И наоборот,
отсутствие коммуникации предотвращает стадное поведение, но лишает
агентов полезных открытий, сделанных соседями. Задача проектирования
--- это, следовательно, компромисс между семантической свежестью и
конкуренцией за занятость целей.

Мы делаем три вклада.
\begin{enumerate}[leftmargin=20pt,itemsep=2pt,topsep=2pt]
\item \textbf{Архитектурное разбиение.} Мы реализуем и сравниваем четыре топологии памяти над одним и тем же символьным субстратом: независимую, общую шину, центральный агрегатор и пиринговую широковещательную рассылку.
\item \textbf{Результат о каденции.} Мы показываем, что пиринговая рассылка обнажает реальный координационный параметр $K$: быстрая рассылка улучшает $P(M^\star|R^\star)$, но вредит $P(C^\star|M^\star)$; метрика эффективности имеет внутренний оптимум.
\item \textbf{Минимальная CSM.} Мы реализуем коллективную семантическую память с тремя явными правилами: слияние, доверие и устаревание. Минимальная CSM достигает новой точки на плоскости M$\times$C и строго доминирует над пиринговыми базовыми вариантами с фиксированной каденцией по доле успеха на бенчмарке дефицита.
\end{enumerate}

\paragraph{Связь с сопутствующей статьёй о Q/R/M/C.}
Сопутствующая статья рассматривает Q/R/M/C как диагностический
фреймворк. Здесь эти показатели стадий используются как измерительные
инструменты. Основной объект иной: само устройство распределённой
символьной памяти.

\section{Субстрат памяти и задача}
\label{p4:sec:substrate}

Каждый агент наблюдает грид-среду и порождает семантические
свидетельства о местах, такие как \texttt{water\_source},
\texttt{hazard\_edge} и \texttt{open}. Локальная память хранит записи
о местах, уверенности в тегах, счётчики посещений, свежесть и
статистику переходов. Запросы планировщика извлекают
места-кандидаты путём сопоставления обязательных, предпочтительных и
штрафных тегов. Выбранный кандидат материализуется в конкретную клетку
и исполняется локальным контроллером.

Многоагентная задача использует сетку $12{\times}10$. Каждый агент
должен достичь клетки с тегом \texttt{water\_source}. Цели дефицитны:
$N{\in}\{3,5,8\}$ агентов, $T{\le}N$ целей, три планировки, три
плотности опасностей и 40 сидов для основной развёртки по каденции.
Каждый запуск логирует события Q/R/M/C по каждому агенту:
\[
Q^\star \rightarrow R^\star \rightarrow M^\star \rightarrow C^\star,
\]
где $M^\star$ означает, что семантический кандидат был зафиксирован
как исполнимая цель, а $C^\star$ --- что цель была достигнута. Эти логи
не являются архитектурным вкладом; они позволяют увидеть, где именно
отказывает каждая топология.

\section{Четыре топологии коллективной памяти}
\label{p4:sec:archs}

\begin{figure}[t]
  \centering
  \includegraphics[width=\linewidth]{figures/fig_architectures.pdf}
  \caption{Четыре топологии памяти. \textbf{(a)} Независимая: приватная память у каждого агента, обмена нет. \textbf{(b)} Общая шина: одна глобальная память, каждый агент читает и пишет на каждом такте. \textbf{(c)} Центральный агрегатор: приватные памяти плюс хаб, сливающий снимки в консенсусный отчёт. \textbf{(d)} Пиринговая рассылка: приватные памяти с обменом снимками каждые $K$ тактов и слитыми представлениями у каждого агента.}
  \label{p4:fig:architectures}
\end{figure}

Четыре топологии используют одну и ту же локальную память и
планировщик. Они различаются только тем, где происходит слияние.

\begin{description}[leftmargin=10pt,topsep=2pt,itemsep=2pt]
\item[Независимая.] Каждый агент владеет приватным хранилищем событий и графом концептов. Канала коммуникации нет.
\item[Общая шина.] Все наблюдения дописываются в единую глобальную шину событий. Один общий граф концептов перестраивается из объединения всех событий.
\item[Центральный агрегатор.] Каждый агент владеет приватным графом концептов; центральный консенсусный слой забирает снимки и вычисляет один взвешенный по доверию отчёт для всех.
\item[Пиринговая рассылка.] Каждый агент владеет приватным графом концептов и приватным слитым представлением. Каждые $K$ тактов агенты рассылают снимки; каждый агент сливает свою память с последним снимком каждого соседа.
\end{description}

Ключевое различие --- не в примитиве кластеризации. Оно в том, на
каком уровне этот примитив инстанцируется:
\[
\begin{array}{rcl}
\text{независимая} &:& N \text{ графов, }0\text{ слияний}\\
\text{общая шина} &:& 1 \text{ граф, }1\text{ неявное глобальное слияние}\\
\text{центральная} &:& N \text{ графов }+1\text{ центральное слияние}\\
\text{пиринговая} &:& N \text{ графов }+N\text{ приватных слияний.}
\end{array}
\]

\section{Каденция обнажает механизм отказа}
\label{p4:sec:cadence}

У пиринговой топологии есть один явный параметр: каденция рассылки $K$.
Основная развёртка перебирает четыре архитектуры и восемь пиринговых
каденций $K{\in}\{1,2,4,8,16,32,48,64\}$, всего 35\,640 запусков.
Рисунок~\ref{p4:fig:conditional-rates} суммирует результат.

\begin{figure}[t]
  \centering
  \includegraphics[width=\linewidth]{figures/fig_conditional_rates.pdf}
  \caption{Условные показатели и среднее время до успеха по развёртке из 35\,640 запусков. Быстрый обмен улучшает материализацию целей, но вредит завершению из-за конкуренции; медленный обмен имеет противоположный профиль. Среднее время до успеха имеет внутренний оптимум при $K{=}8$.}
  \label{p4:fig:conditional-rates}
\end{figure}

Быстрая пиринговая рассылка и централизованный обмен живут в режиме
насыщения M и коллапса C: агенты надёжно фиксируют цели, но многие
фиксации указывают на цели, которые преследуют и соседи. Независимая
топология и медленная пиринговая рассылка живут в режиме дефицита M и
насыщения C: меньше агентов фиксируются на хороших целях, но те, кто
фиксируется, реже сталкиваются с занятостью со стороны соседей.

\begin{figure}[t]
  \centering
  \includegraphics[width=0.85\linewidth]{figures/fig_bottleneck_shift.pdf}
  \caption{Сдвиг узкого места по пиринговым каденциям. $P(M^\star|R^\star)$ падает, а $P(C^\star|M^\star)$ растёт с ростом $K$. Практическая метрика эффективности, среднее время до успеха, минимизируется при $K{=}8$.}
  \label{p4:fig:bottleneck-shift}
\end{figure}

По всем пиринговым запускам корреляция Спирмена $P(M^\star|R^\star)$ с
$K$ равна $-0.608$, а $P(C^\star|M^\star)$ с $K$ равна $+0.420$ (обе с
$p<10^{-4}$). Более сильная сигнатура видна при фиксированной каденции:
внутрикаденционная корреляция Пирсона между звеньями M и C достигает
пика $-0.61$ при $K{=}4$ и затухает к шуму с ростом $K$. Каденция,
следовательно, не просто настраивает <<объём коммуникации>>. Она
перемещает узкое место между семантическим принятием обязательства и
завершением на разделяемом ресурсе.

\section{Минимальная коллективная семантическая память}
\label{p4:sec:csm}

Развёртка по каденции показывает, что обмен снимками с фиксированной
частотой обнажает лишь одну поверхность управления. Коллективная
семантическая память должна предоставлять как минимум три:
\begin{description}[leftmargin=14pt,itemsep=1pt,topsep=2pt]
\item[Слияние.] Как комбинируются свидетельства соседей об одном и том же семантическом месте.
\item[Доверие.] Как взвешиваются свидетельства разных соседей, когда они конфликтуют или становятся ненадёжными.
\item[Устаревание.] Как старые свидетельства затухают, прежде чем повлиять на извлечение.
\end{description}

Минимальная CSM использует пиринговую рассылку при $K{=}8$ и
накладывает поверх три правила:
\begin{description}[leftmargin=14pt,itemsep=1pt,topsep=2pt]
\item[\textit{Слияние.}] Оценка места равна собственному свидетельству с весом 1 плюс свидетельства соседей с весом $\mathrm{trust}_i \cdot e^{-\alpha \cdot \mathrm{age}}$, где $\alpha=0.05$. Порог включения $\tau=0.30$.
\item[\textit{Доверие.}] Экспоненциальное скользящее среднее (EMA) по каждому соседу, отражающее согласованность успеха извлечения. Соседи, чей последний снимок поддерживает локально зафиксированную цель, наращивают доверие; остальные затухают к $0.4$. Скорость обновления $\beta=0.10$.
\item[\textit{Устаревание.}] Каждый широковещательный снимок снабжён меткой времени; вес слияния затухает как $e^{-\alpha\cdot\mathrm{age}}$. Снимки старше $10K$ тактов вычищаются.
\end{description}

Реализация является полностью совместимой заменой пиринговой памяти с
фиксированной каденцией:
\method{experiments/collective\_semantic\_memory/csm\_memory.py}.

\subsection{Результаты}

Мы сравниваем CSM с пиринговой рассылкой с фиксированной каденцией при
$K \in \{1,4,8,16,64\}$ на ячейках дефицита
$(N,T){\in}\{(3,2),(5,3),(8,5)\}$, трёх планировках, трёх плотностях
опасностей и 20 сидах: всего 3\,240 запусков.

\begin{table}[ht]
\caption{Минимальная CSM против пиринговой рассылки с фиксированной каденцией на бенчмарке дефицита. Нижняя граница 95\%-го бутстреп-ДИ CSM превышает верхнюю границу ДИ каждого peer-$K$.}
\label{p4:tab:csm-results}
\centering
\small
\begin{tabular}{l c c c}
\toprule
Архитектура & $P(M^\star|R^\star)$ & $P(C^\star|M^\star)$ & Доля успеха (95\% ДИ) \\
\midrule
peer ($K{=}1$)  & 0.991 & 0.586 & 0.577 [0.572, 0.586] \\
peer ($K{=}4$)  & 0.955 & 0.623 & 0.581 [0.572, 0.590] \\
peer ($K{=}8$)  & 0.730 & 0.871 & 0.599 [0.591, 0.607] \\
peer ($K{=}16$) & 0.697 & 0.892 & 0.597 [0.589, 0.606] \\
peer ($K{=}64$) & 0.684 & 0.900 & 0.601 [0.594, 0.609] \\
\textbf{CSM (минимальная)} & \textbf{0.782} & \textbf{0.825} & \textbf{0.617 [0.610, 0.624]} \\
\bottomrule
\end{tabular}
\end{table}

CSM занимает третий режим. Быстрые рассылки обладают превосходной
материализацией, но плохим завершением; медленные --- наоборот. CSM
удерживает оба условных звена выше 0.78. Абсолютный прирост успеха
скромен, от $+1.7$ до $+4.0$ процентных пунктов, но стадийная
декомпозиция показывает, почему он содержателен: архитектура меняет
совместную рабочую точку M$\times$C, а не просто перемещается вдоль
кривой фиксированной каденции.

\paragraph{Почему CSM выигрывает: вентиль, а не обмен.} Провенансная
аннотация тех же запусков даёт механизм. CSM обменивается не
\emph{лучше}, чем пиринговые варианты с фиксированной каденцией, ---
она \emph{отказывается действовать по устаревшим или недоверенным
чужим свидетельствам}. Правило включения из раздела~\ref{p4:sec:csm}
влечёт замкнутую формулу горизонта, за которым свидетельство соседа
уже не может сформировать намерение,
$\mathrm{age}_{\max}(\mathrm{trust}_i) = \alpha^{-1}
\ln(\mathrm{trust}_i \cdot c/\tau)$, и этот горизонт эмпирически
точен: измеренные пороги формирования намерений совпадают с формулой
такт в такт по всем уровням доверия (23/23.05, 18/18.59, 12/12.84 при
$\mathrm{trust}=1.0/0.8/0.6$), включая контрольную проверку по
зарегистрированному предсказанию на шести никогда не запускавшихся
ячейках ($6/6$ точных целочисленных точек излома) и дискретное
переключение доверия: $\mathrm{trust}\le\tau/c$ полностью обнуляет
обязательства по чужим свидетельствам. У пиринговых вариантов с
фиксированной каденцией такого вентиля нет --- их чужие свидетельства
никогда не истекают, --- и именно поэтому их завершение обрушивается
при быстром обмене. Полный анализ, обусловленный провенансом,
приводится в сопутствующей рукописи; здесь он служит механистическим
прочтением таблицы~\ref{p4:tab:csm-results}.

\section{Тождество мест из смысла: устранение общей системы координат}
\label{p4:sec:semantic-identity}

Всё изложенное выше --- и категорное прочтение ниже --- предполагает
общий пространственный базис $X$: <<то же место>> означает <<те же
$(x,y)$ в системе координат, которую среда выдаёт каждому агенту>>.
Этот раздел устраняет это предположение и сообщает, что оно скрывало.
Каждый агент теперь держит \emph{приватную} систему координат
(истинную сетку под секретным сдвигом, а в последнем эксперименте ---
и под секретным поворотом, кратным $90^\circ$); широковещательные
снимки несут приватные координаты отправителя, бессмысленные для
получателя, пока тот не решит, что места отправителя \emph{значат}.

\paragraph{Механизм.} \emph{Отпечаток места} (place fingerprint) ---
это локальная семантическая констелляция вокруг посещённой клетки:
какие содержательные теги (опасности, стены, вода) находятся на каких
относительных смещениях --- инвариантно к сдвигу и без координат. Два
агента сопоставляют места по \emph{взаимно уникальному} соответствию
отпечатков (каждая сторона --- единственный кандидат другой;
однонаправленная уникальность вклеивает далёкие подобия в окрестность
получателя, а узоры, состоящие только из границ, исключаются как
описания формы мира, а не места). Межагентная система координат ---
это модальный сдвиг по $\ge 3$ согласующимся сопоставленным парам; при
поворотах по одной гипотезе на каждые $90^\circ$ переиспользуется тот
же выравниватель, и победитель обязан превосходить консенсус
ближайшего конкурента в $2\times$. Чужие свидетельства --- включая
места, которых получатель никогда не видел, --- затем переносятся
через восстановленную систему координат.

\paragraph{Результаты.} Зарегистрированы до выполнения; полный
протокол в сопутствующей рукописи.
\begin{enumerate}[leftmargin=20pt,itemsep=2pt,topsep=2pt]
\item \textbf{Точное восстановление.} Свидетель--путник: секретный
  сдвиг восстанавливается точно в $100\%$ запусков в мирах с
  ориентирами, а при полных $SE(2)$-системах координат --- в $80/80$
  комбинаций поворот$\times$сдвиг$\times$сид, со $100\%$-м завершением
  на фиксациях, принятых по чужим свидетельствам.
\item \textbf{Асимметрия отказов.} Там, где идентификация невозможна
  --- в безликом регионе или в $180^\circ$-симметричной планировке,
  где система координат невосстановима \emph{в принципе}, ---
  семантическое тождество не делает \emph{ни одного} обязательства
  (отказывает закрытым образом, fail closed), тогда как координатное
  тождество продолжает уверенно фиксировать неверные клетки ($100\%$
  неверных фиксаций; отказывает открытым образом, fail open).
\item \textbf{Полный стек.} При $N{=}4$ агентах в различных секретных
  системах координат выравнивание складывается онлайн (возникновение в
  $100\%$ эпизодов, средний такт ${\approx}15$) и восстанавливает
  $88\%$ разрыва производительности с общей системой координат (успех
  $0.444$ против $0.456$ у оракула против $0.350$ у координатного
  варианта); при случайных $SE(2)$-системах каждое семантическое
  обязательство указывает на настоящую воду ($8/8$), тогда как
  координатный стек производит $992$ фантомные фиксации ($0$
  настоящих).
\item \textbf{Вентиль ортогонален.} Повторный прогон анализа
  $\mathrm{age}_{\max}$ выше с правилом включения CSM ПОВЕРХ
  семантического тождества, при сдвиге системы координат и при
  повороте на $90^\circ$, воспроизводит ТЕ ЖЕ целочисленные точки
  излома ($20/14/9$; $\mathrm{bp}=20$ при повороте) --- вентиль
  доверие$\times$устаревание и восстановление системы координат суть
  независимые слои.
\end{enumerate}

Более сильное утверждение, обещанное в ограничениях ранних версий,
тем самым конструктивно исполнено: тождество мест (place identity)
может нестись одним лишь смыслом, оно компонуется со
слиянием/доверием/устареванием без изменений и --- в отличие от
координатной конвенции, которую оно заменяет, --- знает, когда оно не
знает.

\paragraph{Категорный индикатор: дефект кругового обхода сопряжения.}
Восстановление системы координат --- это в точности та операция,
которая нужна категорному прочтению из \S\ref{p4:sec:categorical}, как
только общий базис $X$ отброшен: при приватных системах координат
каждый агент несёт собственную категорию мест, а коммуникация --- это
пара функторов восстановления $F=\textsc{align}(A,B)$ (система
координат отправителя $B$, выраженная в координатах получателя $A$) и
$G=\textsc{align}(B,A)$. Если они действительно сопряжены --- системы
координат склеены, --- круговой обход $\eta_X: X\to G(F(X))$ есть
тождество, так что \emph{дефект сопряжения}
\[
\mathrm{def}=\lVert F.\mathrm{offset}+G.\mathrm{offset}\rVert_1
\ \text{(сдвиг)},\qquad
r_F{+}r_G\equiv 0\,(\mathrm{mod}\,4)\ \text{и}\
\lVert R_{r_G}F.\delta+G.\delta\rVert_1\ (SE(2)),
\]
равен нулю \citep{maclane1998categories}. Измеренный на описанном выше
методе восстановления (\texttt{experiments/warp/alignment\_defect.py},
3 сдвига $\times$ 10 сидов на ячейку): дефект \emph{в точности} равен
$0$ в $30/30$ запусков для сдвиговых систем координат и в $30/30$ для
каждого из четырёх поворотов на $90^\circ$ --- восстановленные функторы
являются строгими сопряжениями. Там, где восстановление обязано
отказать (агент, наблюдавший только безликую внутреннюю полосу), $F$
или $G$ не определён в $30/30$ запусков с \emph{нулём} ложных склеек
($\mathrm{def}>0$): сопряжения не существует, и выравниватель
отказывается его строить. Точный дискретный выравниватель, таким
образом, реализует дихотомию --- строгое сопряжение
($\mathrm{def}{=}0$) или отсутствие сопряжения --- и никогда не даёт
уверенной неверной склейки; <<отказ закрытым образом>> выше --- это в
точности \emph{отказ от несопряжённой пары}. Это и есть та величина,
которую категорная модель вычисляет, а простое переописание --- нет:
скаляр, равный нулю тогда и только тогда, когда локальные системы
координат коллектива когерентно склеены, и именно эта склейка --- а не
общая координатная конвенция --- была молчаливой предпосылкой слияния
из \S\ref{p4:sec:categorical}.

\section{Категорное BDI-прочтение коллективной семантической памяти}
\label{p4:sec:categorical}

Этот раздел даёт компактное теоретическое прочтение примера поиска
воды в терминах убеждения, желания и намерения (belief, desire,
intention). Цель --- не ввести новый планировщик, а сделать явной
математическую структуру, уже используемую архитектурой коллективной
семантической памяти.

Мы предполагаем конечное множество агентов $A=\{1,\ldots,N\}$, общий
пространственный базис $X$ и общий семантический алфавит $\Sigma$. В
текущей грид-реализации $X$ --- это общая система координат, а
$\Sigma$ содержит теги вроде \texttt{water\_source},
\texttt{water\_candidate}, \texttt{hazard\_edge} и \texttt{open}. Это
явное модельное предположение. Агенты могут иметь разные наблюдения,
разные памяти и разные значения уверенности, но они интерпретируют
места в общей координатной системе и выражают семантические
свидетельства в общем алфавите тегов. Без первого предположения
пиринговое слияние требует сначала решить задачу выравнивания мест ---
раздел~\ref{p4:sec:semantic-identity} решает её конструктивно (консенсус
отпечатков восстанавливает межагентную систему координат), а
раздел~\ref{p4:sec:colimit} формулирует, чем эта конструкция является
категорно. Выравнивание символов (общий $\Sigma$) остаётся
предположением.

У этого предположения есть, кроме того, измеримая \emph{внутренность}:
внутри фиксированного $\Sigma$ вопрос о том, какие классы тегов
фактически несут соответствие мест, --- эмпирическое свойство, а не
модельный выбор. Зарегистрированная абляция словаря в сопутствующей
статье о переносе маршрутов делает сам алфавит воздействием и
обнаруживает, что механизм идентичности \emph{безопасен в $\Sigma$}:
обеднение словаря вплоть до одного содержательного класса превращает
восстановление системы отсчёта в отказ и никогда --- в неверное
выравнивание (ноль fail-open переносов на $252$ ячейках), а
консенсусная масса разложима по классам (опасности $0.56$, стены
$0.44$, границы мира $0.00$). Чего пока не делает ни один механизм ---
это \emph{отбора} $\Sigma$: здесь признаки получают статус ориентира
от дизайна плюс готовности мира сотрудничать. Правило, по которому
признаки \emph{зарабатывают} этот статус участием в успешных
выравниваниях --- алфавит, формируемый координацией, а не
предполагаемый, --- самый острый открытый конец этого предположения и
естественная точка встречи с программой выравнивания символов
(\S\ref{p4:sec:symbol-alignment}).

В момент $t$ агент $i$ поддерживает локальный семантический контекст
$\mathcal{K}_t^i=(G_t^i,\Sigma,I_t^i)$, где $G_t^i$ --- множество
записей о местах в приватной памяти агента, а
\[
I_t^i \subseteq G_t^i \times \Sigma
\]
--- пороговое отношение инцидентности, утверждающее, что место $g\in
G_t^i$ несёт тег $\sigma\in\Sigma$ с достаточной уверенностью.
Градуированная реализация использует уверенности тегов и взвешенные
оценки; чёткое отношение $I_t^i$ --- это пороговый скелет, нужный для
теоретического описания.

Этот формальный контекст индуцирует соответствие Галуа между
множествами тегов и множествами мест:
\[
\operatorname{ext}_t^i(B)=\{g\in G_t^i:\forall\sigma\in B,\ (g,\sigma)\in I_t^i\},
\qquad
\operatorname{int}_t^i(S)=\{\sigma\in\Sigma:\forall g\in S,\ (g,\sigma)\in I_t^i\}.
\]
Для задачи поиска воды желаемый семантический предикат ---
$B_{\mathrm{water}}=\{\texttt{water\_source}\}$ или, в реализованной
взвешенной форме, предикат, комбинирующий положительные свидетельства
для \texttt{water\_source}, \texttt{water\_candidate} и
\texttt{water\_nearby} со штрафами за \texttt{hazard\_edge} или
близкий риск. Извлечение успешно для агента $i$ в момент $t$ ровно
тогда, когда
\[
\operatorname{ext}_t^i(B_{\mathrm{water}})\neq\varnothing .
\]
Таким образом, стадия извлечения --- это не поиск по координатам. Это
проверка того, что запрошенный семантический смысл уже имеет непустой
объём в реализованной памяти агента.

Теперь можно определить компоненты BDI.

\paragraph{Убеждение.} Состояние убеждений агента $i$ --- это
структурированная семантическая память, индуцированная
$\mathcal{K}_t^i$. Эквивалентно, это решётка концептов
$L_t^i=\operatorname{Concepts}(\mathcal{K}_t^i)$, узлы которой
сопрягают объёмы мест с содержаниями тегов. Убеждение о воде --- это
не пропозиция <<вода находится в координате $x$>>. Это область
$L_t^i$, чьё содержание включает водный предикат. Если эта область
пуста, агент не просто плохо выбирает; соответствующий семантический
объект отсутствует в его текущей памяти.

\paragraph{Желание.} Желание --- это обусловленный состоянием
семантический предикат над решёткой убеждений. В примере с водой
внутреннее состояние агента содержит переменную жажды. Активное
желание можно записать как
\[
D_t^i =
\begin{cases}
B_{\mathrm{water}}, & \text{если } \mathrm{thirst}_t^i\geq\theta_{\mathrm{on}},\\
B_{\mathrm{goal}}, & \text{иначе.}
\end{cases}
\]
Итак, желание --- не дополнительный метафизический объект. Это
правило, выбирающее, какой семантический предикат следует запрашивать
из текущей структуры убеждений.

\paragraph{Намерение.} Намерение --- это обязательство по одному
материализованному кандидату. Как только желание выбрало предикат, а
извлечение дало непустой объём, планировщик применяет
удовлетворяющую (satisficing) функцию выбора
\[
\chi_t^i:\mathcal{P}(G_t^i)\rightarrow G_t^i\cup\{\bot\},
\qquad
g_t^{i,\star}=\chi_t^i\!\big(\operatorname{ext}_t^i(D_t^i)\big),
\]
где $\bot$ означает неудачу материализации цели. В терминах Q/R/M/C
желание определяет запрос $Q^\star$, убеждение поддерживает извлечение
$R^\star$, а намерение соответствует материализации $M^\star$.
Завершение $C^\star$ находится ниже по потоку от этой семантической
структуры: оно зависит от того, достигает ли физический контроллер
зафиксированной цели и остаётся ли цель доступной при многоагентной
занятости.

Это даёт компактное категорное прочтение. Пусть $\mathbf{Ctx}_\Sigma$
--- категория, объекты которой суть семантические контексты
$(G,\Sigma,I)$ над фиксированным алфавитом $\Sigma$, а морфизмы
сохраняют семантическую инцидентность. Траектория локальной памяти
агента $i$ --- это индексированная временем диаграмма
\[
\mathcal{K}_0^i\rightarrow\mathcal{K}_1^i\rightarrow\cdots
\rightarrow\mathcal{K}_t^i\rightarrow\mathcal{K}_{t+1}^i\rightarrow\cdots
\]
в $\mathbf{Ctx}_\Sigma$. Каждая стрелка индуцирована наблюдением и
обновлением памяти. Когда новое свидетельство консолидируется,
формальный контекст меняется; следовательно, меняется его решётка
концептов. Таково порядково-теоретическое содержание роста памяти.

В терминах Q/R/M/C пустой объём
$\operatorname{ext}_t^i(D_t^i)=\varnothing$ --- это отказ
убеждения/извлечения, непустой объём без устойчивой фиксации --- отказ
материализации, а успешная фиксация, которую невозможно достичь, ---
отказ завершения: семантическая и физическая части состояния агента
отказывают по-разному и измеримы по отдельности.

Коллективный случай поднимает это на одну конструкцию выше. В моменты
рассылки получатель $i$ держит диаграмму локальных объектов памяти
$\mathcal{D}_t^i=(\mathcal{K}_t^i,\mathcal{K}_t^{j_1},\ldots,
\mathcal{K}_t^{j_m})$ вместе с картами выравнивания и запрашивает
слитый консенсусный объект
$\widetilde{\mathcal{K}}_t^i=\operatorname{Merge}_{\mathrm{trust},
\mathrm{stale}}(\mathcal{D}_t^i)$, где свидетельства соседей
взвешиваются в точности по правилам раздела~\ref{p4:sec:csm}
($\mathrm{trust}_{j\to i}\cdot e^{-\alpha\,\mathrm{age}_j}$). Те же
BDI-определения применимы к слитому контексту: намерение теперь
выбирается из $\operatorname{ext}_{\widetilde{\mathcal{K}}_t^i}(D_t^i)$,
а не из одного лишь приватного объёма.

Это объясняет, почему каденция коммуникации порождает сдвиг узкого
места. Более быстрая рассылка повышает свежесть свидетельств соседей.
Это, как правило, расширяет или улучшает водный объём
$\operatorname{ext}_{\widetilde{\mathcal{K}}_t^i}(B_{\mathrm{water}})$,
так что агенты материализуют цели надёжнее; в терминах Q/R/M/C это
повышает $P(M^\star\mid R^\star)$. Но то же свежее общее убеждение
синхронизирует и намерения. Несколько агентов могут зафиксироваться на
одном и том же дефицитном источнике воды. Память семантически точна,
но физическая цель теперь оспариваема. Это понижает $P(C^\star\mid
M^\star)$. Узкое место, следовательно, физически обосновано. Это не
формальный артефакт декомпозиции: оно возникает потому, что
семантическое убеждение разделяется через память, тогда как завершение
ограничено занятостью и дефицитными ресурсами. Коллективная
семантическая память должна поэтому регулировать не только объём
разделяемой информации, но и то, как она сливается, чьим
свидетельствам доверяют и когда устаревшие свидетельства должны
перестать влиять на формирование намерений.

В этом смысле <<коллектив>> --- не центральный разум и не глобальная
карта, а стабилизированное семейство индексированных соседями
процессов; следующий подраздел делает это утверждение точным.

Двойственное, процессное прочтение --- агенты как коалгебры (codata),
коллектив как финальная коалгебра, коммуникация как корекурсия глубины
один, а сдвиг узкого места как соревнование между бисимилярностью
убеждений и бисимилярностью намерений --- развито в
приложении~\ref{p4:sec:coalgebra}. Это линза, а не механизм: ничто ниже
по тексту от него не зависит, поэтому оно и живёт в приложении.

\subsection{Слияние как копредел и его инварианты}
\label{p4:sec:colimit}

Объект слияния выше можно описать точнее. Работая в версии
$\mathbf{Ctx}_\Sigma$, обогащённой (enrichment) метками времени и
провенансом, слияние снимков от агентов $1,\dots,N$ есть диаграмма
локальных объектов памяти плюс карты выравнивания, индуцированные
отношением тождества мест. Консенсусный объект
$\widetilde{\mathcal{K}}_t^i$ --- это копредел (colimit) этой
диаграммы, конкретно реализуемый слиянием union-find над выровненными
концептами. Такое прочтение даёт три полезных ограничения.
\begin{enumerate}[leftmargin=20pt,itemsep=2pt,topsep=2pt]
\item \textbf{Независимость от порядка.} Слитый объект единствен с точностью до изоморфизма, поэтому результат не является артефактом обработки снимков агентов в конкретном порядке.
\item \textbf{Цепочки ожидаемы.} Если концепт A выравнивается с B, а B --- с C, union-find выполняет транзитивное замыкание. Чрезмерное слияние, следовательно, не баг структуры данных; это цена выбранного отношения эквивалентности.
\item \textbf{Затухание --- не морфизм.} Устаревание хранится как поле веса вне морфизмов выравнивания. Иначе <<состарить-потом-слить>> и <<слить-потом-состарить>> не обязаны коммутировать, что разрушило бы чистую интерпретацию широковещательных снимков как одной диаграммы.
\end{enumerate}

В ранних версиях этот раздел заканчивался признанием, что модель не
отвечает на вопрос, какие перцептивные признаки должны отождествлять
два места --- карты выравнивания предполагались, а не строились. Этот
пробел теперь закрыт разделом~\ref{p4:sec:semantic-identity}, и это
закрытие даёт копредельному прочтению его недостающую научную функцию:
\begin{enumerate}[leftmargin=20pt,itemsep=2pt,topsep=2pt]
\item \textbf{Морфизмы выравнивания строятся, а не предполагаются.}
  Консенсус отпечатков И ЕСТЬ алгоритм, порождающий карты выравнивания
  диаграммы --- взаимно уникальные сопоставления ориентиров определяют
  преобразование системы координат, через которое вкладывается каждый
  локальный объект памяти.
\item \textbf{Существование становится эмпирическим предикатом.}
  Копредел существует тогда и только тогда, когда существует
  консенсус. В $180^\circ$-симметричном мире морфизм выравнивания не
  существует \emph{в принципе} (две несравнимые диаграммы
  поддерживают равный консенсус), и реализация обнаруживает именно
  это и отказывается сливать --- поведение <<отказа закрытым образом>>
  из раздела~\ref{p4:sec:semantic-identity} есть операционализированное
  категорное несуществование. Этот предикат измеряется, а не
  постулируется: дефект кругового обхода сопряжения из
  \S\ref{p4:sec:semantic-identity} в точности равен $0$, когда попарное
  выравнивание --- строгое сопряжение (во всех восстановленных
  случаях сдвига и поворота на $90^\circ$, $30/30$), и не определён,
  когда сопряжения не существует (отказ закрытым образом, $30/30$, с
  нулём ложных склеек) --- скалярный свидетель существования
  копредела.
\item \textbf{Затухание остаётся вне морфизмов} (инвариант 3), и
  эмпирически вентиль доверие$\times$устаревание не меняется при
  восстановлении системы координат --- те же целочисленные точки
  излома с общей системой координат и без неё, --- что подтверждает
  коммутирование двух слоёв, как того требует диаграммное прочтение.
\end{enumerate}

\subsection{Коллективная память как когерентная склейка локальных категорий мест}
\label{p4:sec:gluing}

Копредельное слияние из \S\ref{p4:sec:colimit} \emph{предполагало} карты
выравнивания, а \S\ref{p4:sec:semantic-identity} их \emph{строит}; теперь мы можем
сформулировать в качестве центрального категорного результата, каким объектом
является коллективная память, когда агенты держат \emph{разные} системы
координат, и в точности когда она существует точным (faithful) образом.

\paragraph{Постановка.} Реализованная память агента $i$ --- это тонкая
категория (частично упорядоченное множество): её решётка концептов
$L_i=\operatorname{Concepts}(\mathcal{K}_i)$ над приватной системой координат
$i$, объекты --- места, порядок --- включение объёмов из
\S\ref{p4:sec:categorical}. Восстановление системы координат
(\S\ref{p4:sec:semantic-identity}) даёт на сонаблюдаемом подпосете монотонное
отображение $F_{ij}:L_i\to L_j$, переносящее места $i$ в систему координат
$j$, и обратное к нему $F_{ji}$. Поскольку категории \emph{тонкие},
восстановленная пара $(F_{ij},F_{ji})$ является \emph{сопряжённой парой} ровно
тогда, когда она есть соответствие Галуа, с единицей
$\eta^{ij}:\operatorname{Id}_{L_i}\Rightarrow F_{ji}\!\circ\!F_{ij}$ ---
дефектом сопряжения, измеренным в \S\ref{p4:sec:semantic-identity}.

\paragraph{Определение (когерентная система выравнивания).} Семейство
$\{F_{ij}\}$ над графом коммуникации \emph{когерентно}, если (i) каждая пара
--- сопряжённое порядковое вложение ($\eta^{ij}$ --- изоморфизм, т.е.\ дефект
$0$: ни одно место не вклеивается в другое), и (ii) оно удовлетворяет условию
коцикла $F_{jk}\!\circ\!F_{ij}=F_{ik}$ на тройных перекрытиях.

\paragraph{Предложение (два агента).} Для сопряжённой пары $F\dashv G$ между
$L_A,L_B$ категория-коллаж (запятая) $\mathrm{Gl}(F)$ --- объекты
$L_A\sqcup L_B$, перекрёстные морфизмы $\mathrm{Gl}(F)(X,Y)=L_B(FX,Y)\cong L_A(X,GY)$
\citep{maclane1998categories} --- содержит $L_A,L_B$ как полные подпосеты и
является двухагентной коллективной памятью; дефект $\eta_X$ --- препятствие к
тому, чтобы $X$ вернулся в себя круговым обходом, и склейка отождествляет
различные места тогда и только тогда, когда $\eta$ не изоморфизм.

\paragraph{Теорема (коллективная память = копредел диаграммы выравнивания).}
Пусть $\mathbf{D}$ --- диаграмма в $\mathbf{Pos}$ (эквивалентно, в обогащённой
$\mathbf{Ctx}_\Sigma$), вершины которой --- $L_i$, а рёбра --- $F_{ij}$.
Коллективная память есть $\operatorname{colim}\mathbf{D}$. Так как
$\mathbf{Pos}$ кополна, этот копредел всегда существует; он является
\emph{точной} склейкой --- каждая $L_i$ вкладывается и никакие два различных
места не отождествляются --- \emph{тогда и только тогда}, когда система
выравнивания когерентна. Когда все $F_{ij}=\operatorname{id}$ (общая система
координат), $\mathbf{D}$ --- постоянная диаграмма и
$\operatorname{colim}\mathbf{D}$ --- в точности копредельное union-find-слияние
из \S\ref{p4:sec:colimit}; общее утверждение устраняет предположение об общей
системе координат, заменяя тождества восстановленными сопряжёнными вложениями.

\emph{Аргумент.} Кополнота $\mathbf{Pos}$ реализует копредел как дизъюнктное
объединение, факторизованное по предпорядку, порождённому
$F_{ij}$-отождествлениями. Каждая каноническая инъекция
$L_i\hookrightarrow\operatorname{colim}\mathbf{D}$ точна тогда и только тогда,
когда эти отождествления не добавляют схлопываний внутри $L_i$, что имеет
место тогда и только тогда, когда каждое $F_{ij}$ --- порядковое вложение
(единица --- изоморфизм, дефект $0$) и отождествления согласованы на
перекрытиях (коцикл) --- т.е.\ когерентность. Случай общей системы координат
--- постоянная диаграмма, чей копредел есть покомпонентное объединение, т.е.\
union-find. \hfill$\square$

\paragraph{Следствие (два вычислимых свидетеля).} Когерентность проверяема без
решения общей задачи семантического выравнивания. (a)~Попарный \emph{дефект
сопряжения} $\eta$ (\S\ref{p4:sec:semantic-identity}) равен нулю тогда и только
тогда, когда каждое ребро --- строгое сопряжённое вложение --- измерен в
точности $0$ в $30/30$ восстановленных случаев сдвига и поворота --- и
защищает от \emph{чрезмерного слияния}. (b)~\emph{Связность} графа
достижимости доверие$\times$каденция (\S\ref{p4:sec:phase-diagram},
рис.~\ref{p4:fig:phase-diagram}) решает, покрывает ли один
$\operatorname{colim}\mathbf{D}$ всех $N$ или диаграмма распадается на
копроизведение под-консенсусов (\emph{фрагментация}). Эти два механизма отказа
категорно различны и измеряются по отдельности: дефект $>0$ --- это
схлопывание, несвязность --- фрагментация: две оси
рис.~\ref{p4:fig:phase-diagram} и поведение отказа закрытым образом из
\S\ref{p4:sec:semantic-identity}.

\paragraph{Границы применимости.} Доказательства --- стандартная кополнота
$\mathbf{Pos}$ и конструкция коллажа \citep{maclane1998categories}; работа
теоремы в том, чтобы назвать объект и привязать его существование к двум
измеренным величинам. Эмпирическое содержание --- что восстановление по
отпечаткам действительно даёт сопряжённые вложения --- живёт в
\S\ref{p4:sec:semantic-identity}, а не в теореме. Одно предположение \emph{не}
снято: общий алфавит тегов $\Sigma$; при разных алфавитах $F_{ij}$ должны были
бы переводить ещё и $\Sigma$, что выходит за пределы текущего восстановления и
является честной следующей границей.

\subsection{Параметры протокола как категорная структура: фазовая диаграмма доверие\texorpdfstring{$\times$}{x}каденция}
\label{p4:sec:phase-diagram}

Прочтение до сих пор моделировало \emph{память} (объекты --- семантические
контексты, слияние --- их копредел) при общей системе координат и полном
обмене. Его предсказательная польза приходит, когда мы направляем линзу на
\emph{структуру коммуникации} и спрашиваем, \emph{какую форму принимает
копредел} как функцию двух параметров протокола --- чтобы категория вычисляла
структурный факт, а не переописывала алгоритм. (В $\mathbf{Pos}$ копредел
диаграммы \emph{существует} всегда --- по кополноте, связна диаграмма или
нет; структурный вопрос в том, является ли он единой консенсусной
компонентой или копроизведением под-консенсусов --- в точности дихотомия
\S\ref{p4:sec:gluing}.) Два параметра играют две
структурные роли:
\begin{description}[leftmargin=12pt,topsep=2pt,itemsep=1pt]
\item[Доверие $\tau$ управляет морфизмами.] Ребро $i\!\to\!j$ (агент $i$
включает свидетельства $j$) существует только там, где доверие $i$ к $j$ не
ниже $\tau$. Повышение $\tau$ прореживает граф доверия.
\item[Каденция $K$ управляет глубиной композиции.] На горизонте $H$ есть
$R=\lfloor H/K\rfloor$ раундов рассылки, поэтому составная цепочка
$i\!\to\!k\!\to\!\cdots\!\to\!j$ может иметь длину не более $R$ морфизмов.
Низкое $K$ (много раундов) допускает глубокую композицию; высокое $K$ ---
только короткие пути.
\end{description}
\emph{Охватывающий} (consensus-spanning) копредел --- единая связная
консенсусная компонента, сплавляющая и обслуживающая свидетельства каждого
агента --- получается тогда и только тогда, когда $R$-раундовая достижимость
$\tau$-прореженного графа доверия является все-ко-всем; иначе копредел
вырождается в копроизведение под-консенсусов.

Рисунок~\ref{p4:fig:phase-diagram} перебирает $\tau\times K$ для $N{=}8$ агентов,
чьё попарное доверие порождено собственным EMA-правилом системы
(\method{peer\_memory.peer\_trust.AsymmetricTrust}) при кольцевой топологии
сонаблюдения (соседи сонаблюдают часто и зарабатывают высокое доверие; далёкие
агенты --- нет), усреднённой по $12$ сидам
(\method{experiments/collective\_semantic\_memory/phase\_diagram\_trust\_cadence.py}).
Результат --- настоящий фазовый переход с диагональной границей. При низком
$\tau$ граф доверия плотен, и копредел охватывает всех агентов при любой
каденции. Повышение $\tau$ прореживает дальние морфизмы до почти-соседского
кольца; коллектив тогда сохраняет охватывающий копредел \emph{только если}
каденция достаточно низка, чтобы многошаговая композиция через всё ещё
доверенных посредников вновь его связала --- $\tau{=}0.55$ удерживает
охватывающий копредел вплоть до $R{=}2$,
$\tau{=}0.80$ требует $R\ge4$ (диаметр кольца), а $\tau{=}0.90$ фрагментирован
повсюду. Каденция и доверие, следовательно, не непрозрачные ручки: вместе они
определяют категорный структурный инвариант, и фазовая граница --- это в
точности место, где коллектив переходит от единого консенсусного объекта к
фрагментированным под-консенсусам.

\begin{figure}[t]
  \centering
  \includegraphics[width=\linewidth]{figures/fig_phase_diagram.pdf}
  \caption{Фазовая диаграмма структуры коммуникации по порогу доверия $\tau$
  (строки) и каденции $K$ (столбцы; $R=\lfloor H/K\rfloor$ раундов рассылки),
  $N{=}8$ агентов на кольце взаимодействия, $12$ сидов. \textbf{Слева:}
  плотность достижимости (доля упорядоченных пар агентов, чьи свидетельства
  могут скомпоноваться за $R$ раундов). \textbf{Справа:}
  $P(\text{копредел охватывает всех }N)$, т.е.\ достижимость
  все-ко-всем. Диагональная граница --- структурный переход: коллектив
  допускает единый консенсусный объект в области низких $\tau$/низких $K$ и
  фрагментируется в углу высоких $\tau$/высоких $K$, причём более высокое
  $\tau$ требует более низкого $K$ (больше раундов), чтобы остаться связным.}
  \label{p4:fig:phase-diagram}
\end{figure}

Именно в этом смысле категорная модель операциональна, а не описательна:
каждый параметр протокола отображается в конкретное изменение диаграммы
коммуникации, а структурный инвариант диаграммы (охватывает ли копредел всех
агентов или распадается в копроизведение) предсказывает качественный режим
коллектива. \emph{Границы применимости.}
Матрица доверия порождена настоящим EMA-правилом при смоделированной
кольцевой топологии сонаблюдения, а <<охватывающий консенсус>>
операционализирован как $R$-раундовая достижимость все-ко-всем; количественно
установить, что пересечение этой структурной границы предсказывает падение
наблюдаемой Q/R/M/C-производительности --- естественный следующий эксперимент.

\subsection{К выравниванию символов: заземление алфавитов тегов через общие места}
\label{p4:sec:symbol-alignment}

Теорема о склейке (\S\ref{p4:sec:gluing}) всё ещё предполагает общий алфавит
$\Sigma$. Отбрасывая его, каждый агент держит контекст $(G_i,\Sigma_i,I_i)$
над \emph{собственным} словарём, и правильный морфизм --- уже не функтор над
одним $\Sigma$, а \emph{инфоморфизм} \citep{barwise1997information}: пара
$f:G_A\!\to\!G_B$, $\hat f:\Sigma_B\!\to\!\Sigma_A$ с условием
$I_B(f g,\,y)\Leftrightarrow I_A(g,\,\hat f y)$. Инфоморфизмы компонуются и
имеют копределы (каналы Барвайза--Селигмана), поэтому коллективная память над
гетерогенными алфавитами --- это \emph{ядро канала} --- копредел диаграммы
инфоморфизмов, --- что дословно расширяет \S\ref{p4:sec:gluing} (функторы над
одним $\Sigma$ становятся инфоморфизмами над $\Sigma_i$; общий $\Sigma$ с
$\hat f=\mathrm{id}$ --- прежний случай).

\paragraph{Механизм заземления.} Символы выравниваются через общие референты:
на местах, уже заякоренных восстановлением системы координат
(\S\ref{p4:sec:semantic-identity}), совместная встречаемость $A$-тега и $B$-тега
--- свидетельство того, что они именуют один и тот же признак, и $\hat f$
считывается с заякоренной совместной встречаемости. (Когда алфавиты
различаются настолько, что сами отпечатки не могут посеять якоря, начальная
инициализация идёт от алфавитно-независимых структурных констелляций с
чередованием восстановления $f$ и $\hat f$.) \emph{Дефект выравнивания
символов} $\mathrm{def}_\Sigma$ --- доля заякоренных пар
$(\text{место},y)$, нарушающих бикондиционал инфоморфизма, --- символьный
аналог дефекта сопряжения систем координат из \S\ref{p4:sec:semantic-identity}.

\paragraph{Предварительный результат.} Контролируемое исследование
(\method{experiments/warp/symbol\_alignment.py}; $|\Sigma_A|{=}8$, $120$ мест,
$20$ сидов) даёт агенту $B$ секретную переразметку алфавита $A$ с шумом
наблюдения и восстанавливает $\hat f$ из заякоренной совместной
встречаемости.

\begin{table}[h]
\caption{Выравнивание символов заземлением против предположения об идентичных
алфавитах. Target-ID F1 --- идентификация целевого признака $A$ по
(переведённым) свидетельствам $B$. <<naive>> предполагает совпадение алфавитов
по индексу; <<recovered>> выучивает $\hat f$; <<oracle>> --- истинное
отображение.}
\label{p4:tab:symbol-alignment}
\centering
\small
\begin{tabular}{c c c c c c}
\toprule
шум & точн.\ перевода & $\mathrm{def}_\Sigma$ & F1 naive & F1 recovered & F1 oracle \\
\midrule
0.00 & 1.00 & 0.000 & 0.56 & 1.00 & 1.00 \\
0.10 & 1.00 & 0.099 & 0.55 & 0.92 & 0.92 \\
0.30 & 0.89 & 0.304 & 0.56 & 0.69 & 0.73 \\
\bottomrule
\end{tabular}
\end{table}

Предположение об идентичных алфавитах проваливается
($\text{F1}\approx0.55$: переразметка отправляет целевой тег $B$ на неверный
признак); заземление восстанавливает отображение (точность перевода $1.00$,
$\mathrm{def}_\Sigma$ отслеживает уровень шума) и достигает уровня оракула в
идентификации цели, устойчиво вплоть до $15\%$ заякоренных мест (точность
$0.95$). Огрубление двух признаков в один $B$-тег ограничивает достижимую
точность --- точной биекции не существует, поэтому ядро канала огрубляет, а
не ошибается. Выравнивание символов, таким образом, не полностью открытая
проблема, а тот же рецепт <<заземление плюс дефект>> на один алфавит выше;
полное рассмотрение (подлинная несоизмеримость, онлайн-дрейф) --- предмет
следующей статьи.

\section{Ограничения и следующие эксперименты}
\label{p4:sec:limitations}

\textbf{Грид-субстрат.} Текущий результат намеренно контролируем: все
эксперименты проведены в грид-мирах. Субстрат с непрерывным состоянием, такой
как VMAS, Habitat, ProcTHOR или MiniWorld, --- следующий шаг переносимости.

\textbf{Координатный якорь.} Устранён в
разделе~\ref{p4:sec:semantic-identity}: локальные сигнатуры-отпечатки
точно восстанавливают межагентную систему координат (сдвиг и повороты
на $90^\circ$), и результаты по слиянию/доверию/устареванию не
меняются поверх восстановленных систем. Открытым остаётся более
трудный конец той же оси: непрерывное $SE(2)$, шум внешнего вида
(текущий сопоставитель использует точное равенство констелляций,
поскольку субстрат бесшумен). Выравнивание символов --- агенты с
\emph{разными алфавитами тегов} --- предварительно рассмотрено в
\S\ref{p4:sec:symbol-alignment} (заземление через общие места,
восстановление в виде инфоморфизма); подлинная несоизмеримость и дрейф
словаря остаются открытыми.

\textbf{Минимальная CSM.} Правила CSM намеренно просты. Адаптивная
каденция $K_t$ и доверие, обусловленное намерением, --- следующие
архитектурные варианты. У селективной консолидации недопредставленных
тегов теперь есть измеренный пилот: в сопутствующей статье о переносе
маршрутов чисто реляционная \emph{песня} (сигнатуры ориентиров $+$
такт-векторы, без глобальных координат) выполняет тот же перенос
никогда не виденной цели, что и полный снапшот, за $6.4\%$ битовой
стоимости, причём безопасно --- одни такты отказывают открытым
образом, одни сигнатуры не достигают невиданного. Замена обмена
снапшотами на бюджетированные песни, сливаемые теми же тремя правилами
CSM, перестала быть спекуляцией; открытой остаётся вторая половина ---
сделать грамматику песни (что триггерит куплет, что входит в
сигнатуру) предметом отбора, а не дизайна.

\textbf{Обучаемые базовые варианты.} Обучаемая коммуникация в духе
CommNet --- полезное сравнение, но более релевантный будущий базовый
вариант --- многоагентная система с явной LLM-памятью или векторной
памятью, измеренная на той же задаче.

\section{Заключение}

Коллективная память не монотонна. В дефицитной многоагентной
навигации мгновенный обмен информацией может снизить завершение,
заставляя агентов сходиться на одних и тех же целях. Пиринговая
каденция обнажает этот компромисс, но одна лишь каденция --- одномерная
поверхность управления. Коллективная семантическая память должна
сделать слияние, доверие и устаревание явными. Минимальная CSM делает
ровно это и переводит систему в новую рабочую точку: не насыщение M
при быстрой рассылке и не насыщение C при медленной, а более
сбалансированный профиль M$\times$C с более высоким успехом, чем у
каждой фиксированной пиринговой каденции на бенчмарке дефицита.

\begin{ack}
Рукопись подготовлена в анонимизированной форме для рецензирования.
Описанные эксперименты выполнены на том же бенчмарк-конвейере, что и в
сопутствующей статье об измерениях Q/R/M/C, но архитектурное
утверждение здесь ограничено коллективной семантической памятью.
\end{ack}

{\small

}

\FloatBarrier
\begin{center}\rule{0.55\linewidth}{0.4pt}\\[2pt]\textbf{Приложения к этой части}\end{center}


\section{Коалгебраический слой: индивиды как ко-данные, коллектив как финальная коалгебра}
\label{p4:sec:coalgebra}

Прочтение из \S\ref{p4:sec:categorical} фиксирует \emph{содержание}
убеждений: в данном состоянии решётка концептов
$L_t^i=\operatorname{Concepts}(\mathcal{K}_t^i)$ говорит, что извлекаемо,
а слияние из \S\ref{p4:sec:colimit} --- как содержания соседей комбинируются
в $\widetilde{\mathcal K}_t^i$. Оба \emph{алгебраичны}: они строят объект
убеждения. Ни одно из них не говорит об агенте как о незавершающемся
\emph{процессе}, который обязан выдавать действие и сигнал на каждом
такте в движущейся среде, и ни одно --- о том смысле, в котором
\emph{популяция} приходит к общей конвенции. Эта процессная сторона ---
\emph{коалгебраическая} (coalgebra)~\citep{rutten2000universal,jacobs2016coalgebra}.
Мы делаем её точной на примере поиска воды и показываем, что две стороны
двойственны и встречаются внутри одного перехода состояния. (Словарь
убеждение/желание/намерение следует традиции BDI-агентов
\citep{rao1995bdi}, здесь наделённой коалгебраическим содержанием.)

\paragraph{Сигнатура.} Зафиксируем множество действий $\mathrm{Act}$
(двигаться, копать, ничего не делать), множество наблюдений $O_i$ на
агента (локальный перцепт агента $i$ --- его \emph{уникальная точка
зрения}) и множество \emph{остензий}
\[
  M \;=\; 1 \,+\, (X\times\Sigma),
\]
значения которого --- либо \emph{молчание} (слагаемое $1$), либо
указатель <<место $x$ несёт тег $\sigma$>> (свистнуть и указать на
клетку с меткой воды). То, что молчание --- полноправное значение $M$,
существенно: остензивный канал \emph{может быть информативно пуст}, и он
отделён от обновления состояния ниже. Это ровно та наблюдаемость,
которую разрушает обучаемый вещественный вектор сообщения (его канал
никогда не пуст и никогда не отделяет сигнализацию от шага), и именно
поэтому символьные стадии остаются читаемыми.

\paragraph{Состояние.} В качестве носителя коалгебры возьмём
высокоуровневое состояние H-MDP из \S\ref{p4:sec:categorical},
\[
  \mathcal{S}_i \;=\; \big\{\, h=(\mathcal{K}^i,\,D^i,\,g^{i,\star},\,s^i) \,\big\},
\]
несущее ту же информацию, что и $(L^i,D^i,g^{i,\star},s^i)$, поскольку
$L^i=\operatorname{Concepts}(\mathcal K^i)$. Таким образом,
\emph{убеждение живёт в состоянии}; решётка --- его проекция.

\paragraph{Определение (ко-данные: котип агента).}
Наблюдаемое поведение агента $i$ --- это коалгебра для эндофунктора
\[
  B_i(S)\;=\;\underbrace{\mathrm{Act}}_{\text{действовать сейчас}}\;\times\;
            \underbrace{M}_{\text{сигналить сейчас}}\;\times\;
            \underbrace{S^{\,O_i\times M^{A\setminus\{i\}}}}_{\substack{\text{продолжить при (мой следующий}\\ \text{перцепт, сигналы соседей)}}},
\]
т.е.\ отображение $c_i=\big\langle\, \method{observe\_act}_i,\ \method{observe\_ostension}_i,\ \method{step\_intension}_i \,\big\rangle : \mathcal{S}_i \to B_i(\mathcal{S}_i)$ с тремя деструкторами
\[
  \method{observe\_act}_i:\mathcal{S}_i\to\mathrm{Act},\quad
  \method{observe\_ostension}_i:\mathcal{S}_i\to M,\quad
  \method{step\_intension}_i:\mathcal{S}_i\to \mathcal{S}_i^{\,O_i\times M^{A\setminus\{i\}}} .
\]
Здесь $M^{A\setminus\{i\}}=\prod_{j\neq i}M$ --- кортеж остензий соседей,
доставленных на этом такте. Операционально
$\method{step\_intension}_i(h)(o,(m_j)_{j\neq i})$
(i)~вкладывает перцепт $o$ и указатели соседей $(m_j)$ в контекст через
слияние с доверием/устареванием,
$\widetilde{\mathcal K}_t^i=\operatorname{Merge}_{\mathrm{trust,stale}}(\cdots)$
из \S\ref{p4:sec:colimit}; (ii)~пересчитывает решётку; (iii)~переоценивает
предикат желания $D^i$; и (iv)~перефиксирует намерение
$g^{i,\star}=\chi_t^i(\operatorname{ext}(D^i))$ из \S\ref{p4:sec:categorical}.
Итак, $\operatorname{Merge}$ и $\chi$ --- \emph{компоненты одного
деструктора}. Агент \emph{и есть} ко-данные, порождаемые $c_i$: у него
нет конечного терминального плана, лишь постоянная способность выдавать
действие и остензию и продолжать. Это <<бесконечный генератор
выживания>> из мотивирующего примера.

\paragraph{Определение (индексированное семейство).}
Система --- это $A$-индексированное семейство
$\{(\mathcal{S}_i,c_i)\}_{i\in A}$, коалгебра в $\mathbf{Set}^{A}$.
Функтор $B_i$ зависит от индекса $i$ через $O_i$ (и через множество
соседей $A\setminus\{i\}$): один и тот же котип выживания
инстанцируется у каждого агента через локальную точку зрения этого
агента. Агент~1 на дюне порождает из <<я вижу горизонт>>; агент~2 в
низине порождает из <<я чувствую почву>>. Индекс --- это точка зрения.

\paragraph{Определение (козависимость: коллективная коалгебра).}
Пусть среда доставляет совместный перцепт $o=(o_i)_{i\in A}\in\prod_i O_i$.
Коллективное однотактовое отображение на $\prod_i\mathcal{S}_i$
подключает выход остензии каждого агента ко входам шага остальных,
\[
  \Phi\big((h_i)_i\big)(o)\;=\;
  \Big(\ \method{step\_intension}_i(h_i)\big(o_i,\,(\method{observe\_ostension}_j(h_j))_{j\neq i}\big)\ \Big)_{i\in A}.
\]
Поскольку следующее состояние $i$ есть функция \emph{выходов
деструкторов} остальных, семейство --- не $N$ независимых процессов, а
одна коалгебра для функтора-произведения
$\mathbf{B}(S)=\big(\prod_i(\mathrm{Act}\times M)\big)\times S^{\prod_i O_i}$
на $\prod_i\mathcal{S}_i$. Эта проводка и есть \emph{козависимость}:
агент~1 не может сделать шаг, не пропустив остензию агента~2 через свой
собственный $\method{step\_intension}$. Каденция рассылки входит сюда
как вентиль: на тактах без рассылки доставляемый кортеж $(m_j)$ --- это
последний полученный (устаревший) кортеж или молчание; $K$ --- период
доставки.

\paragraph{Определение (поведение и бисимуляция).}
$B_i$ --- контейнерный (полиномиальный) функтор, поэтому у него есть
финальная коалгебра
$(\nu B_i,\ \mathrm{beh}_i)$~\citep{aczel1988nonwellfounded,rutten2000universal},
элементы которой --- поведенческие деревья с ветвлением по
$\big(O_i\times M^{A\setminus\{i\}}\big)$ и метками
$(\mathrm{Act}\times M)$, и существует единственный морфизм коалгебр
$!_i:\mathcal{S}_i\to\nu B_i$, отправляющий состояние во всё его будущее
поведение. \emph{$B_i$-бисимуляция} (bisimulation) --- это отношение
$R\subseteq\mathcal{S}_i\times\mathcal{S}_i$, при котором связанные
состояния выдают одинаковое действие и одинаковую остензию и на каждом
входе переходят в $R$-связанные состояния. Контейнерные функторы
сохраняют слабые обратные образы, поэтому \emph{наибольшая бисимуляция}
$\sim_i$ существует, является эквивалентностью и равна $\ker(!_i)$: два
состояния бисимилярны тогда и только тогда, когда у них одинаковое
поведение~\citep{sangiorgi2011bisimulation}. То же верно для
коллективной коалгебры
$(\prod_i\mathcal{S}_i,\ \langle(\method{observe\_act}_i)_i,(\method{observe\_ostension}_i)_i,\Phi\rangle)$
с наибольшей бисимуляцией $\approx$.

\paragraph{Определение (общество / культура).}
\emph{Коллектив} --- это образ состояний в финальной коалгебре
$\nu\mathbf{B}$, эквивалентно фактор
$\big(\prod_i\mathcal{S}_i\big)/\!\approx$. Это ни центральный сервер,
ни общая карта: это класс бисимилярности совместного поведения.
\emph{Конвенция поиска воды} --- этот инвариант, прочитанный на водной
подсигнатуре. Пусть $M_{\mathrm{water}}\subseteq M$ --- указатели, чей
тег лежит в водном семействе, и пусть
$\nu B_i\to\nu B^{\mathrm{w}}$ --- забывающее отображение в
поведенческие деревья, ограниченные приёмом $M_{\mathrm{water}}$ и
выдачей соответствующей фиксации намерения. Популяция \emph{держит
конвенцию}, когда достижимые поведения всех агентов согласуются под
этим отображением: услышав <<укажи-и-копай у отмеченного водой места>>,
деструктор намерения каждого агента фиксирует это место (с точностью до
занятости). Управление тогда автоматично в обоих направлениях. Новый
агент (четвёртый агент из примера) --- член общества ровно тогда, когда
его поведение попадает в класс $[\cdot]_{\approx}$ на
$M_{\mathrm{water}}$; иначе его остензии не понимаются и его собственный
шаг никогда не сходится (сверху-вниз); а сам класс был достигнут
локальными, связанными шагами первых агентов у одного зелёного куста
(снизу-вверх).

\paragraph{Замечание (коммуникация = ленивая корекурсия глубины один).}
Поведение $!_i(h)$ определяется бесконечным развёртыванием $c_i$.
Получив $m_j=\method{observe\_ostension}_j(h_j)$, агент $i$ использует
внутри $\method{step\_intension}_i$ только этот выход деструктора
\emph{глубины один}; он \emph{не} вычисляет $!_j(h_j)\in\nu B_j$ ---
полное поведение соседа, каковое и есть регресс <<я думаю, что ты
думаешь\ldots>>. Коммуникация --- это, следовательно, корекурсия,
усечённая до наименьшей глубины, разрешающей собственный следующий шаг
получателя, --- конечное приближение отображения в финальную коалгебру
--- с релевантностно-теоретическим правилом остановки
\citep{sperber1995relevance}, решающим эту глубину. Регресс схлопывается
в одно эффективное действие.

\paragraph{Замечание (мост алгебра--коалгебра к \S\ref{p4:sec:categorical}--\ref{p4:sec:colimit}).}
Убеждение алгебраично: FCA-решётка $L_t^i$, построенная копредельным
слиянием, --- это \emph{содержание в состоянии}. Агент коалгебраичен:
ко-данные, наблюдаемые тремя деструкторами, --- это \emph{процесс над
состояниями}. Они двойственны и встречаются в одном месте ---
копредельное слияние и удовлетворяющий выбор $\chi$ суть внутренности
$\method{step\_intension}$. Диаграмма роста памяти
$\mathcal{K}_0^i\to\mathcal{K}_1^i\to\cdots$ из \S\ref{p4:sec:categorical}
--- это в точности проекция-на-убеждения траектории коалгебры
$h_0^i\mapsto h_1^i\mapsto\cdots$.

\paragraph{Коалгебраическое прочтение сдвига узкого места.}
Результат \S\ref{p4:sec:cadence} и таблицы~\ref{p4:tab:csm-results} читается
здесь одной строкой. Соревнуются два вида бисимилярности.
Бисимилярность на подсигнатуре \emph{остензий / убеждений} --- желанный
инвариант: это общая водная конвенция. Бисимилярность на сыром выходе
\emph{намерений} над \emph{дефицитной} целью --- нежеланный: несколько
состояний, отображающихся в одно и то же $g^{\star}$, --- это коллизия.
Занятость $\mathcal{O}_{-i}$ --- внекоалгебраическое ограничение,
штрафующее второе. Каденция $K$ настраивает, насколько сильно связка
$\Phi$ гонит коллектив к \emph{полной} наибольшей бисимуляции
$\approx$: малое $K$ доставляет свежие остензии часто, расширяя слитый
водный объём, так что намерение фиксируется надёжно
($P(M^\star\mid R^\star)\uparrow$), но пересинхронизирует деструктор
намерения, так что агенты сталкиваются
($P(C^\star\mid M^\star)\downarrow$). Внутренний оптимум $K{=}8$ ---
это точка, где убеждения сходятся, а намерения не схлопываются на одну
клетку; минимальная CSM добавляет доверие и устаревание именно затем,
чтобы $\method{step\_intension}$ достигал бисимилярности убеждений
(свежие, доверенные, выровненные свидетельства), гася бисимилярность
намерений (устаревшие свидетельства затухают прежде, чем успеют согнать
всех агентов на одну и ту же фиксацию). Мы формулируем это как
прочтение, согласующееся с данными, а не как теорему: контейнерные
функторы, их финальные коалгебры и бисимуляция --- стандартные
конструкции, используемые здесь как линза, а
релевантностно-теоретическое правило остановки --- интерпретирующая
глосса, а не подогнанный механизм.

\begin{table}[t]
\caption{Термины мотивирующего примера как формальные объекты
коалгебраического прочтения. Каждая строка --- смысл конструкции, уже
использованной в \S\ref{p4:sec:coalgebra}, в задаче поиска воды.}
\label{p4:tab:coalg-dictionary}
\centering
\small
\begin{tabular}{p{2.9cm} p{4.3cm} p{4.4cm}}
\toprule
Неформальный термин & Формальный объект & Смысл в поиске воды \\
\midrule
Котип (ко-данные) & $B_i$-коалгебра $(\mathcal{S}_i,c_i)$ & незавершающийся генератор действий/сигналов агента \\
Индексированный тип & семейство над $A$; зависимость $B_i$ от $O_i$ & уникальная локальная точка зрения агента (дюна против низины) \\
Козависимый тип & проводка $\method{observe\_ostension}_j\!\to\!\method{step\_intension}_i$ в $\Phi$ & <<я шагаю только через твою остензию>> \\
Индивид (экземпляр) & состояние $h\in\mathcal{S}_i$ и его поведение $!_i(h)$ & работающий агент, копающий и свистящий сейчас \\
Деструкторы (когенераторы) & $\method{observe\_act}$, $\method{observe\_ostension}$, $\method{step\_intension}$ & шагни-и-копай; свистни-и-укажи; вложи сигнал соседа в следующее состояние \\
Ленивое вычисление & корекурсия глубины один + релевантностная остановка & развернуть коалгебру соседа на один шаг, а не всё его поведение \\
Финальное коиндуктивное ядро / общество & $\nu\mathbf{B}$; наибольшая бисимуляция $\approx$ (на $M_{\mathrm{water}}$) & стабильная распределённая культура поиска воды \\
\bottomrule
\end{tabular}
\end{table}


\section{Команды воспроизведения}
\label{p4:app:repro}

\begin{itemize}[leftmargin=16pt]
\item Основная развёртка по каденции:
\path{python experiments/big_experiment/exp_cadence_phase.py --mode full --workers 8 --out_dir tmp/big_experiment_qrmc_40}
\item Расширенные каденции:
\path{python experiments/big_experiment/exp_extra_K.py --workers 8 --out_dir tmp/big_experiment_extraK}
\item Анализ Q/R/M/C:
\path{python experiments/big_experiment/analyze_qrmc.py --runs_csv tmp/big_experiment_qrmc_40/runs.csv --out_dir tmp/big_experiment_qrmc_40}
\item Бенчмарк минимальной CSM:
\path{python -m experiments.collective_semantic_memory.run_csm_benchmark --out_dir tmp/csm_benchmark}
\item Развёртка переносимости на MiniGrid:
\path{python -m experiments.minigrid_multiagent_wrapper.run_portability_sweep --out_dir tmp/minigrid_portability}
\item Дефект кругового обхода сопряжения при восстановлении системы координат (раздел~\ref{p4:sec:semantic-identity}):
\path{python experiments/warp/alignment_defect.py}
\item Структурная фазовая диаграмма доверие$\times$каденция (раздел~\ref{p4:sec:phase-diagram}, рис.~\ref{p4:fig:phase-diagram}):
\path{python experiments/collective_semantic_memory/phase_diagram_trust_cadence.py --fig figures/fig_phase_diagram.pdf}
\end{itemize}

\section{Замечание об активах и лицензиях}
\label{p4:app:assets}

В экспериментах используются MiniGrid и MiniWorld~\citep{MinigridMiniworld23}. Локально реализованный код бенчмарка и памяти будет выпущен под лицензией MIT в составе артефакт-пакета.

\clearpage
\FloatBarrier
\FloatBarrier

\clearpage
\part{Грамматика песни и Utility-Certified Memory: онтогенез запоминаемого}
\noindent\textit{Формирование памяти как двухосевое решение (контрфактическая полезность × простота аналогии): измерено детерминированно, выучено бандитом, обыгравшим проектировщика, эволюционировано к конечному бюджету и расценено против семи политик на гетерогенных воплощениях.}

\paragraph{Аннотация части.} \textbf{Проблема.} Предыдущие части серии закрыли
\emph{эксплуатацию и координацию} уже сформированной распределённой
памяти (контракты провенанса, единицы переноса, правила
merge/trust/staleness). Открытым остался её \emph{онтогенез}: как
агенты в ходе совместной деятельности обучаются тому, что достойно
запоминания, что становится ориентиром и какой фрагмент опыта ---
песней.
\textbf{Модель.} Мы определяем формирование памяти как ограниченный
выбор из пяти операций --- слияние, исключение, новая схема, повтор,
удаление, --- управляемый двумя \emph{независимыми} осями:
контрфактической маргинальной полезностью $U(m \mid M) =
\mathrm{cost}(M) - \mathrm{cost}(M \cup \{m\})$, измеряемой
детерминированным replay с воспоминанием, включённым и выключенным, и
простотой структурной аналогии (перекрытие сигнатур плюс
\emph{decision distortion} --- куда две песни в итоге посылают
идущего). Хранимые схемы несут \emph{сертификаты полезности}
(условия, наблюдавшийся выигрыш, профиль полезности по ролям,
свидетельства, отказы); приёмники пересчитывают собственную
полезность, а не доверяют числу отправителя.
\textbf{Результаты.} Детерминированный прототип проходит все четыре
зарегистрированных предсказания: дубликаты бесплатны по построению,
операция исключения бьёт и слияние с перезаписью ($12{,}3$ против
$62{,}6$ шага на испорченном мире), и recency-затенение, а сходство
измеримо не является критерием консолидации. Контекстный бандит
\emph{выучивает} матрицу решений из одного координационного фидбека и
превосходит спроектированные пороги ($78{,}1$ против $83{,}5$ против
случайного $105{,}9$); его единственное принципиальное отклонение ---
добавлять вместо перезаписывать --- ровно то, что и окупается.
Эволюционный цикл над грамматикой песни отбирает \emph{конечный}
бюджет куплетов без регрессии стоимости и сходится к окрестности
спроектированных порогов. Сравнение семи политик на гетерогенных
воплощениях при $10/100/1000$ эпизодах заканчивается четырьмя
зарегистрированными провалами, механизмы которых и есть результат:
структурная память бьёт эмбеддинговое сходство,
recency-importance-relevance и полные снапшоты в $1{,}7$--$4{,}5$
раза на каждом масштабе; консолидация покупает \emph{биты}, а не
успех ($154$ тыс.\ против $339$ тыс.\ у raw и $5{,}2$ млн у
снапшотов на $1000$ эпизодах); а отказ длинного горизонта ---
проблема \emph{устаревания} с измеренной лестницей композиции:
консолидация $-14\%$, $+$свежесть $-3\%$, $+$trust-flip $-0{,}9\%$
--- живучая роль закрыта против raw-контекста.
\textbf{Поздние волны.} Цикл повышения ориентиров замкнут причинно
над суперсетом кандидатов (нарочно вредный мусор пойман при
$\Delta=-0{,}67$; граничный признак с нулевой атрибуцией оказался
каузально несущим --- атрибуция и причинная необходимость суть разные
величины); кросс-семейный кодбук сигнатур монетизирует повторяющиеся
мотивы ($-16\%$ бит на $1000$ эпизодах, с ростом от масштаба);
graph-matching движок аналогии режет память на $24\%$ при равной
стоимости; meta-обучение с только-терминальной наградой бьёт и
плотно-шейпленного бандита, и ручную матрицу; end-to-end цикл с живой
LLM показывает превосходство реконструкции над проигрыванием ---
сырые транскрипты деградируют с $0{,}42$ до $0{,}08$
(\texttt{llama3.1}) и с $1{,}0$ до $0{,}375$ (\texttt{Qwen2.5-3B})
с ростом конфликтной истории, тогда как реконструированный активный
контекст держит $1{,}0$ за $6$--$12\%$ payload; а социальный
эксперимент возвращает сильнейший честный негатив программы: на
длинном горизонте в глобально дрейфующем мире независимые агенты бьют
каждое плечо обмена --- время трансляции не есть время мира, закон
серии «гейт, а не обмен» перевыведен на уровне формирования.
\textbf{Тезис.} Что запоминать --- решает не описательная
заметность, а контрфактическая полезность для будущих индивидуальных
и коллективных решений, скорректированная на стоимость структурного
переноса между локальными мирами агентов; и само правило решения
выучиваемо, а параметры грамматики под ним --- эволюционируемы в
пределах фиксированного семейства грамматик.
\textbf{Границы.} Сеточные субстраты; детерминированный replay;
измерительная программа с зарегистрированными предсказаниями, а не
общая теория.

\FloatBarrier
\section{Введение: от координации к онтогенезу}
\label{p5:sec:intro}

Серия до сих пор имела дело с памятью, которая уже существует.
Аннотация провенанса расценивает чужие свидетельства в момент
обязательства; песня переносит маршруты и идентичность без общих
систем координат; тройка merge/trust/staleness заставляет коллектив
локальных памятей координироваться. Чего не делает ни один из этих
механизмов --- он не \emph{формирует} память: не решает, какое
наблюдение становится записью, какая запись --- схемой, какой признак
зарабатывает статус ориентира и что вообще стоит транслировать.
Острейший одноагентный диагноз серии --- генерация кандидатов
голодает на $91\%$ точек решения при почти оптимальном ранжировании
--- это в точности проблема формирования, и список открытых проблем
рукописи-предка называет селективную консолидацию её прямой атакой.

Эта часть формулирует проблему формирования как ограниченное
решение, даёт ей измеренное детерминированное решение, а затем
подвергает и правило решения, и грамматику под ним обучению и
селекции. Всюду сохраняется методологическое обязательство серии:
предсказания регистрируются до прогонов, а провалы публикуются с
механизмами --- несколько самых информативных результатов ниже ---
именно зарегистрированные провалы.

\paragraph{Ключевая идея.} Полезность и простота аналогии --- две
\emph{независимые} оси, и никакой одиночный скаляр (новизна,
сходство, важность, награда) не заменяет их произведение:

\begin{center}
\small
\begin{tabular}{lll}
\toprule
Полезность & Аналогия & Операция памяти \\
\midrule
высокая & простая & \textbf{MERGE} в существующую схему \\
высокая & простая, но противоречит решению & \textbf{EXCEPTION}:
контрпример, не перезапись \\
высокая & сложная & \textbf{NEW\_SCHEMA} \\
низкая & простая & \textbf{REPEAT}: обновить статистику \\
низкая & сложная & \textbf{DROP} \\
\bottomrule
\end{tabular}
\end{center}

Пятая ячейка --- несущая: полезное событие, противоречащее схеме,
нельзя ни слить (испортит), ни удалить (это свидетельство); оно
хранится как исключение, пока схема не расщепится. Это и защищает
память от ложной генерализации.

\section{Модель}
\label{p5:sec:model}

\paragraph{Контрфактическая маргинальная полезность.} Для
воспоминания $m$, агента $i$, намерения $\iota$ и текущей памяти $M$:
\[
U_{i,\iota}(m \mid M) \;=\;
\mathrm{cost}_i(M, \iota) \;-\; \mathrm{cost}_i(M \cup \{m\}, \iota),
\]
измеряется \emph{детерминированным replay} с включённым и выключенным
воспоминанием --- приём source-masked replay из варп-анализа,
обобщённый с «чужое свидетельство вкл/выкл» до «это воспоминание
вкл/выкл». Полезность контекстна (по агенту и намерению) и
маргинальна (относительно уже хранимого): дубликат хранимой схемы
имеет $U = 0$ \emph{по построению}, так что повторы ловятся без
всякой эвристики сходства. Командная версия вычитает стоимости
коллизий, коммуникации и ложного переноса.

\paragraph{Стоимость аналогии.} Новый эпизод $e$ отображается на
хранимую схему $s$ через $\phi$ со стоимостью
\[
C_{\mathrm{analog}}(e,s) = L(\phi) + \lambda_d\,
D_{\mathrm{decision}}(e,s) + \lambda_a \operatorname{Defect}(\phi),
\qquad A(e,s) = e^{-C_{\mathrm{analog}}},
\]
где $L$ считает структуру вне соответствия (здесь: куплеты вне
сигнатурного LCS), $D_{\mathrm{decision}}$ меряет, насколько
по-разному две песни \emph{решают} --- операционализировано как
расстояние между их сквозными смещениями по тактам, т.\,е.\ куда они
в итоге посылают идущего, --- а $\operatorname{Defect}$ штрафует
нарушение реляционной структуры. Описательное сходство и совпадение
решений --- разные величины, и $D_{\mathrm{decision}}$ ловит вторую;
это rate-distortion-интуиция \citet{rdd2026}, перенесённая в
многоагентную структурную постановку.

\paragraph{Единица памяти --- условная программа, не факт:}
$m = (G_m, C_m, E_m, U_m, \Phi_m, P_m, F_m)$ --- реляционный
граф/песня, условия применимости, ожидаемые эффекты, профиль
полезности \emph{по ролям}, известные аналогии, провенанс со ссылками
на неизменяемые исходные эпизоды и известные отказы. При обмене схема
едет с \textbf{сертификатом полезности} $(\text{условия}, \Delta V,
\text{неопределённость}, \text{свидетельства}, \text{аналогия},
\text{отказы})$; приёмник переносит песню через собственное
отображение и пересчитывает собственную полезность. Для гетерогенных
воплощений это принципиально: маршрут через поле опасностей может
одновременно иметь $U_{\mathrm{robust}} > 0$ и
$U_{\mathrm{fragile}} < 0$, поэтому воспоминание несёт функцию
полезности по ролям, а не один вес.

\paragraph{Четыре слоя хранения.} (1) \emph{неизменяемый эпизодный
стор} --- исходные эпизоды не переписываются (непрерывная перезапись
сначала помогает, потом вредит \citep{faulty2026}; первичные
свидетельства --- якорь реконсолидации); (2) \emph{граф аналогических
схем} --- песни, ориентиры, исключения под двухосевой матрицей;
(3) \emph{процедурный слой} для схем, продвинутых во фрагменты
политики; (4) \emph{координационный слой} для социальных правил
(резервации, роли, доверие, условия передачи). Активный контекст
агента --- реконструкция: релевантные схемы, минимальные
подтверждающие эпизоды, актуальные исключения и координационный
контракт, а не проигрывание истории.

\section{U1: матрица, измеренная детерминированно}
\label{p5:sec:u1}

Восемь валидных семейств сеточных миров, пять кандидатов на семейство
--- по одному на ячейку матрицы, последовательно: полезная новая
песня (\textsc{new}), её точный дубликат (\textsc{repeat}), песня из
структурно другого мира (\textsc{novel}), песня к цели, нерелевантной
water-намерению (\textsc{irrelevant}), и песня из варианта мира с
тайно перенесённой водой (\textsc{conflict}). Два одноосевых
бейзлайна на том же потоке: similarity-only (хранить новое, слияние с
перезаписью для похожего) и utility-only (хранить полезное,
потреблять newest-first). Зарегистрированы четыре предсказания;
первая регистрация провалилась на двух дефектах конструкции
(«нерелевантная» цель случайно лежала по пути к воде; потребители
спавнились в стенах) --- оба задокументированы, исправленная
регистрация прошла $4/4$:

\begin{itemize}[leftmargin=16pt,itemsep=2pt,topsep=2pt]
\item \textbf{Повторы бесплатны по построению:} маргинальная
  полезность дубликата --- ровно $0$ в каждом мире; итоговая память
  --- $3$ схемы против $5$ у «добавлять всё».
\item \textbf{Сходство --- не критерий консолидации:}
  similarity-бейзлайн инкорпорирует нерелевантную песню и молча
  портит исходный мир, перезаписав его схему конфликтным вариантом
  ($62{,}6$ шага и phantom-first против \textbf{$12{,}3$ шага
  first-try} у двухосевой памяти).
\item \textbf{Исключения бьют и перезапись, и recency:} с хранимым
  исключением вариантный мир стоит $20{,}5$ против $70{,}4$ без него,
  а исходный остаётся корректным с первой попытки; recency-порядок
  utility-only фантомит первым на исходном ($24{,}0$ шага).
\item \textbf{Сертификаты пересчитываемы:} приёмник в другой позиции
  восстанавливает положительную полезность из свидетельств
  сертификата в каждом мире.
\end{itemize}

\section{U2: матрица выучивается --- и улучшает себя}
\label{p5:sec:u2}

Контекстный бандит LinUCB заменяет ручные пороги. Признаки: $(U,
\text{share}, D, \text{новизна}, \text{размер памяти})$; действия ---
пять операций; награда --- \emph{измеренное} изменение стоимости
памяти на двухмировой пробе (мир кандидата и мир происхождения
ближайшей схемы --- второй член учит контроллер не перезаписывать),
минус небольшой налог на хранение.

Первая регистрация провалилась на двух дефектах оценки, которые сами
--- результат. Финальная батарея была циркулярной (каждая политика
оценивалась на мирах, которые сама решила хранить, отчего random
неотличим от спроектированной), а две пары действий были
\emph{reward-эквивалентны}: \textsc{new} и \textsc{exception} имели
идентичное структурное следствие, как и \textsc{drop} с
\textsc{repeat} --- и бандит восстановил ровно эту фактор-структуру,
путая действия строго \emph{внутри} классов эквивалентности и выучив
drop-ячейку на $40/43$. Бухгалтерия без последствий необучаема.
Версия 2 сделала батарею независимой и дала support последствие
(порядок потребления). Результат:

\begin{center}
\small
\begin{tabular}{lccc}
\toprule
 & выученная & ручная матрица & random \\
\midrule
стоимость на held-out батарее & \textbf{78{,}1} & 83{,}5 & 105{,}9 \\
\bottomrule
\end{tabular}
\end{center}

Выученная политика обыгрывает собственного проектировщика на
$6{,}5\%$. Её мажоритарные действия совпадают с ручной матрицей в
$3/5$ ячеек (зарегистрированный промах); расхождения --- две
контекстно-зависимые эквивалентности (слить дубликат $\equiv$
повторить; повторить мусор $\equiv$ удалить) и одно
\emph{принципиальное прибыльное отклонение}: там, где ручная матрица
говорит \textsc{merge}, контроллер добавляет новую версию и отдаёт
арбитраж порядку по support --- он самостоятельно открыл принцип
неизменяемого стора (\S\ref{p5:sec:model}), тот же урок, к которому
литература о консолидации приходит из анализа отказов
\citep{trustmem2026,faulty2026}.

\section{U3: параметры грамматики эволюционируют --- и константы у аттрактора}
\label{p5:sec:u3}

Эволюционная стратегия ($30$ поколений $\times$ $16$ геномов)
отбирает грамматику $\theta = (\text{шаг куплетов},
\text{порог сходства}, \text{порог дисторсии},
\text{порог полезности}, \text{бюджет } v_{\max})$ с фитнесом на
картах, где геном никогда не пел (стоимость $+$ $\lambda\cdot$биты).
Оба зарегистрированных предсказания прошли: эволюционированный геном
на hold-out эквивалентен ручной грамматике в пределах допуска
регистрации и отбирает \emph{конечный} бюджет ($v_{\max} = 10$) без
регрессии стоимости --- ограниченная рациональность, полученная
селекцией, а не навязанная. Сверх регистрации: селекция сошлась к
окрестности спроектированных порогов (сходство $0{,}74 \to 0{,}46$
при ручном $0{,}4$; дисторсия $3{,}9$ при ручной $3$), выбрав более
плотные куплеты и более низкий порог полезности --- «хранить больше,
но короче». Спроектированные константы матрицы лежат у аттрактора
селекции, а не произвольны.

\section{U7: семь политик, два тела, три горизонта}
\label{p5:sec:u7}

Решающее сравнение прогоняет один смешанный поток эпизодов (новые
семейства, точные повторы, варианты внешности с пересеянной
hazard-текстурой над устойчивой структурой стен-и-воды и конфликты с
перенесённой водой) через семь политик командной памяти, потребляемых
двумя воплощениями --- \emph{хрупким} (опасность стоит $12$ шагов) и
\emph{живучим} (опасность $1$, шаг $2$), --- и оценивает на последних
истинных мирах семейств плюс свежих вариантах внешности, которых в
потоке не было. Политики: raw long context (хранить всё,
newest-first), эмбеддинговое сходство сигнатур,
recency-importance-relevance \citep{genagents2023}, обмен полными
снапшотами, utility-only, analogy-only и полная двухосевая память с
ролевыми сертификатами.

\begin{table}[h]
\centering
\caption{Средняя стоимость батареи для хрупкой роли по длинам потока
($12$ сидов) и хранимые биты на $1000$ эпизодах. Все четыре
зарегистрированных предсказания провалились; механизмы --- и есть
результат (см.\ текст).}
\label{p5:tab:u7}
\small
\begin{tabular}{lcccc}
\toprule
Политика & $10$ эп. & $100$ эп. & $1000$ эп. & биты @$1000$ \\
\midrule
raw long context & $128{,}9$ & $\mathbf{53{,}3}$ & $\mathbf{139{,}8}$ & $339\,453$ \\
utility-only & $130{,}5$ & $53{,}6$ & $144{,}6$ & $162\,107$ \\
двухосевая (UCSM) & $132{,}8$ & $60{,}8$ & $145{,}2$ & $\mathbf{153\,774}$ \\
сходство сигнатур & $161{,}1$ & $214{,}7$ & $232{,}1$ & $186\,843$ \\
recency-imp.-relevance & $190{,}7$ & $225{,}4$ & $250{,}1$ & $339\,453$ \\
полные снапшоты & $209{,}4$ & $246{,}6$ & $264{,}5$ & $5\,198\,886$ \\
\bottomrule
\end{tabular}
\end{table}

\paragraph{Что провалилось и что это значит.} Зарегистрированные
предсказания доминирования, сублинейных бит, сохранения ролей и
ложных слияний провалились все. Содержание несут три механизма.
\emph{Первый}: fail-closed слой идентичности делает раздутие дешёвым
в шагах --- устаревшие и чужие песни просто не порождают якорей, а
значит и целей, так что raw-архив бесплатно фильтруется в момент
потребления и платит только битами и вычислениями матчинга.
\emph{Второй}: сублинейный рост памяти невозможен без кросс-семейной
абстракции --- в открытом потоке, где $30\%$ эпизодов основывают
новые структуры, консолидация ограничивает рост внутри семейства
($2{,}2\times$ меньше бит, чем raw, $34\times$ меньше снапшотов), но
не может сжать новизну, для которой у неё нет схемы. \emph{Третий}:
отказ длинного горизонта --- phantom-first около $0{,}95$ у каждой
песенной политики на $1000$ эпизодах --- это каскад
\emph{устаревания} из накопленных конфликтов, который никакая
политика хранения не чинит одним порядком; это родная территория
закона доверие$\times$давность серии, которого ни одна из семи
политик не реализовывала.

\paragraph{Устойчивые позитивы.} На каждом масштабе структурная
память (идентичность $+$ такты, какой бы ни была консолидация) бьёт
мейнстримные эвристики --- эмбеддинговое сходство,
recency-importance-relevance \citep{genagents2023,memgas2025} и
снапшоты --- в $1{,}7$--$4{,}5$ раза; а двухосевая память дешевле
всех песенных политик по битам всюду.

\paragraph{U7b: проверка диагноза устаревания.} Восьмая политика ---
та же двухосевая консолидация, потребляемая newest-first --- была
зарегистрирована и прогнана. На $100$ эпизодах она \emph{проходит}:
совпадает или бьёт raw на обеих ролях ($52{,}0/67{,}9$ против
$53{,}6/68{,}9$) за $54\%$ его бит, закрывая весь разрыв накладных
расходов консолидации одним порядком. На $1000$ эпизодах она
проигрывает $\sim 3\%$: при сильном дрейфе внешности избыточность raw
сохраняет остаточную ценность --- старые варианты песни всё ещё
якорятся там, где новейшая не находит куплетов. Формула программы,
которая выживает: \textbf{консолидация покупает биты, свежесть
покупает успех, и на умеренных горизонтах они компонуются; следующий
рычаг --- настоящий возрастной гейт и кросс-семейные схемы, а не
лучший порядок.}


\section{Закрытие трёх разрывов: устаревание, эмерджентность, абстракция}
\label{p5:sec:gaps}

Вторая зарегистрированная волна атаковала три самых конкретных
разрыва, локализованных первой.

\paragraph{Устаревание в потреблении (U7c).} Контракт
доверие$\times$давность серии вшит в потребление: схемы упорядочены
по $\text{support}\cdot e^{-\alpha\,\text{age}}$ плюс дискретный
\emph{trust-flip} --- схема, накопившая исключение, понижается, и
контрпример пробуется раньше дискредитированного родителя. На $1000$
эпизодах это полностью закрывает живучую роль ($214{,}2$ против
$214{,}3$ у raw) и втрое сокращает разрыв хрупкой ($-0{,}9\%$ против
$-3\%$ у одной свежести); строгая регистрация по обеим ролям всё же
проваливается, честно --- остаточная ценность избыточности при дрейфе
внешности переживает и гейт. Лестница композиции теперь измерена:
одна консолидация $-14\%$, $+$свежесть $-3\%$, $+$trust-flip
$-0{,}9\%$.

\paragraph{Причинно повышенные ориентиры (E1).} Причинный цикл
\emph{предложить $\to$ аблировать $\to$ измерить $\to$ повысить}
замкнут: стартуя с суперсета признаков, включающего производного
кандидата и нарочно вредный frame-вариантный мусор (чётность клетки),
цикл поймал мусор причинно ($\Delta = -0{,}67$; вред проявлялся
отказами --- ноль fail-open за все раунды абляций) и дошёл до
минимального словаря с идеальным счётом ($0{,}33 \to 1{,}0$).
Зарегистрированное ожидание, что выживут контент-классы, провалилось
самым поучительным образом: выжившая пара --- $\{\texttt{стена},
\texttt{пустота}\}$: граничный признак, несущий \emph{ноль}
консенсусной массы в атрибуционном измерении, оказывается каузально
несущим как \emph{анти-алиасинговый страж}, когда словарь обеднён.
Каузальная абляция и атрибуционная масса --- разные величины;
ценность признака неаддитивна и зависит от остального словаря.

\paragraph{Кросс-семейная абстракция первого порядка (X1).} В кодек
хранения добавлен общий кодбук сигнатур --- констелляции,
повторяющиеся через семейства, хранятся однажды и адресуются
индексом, без потерь по построению. Зарегистрировано и пройдено: рост
бит кодбука ($7{,}91\times$ на декаду) меньше плоского кодека
($8{,}42\times$), а абсолютная экономия растёт с масштабом
($-10\%$ на $100$ эпизодах, $-16\%$ на $1000$): кросс-семейные мотивы
существуют и монетизируются. Сублинейность приближена, но не
достигнута --- полная абстракция уровня схем (слоты и роли, а не
только общие констелляции) остаётся открытой половиной.


\paragraph{LLM end-to-end цикл (L1).} Полный обещанный цикл ---
наблюдение $\to$ кандидат памяти $\to$ контроллер $\to$ граф схем
$\to$ реконструированный активный контекст $\to$ действие LLM ---
замкнут на живых моделях (разделение труда по W4: LLM коммитит
семантическую цель, детерминированный waypoint-примитив идёт).
Формирование --- детерминированный UCSM над историями из $4$ или
$12$ эпизодов, в которых воду переносят один или два раза, так что
устаревшие цели держат до $10{:}2$ большинства сырого транскрипта;
затем свежая LLM-сессия коммитит при трёх плечах контекста (без
памяти / полный транскрипт / граф схем с сертификатами и
исключениями). Зарегистрировано и измерено: память необходима
(безпамятное плечо --- $0{,}0$ всюду); реконструкция бьёт
проигрывание --- \texttt{llama3.1-8B} деградирует на сыром
транскрипте с $0{,}42$ до $0{,}08$ с ростом истории (конфликтные
цепочки $0{,}17 \to 0{,}00$), тогда как активный контекст держит
$1{,}0$ всюду за $6$--$12\%$ payload транскрипта;
\texttt{Qwen2.5-3B} читает короткий транскрипт идеально ($1{,}0$), но
падает до $0{,}375$ на $12$-эпизодных цепочках --- активный контекст
снова $1{,}0$. Четырёхмиллиардная thinking-модель проваливает цепочку
supersedes во всех плечах --- потолок способности ниже задачи,
сообщается ради честности. Режим отказа транскрипта --- именно
устаревшее большинство, а не длина (безобидные $12$-эпизодные логи
сильная модель читает): сырой контекст воспроизводит фантомный каскад
U7 в пространстве промпта, а реконструированный граф схем --- его
staleness-гейт в том же пространстве.


\section{Третий эшелон: структура, социум, терминальный кредит}
\label{p5:sec:echelon3}

Оставшиеся три разрыва атакованы в масштабе; один закрылся чисто,
один --- с промахом на волосок, один вернул сильнейший негатив
программы.

\paragraph{Граф-аналогия (G1).} LCS-прокси заменён структурным
выравниванием: Нидлман--Вунш по куплетам с частичным сходством узлов
(Жаккар по ключам сигнатур) и \emph{реляционным дефектом} ---
согласованностью цепочек тактов выровненных пар (идут ли две песни
одним скелетом?). На потоках, насыщенных вариантами внешности
($300$ эпизодов, $12$ сидов), граф-контроллер распознаёт «ту же
структуру в новой текстуре», которую LCS слепо заводил новыми
схемами: на $24\%$ меньше схем и бит ($122$ против $162$; $42$ тыс.\
против $55$ тыс.) при стоимости в пределах $5\%$ на обеих ролях и
без цены безопасности (фантомы $+0{,}012$). Зарегистрированный порог
по схемам ($\le 0{,}75\times$) промахнут на $0{,}01$ --- сообщается
как проваленная регистрация, каждое содержательное условие которой
выполнено.

\paragraph{Терминальный кредит (M1).} Бандит с немедленной наградой
отправлен в отставку: meta-обучение эволюционными стратегиями
оптимизирует тот же класс политик против \emph{только терминальной}
награды --- одно число за весь стрим, после десятков переплетённых
решений. На held-out политика бьёт и плотно-шейпленного бандита
($-97{,}95$ против $-99{,}69$), и ручную матрицу ($-99{,}95$), а на
\textsc{exception}-ячейках берёт store-класс в $91\%$ случаев без
replace-большинства --- третье независимое переоткрытие принципа
неизменяемого стора (после проектировщика и бандита), на этот раз из
одного отложенного скаляра.

\paragraph{Социальный процесс (S1) --- честный негатив.} Шесть
агентов смешанных воплощений, каждый со СВОЕЙ памятью, действуют
одновременно в мире, семейства которого эволюционируют глобально,
обмениваясь сертификатами на каденции. На $20$ эпизодах коммуникация
помогает; на $300$ --- \textbf{независимость выигрывает вчистую}
(хрупкий $148$ против $170$; живучий $171$ против $238$ у каждого
коммуницирующего плеча). Механизм: приёмники накапливают сотни чужих
схем ($741$ против $221$), чьи мировые референты устаревают быстрее,
чем их отфильтрует ролевой гейт или свежесть, --- время трансляции не
есть время мира. Три из четырёх регистраций на этом провалились;
прошедшая --- резервации (S1.3: дублирующиеся коммиты в ноль без
цены). Это центральный закон серии --- \emph{ценность коллективной
памяти несёт гейт, а не обмен} --- перевыведенный в субстрате
формирования: сертификаты без staleness-контракта на мировых часах
не просто недостаточны --- они чисто вредны на горизонте, ровно как
сырой обмен для LLM-коллектива в W6 части II.


\paragraph{Контракт на часах мира (S2) --- необходим, не достаточен.}
Диагноз S1 реализован: каждое семейство несёт версию состояния,
инкрементируемую при реальном изменении его мира; сертификаты несут
(семейство, версия); новости о версии путешествуют в метаданных
сертификата, даже когда сама схема отвергнута; допустимость гейтится
по часам мира и при потреблении, и при консолидации. На парных
потоках назначений ($12$ сидов, $300$ эпизодов) гейт улучшает каждое
коммуницирующее плечо (живучий $238 \to 220$, хрупкий $170 \to 163$),
а резервации по-прежнему зануляют дублирующиеся коммиты бесплатно ---
но реверс S1 не реверсируется: независимые агенты бьют гейтованное
плечо на \emph{каждом} сиде, а памяти почти не сжимаются ($688$
против $708$ записей). Гейт возвращает лишь ${\sim}27\%$ штрафа
коммуникации. Два зарегистрированных негатива теперь зажимают
проблему с обеих сторон: фильтрация по времени сообщений (S1) и по
времени мира (S2) обе не делают социальный обмен нетто-полезным на
горизонте. Остаточный механизм вообще не временной: это
\emph{кросс-семейный алиасинг от самого объёма чужого} --- сотни схем
о семействах, которые приёмник почти не посещает, всё равно частично
матчатся с его локальными констелляциями и эмитят фантомные цели.
Данные теперь называют контроль ДОПУСКА, а не устаревания: принимать
чужое знание лишь там, где оно показало полезность на собственных
визитах, --- сертификат есть свидетельство, а свидетельство, говорят
оба негатива вместе, не является допустимым доказательством.


\paragraph{Контроль допуска (S3) --- реверс реверсируется.} Принцип,
на который указали оба негатива, реализован: каждый чужой сертификат
по прибытии попадает в КАРАНТИН (метаданные о версии путешествуют
всё равно); при визите агента в семейство $f$ карантинные сертификаты
о $f$ валидируются НА МЕСТЕ --- маргинальная полезность чужой песни
проигрывается на реальном текущем мире собственной ролью агента;
прошедшие допускаются с полезностью, \emph{измеренной, а не
рассказанной}, и применимы уже в этом визите; провалившие
отбрасываются; знание о никогда не посещаемых семействах в
потребление не попадает вовсе. Все четыре зарегистрированных
предсказания проходят ($12$ сидов, парные потоки). \textbf{Социальный
обмен впервые нетто-полезен на горизонте}: полный контракт бьёт
независимость на каждом сиде и обеих ролях (хрупкий $118$ против
$148$, $-20\%$; живучий $134$ против $171$, $-21\%$). Механизм
монотонен --- свидетельство $162/220$ $\to$ допуск по одной
релевантности $136/169$ $\to$ допуск по продемонстрированной
полезности $118/134$, --- гигиена наконец материализуется (допущенная
память $233$ против $688$ записей, $0{,}34\times$, при $592$
мусорных сертификатах, безвредно лежащих в карантине), а резервации
компонуются (ноль дублирующихся коммитов за $+0{,}6\%$ цены).
Социальная лестница программы тем самым измерена из конца в конец:
обмен без контрактов вреден (S1), staleness на часах мира необходим,
но недостаточен (S2), а допуск по продемонстрированной полезности на
собственных визитах --- свидетельство не есть доказательство ---
закрывает рубеж (S3).


\section{Один метод: интегрированный рантайм и его бенчмарк}
\label{p5:sec:integration}

Каждый механизм выше был проверен в собственном контролируемом
эксперименте; рецензент вправе спросить, работает ли их композиция.
Поэтому построен \textbf{Songline Memory Runtime v1}: один тип записи
$m = (G, C, E, U, P, T, R, F)$ (граф песни, применимость, ожидаемый
эффект, полезность по ролям, провенанс, время и версия мира,
кадровые соответствия, восстанавливаемые при потреблении,
исключения), один цикл от наблюдения до исполнения, и каждый механизм
--- флаг конфигурации: армы и абляции бенчмарка суть конфигурации
одного артефакта, а не соседние кодовые базы. Учёт провода и
хранения --- полным кодеком (сертификат $64$, провенанс $48$,
время/версия $32$, резервация $24$ бита).

\paragraph{Интеграция одним прогоном (R1).} Десятипунктный тест
завершённости --- маршрут открыт в приватном кадре; песня получена;
сертификат удержан в карантине; допуск по собственной измеренной
полезности приёмника; соответствие ориентиров; исполнение маршрута;
rupture-откат после изменения мира; уступка по резервации;
устаревшие версии никогда не используются --- проходит в ОДНОМ
прогоне одного рантайма, $10/10$ сценариев.

\paragraph{Интеграционный бенчмарк (I1).} Шесть агентов смешанных
воплощений, два намерения (ресурс и отдых), приватные кадры,
глобально дрейфующие семейства, $300$ эпизодов, $12$ парных
test-сидов (номера $100{+}$, никогда не использованные для настройки;
все пороги заморожены с прежних волн). Групповая стоимость
(хрупкий/живучий): полный Songline $\mathbf{120{,}0/125{,}1}$
доминирует над независимостью ($137{,}3/146{,}5$), сырой общей
историей ($160{,}9/204{,}4$), эмбеддинговым векторным сходством
($205{,}5/187{,}5$ --- без тактов оно не переносит невиданное и
вырождается в fail-closed слепой поиск), графом без провенанса
($146{,}3/174{,}5$) и песенной лестницей без допуска
($146$--$157/174$--$196$) --- зарегистрированная H1a держится на
каждом арме; одновременно это единственное плечо с
нулевыми дублирующимися коммитами, наименьшей долей phantom-first и
проводом дешевле сырого ($55$ против $69$ КБ). Преимущество растёт с
масштабом --- $-11\%$, $-13\%$, $-14\%$ при $N = 4, 6, 8$
(зарегистрированная H3: \textbf{pass}) --- и переживает $10\%$ и
даже $20\%$ перцептивного шума ($-13$--$15\%$; клауза преимущества
H2: \textbf{pass}).

\paragraph{Матрица абляций --- и её честная половина.} Удаление
ориентиров или тактов катастрофично (успех $0$, стоимость худшая:
fail-closed по построению); удаление допуска стоит $+22\%/+38\%$ ---
самый несущий компонент, подтверждение S3; удаление часов мира ---
$+9\%$ на живучей роли; удаление резерваций возвращает дублирующиеся
коммиты; удаление utility-гейта раздувает провод и память на
${\sim}40\%$ при равной стоимости. Но три абляции --- провенанс,
исключения, неизменяемый слой --- бесплатны \emph{в этом субстрате}:
зарегистрированная H1 поэтому \textbf{проваливается} как
сформулирована, и механизм провала стоит самого провала. Их ценность
показана в средах, которые их нагружают (коррупция при возврате в
U1, gossip-взрыв в S1, аудит и реконсолидация); live-visit
бенчмарк с дрейфом их удаление не наказывает. В одном артефакте они
--- бесплатная страховка, не мёртвый груз, но клейм композиции обязан
быть ограничен механизмами, которые субстрат действительно нагружает.

\paragraph{Полезность без оракула (UE1).} Ridge-оцениватель только по
runtime-доступным признакам (якорение, локальная дистанция плана,
форма песни, счётчики памяти) достигает Спирмена $0{,}989$ и знаковой
точности $0{,}97$ против истинного контрфакта на held-out мирах, а
полный контроллер, ведомый одной $\widehat U$, совпадает с оракульной
версией (отношение стоимости $0{,}965/0{,}977$) --- критический
эксперимент рецензента: система переживает оценивание.

\paragraph{Урок о марже (N1).} Под независимым шумом ложных
пропусков $\times$ ложных срабатываний покрытие деградирует изящно
($0{,}96 \to 0{,}92$ при $10\%/10\%$: зарегистрировано
\textbf{pass}), но fail-closed контракт НЕ переживает наивного
смягчения порога: с простой косинусной маржой fail-open достигает
$0{,}21$ внутри конверта $20\%$ (N1.1: \textbf{fail}). Точное
равенство было несущим; названное лекарство --- верификация
консенсусом якорей (правило $\ge 3$ согласных пар слоя идентичности,
которое одно-якорное потребление песни обходит) --- машинерия маржи,
а не пониженный порог.


\section{Закрытие acceptance-теста: бюджеты, противник, безопасность,
второй субстрат}
\label{p5:sec:acceptance}

Рецензент предложил конкретный acceptance-тест: при равных бюджетах
полный метод должен бить лучший бейзлайн на невиданных динамических
мирах, не ухудшая безопасность; преимущество должно переживать шумную
перцепцию, приватные кадры, гетерогенность и дрейф на длинном
горизонте; а удаление провенанса, допуска, маршрутной структуры или
контракта времени мира должно воспроизводимо уничтожать
соответствующий выигрыш. Отчитываемся по каждой клаузе.

\paragraph{Статистическая опора.} На интеграционном бенчмарке полный
метод бьёт независимость на \textbf{$12/12$ парных test-сидах}
(средний командный выигрыш $13{,}6\%$, знаковый тест $p = 0{,}0002$),
а его попадания в опасности \emph{ниже} ($5{,}06$ против $5{,}66$ за
визит) --- преимущество не покупает скорость ценой безопасности.

\paragraph{Равные бюджеты (B1).} Ограничив провод и память
собственным следом полного метода ($55$/$45$ КБ) и дав тот же кап
сырой истории, full выигрывает $119{,}8/124{,}5$ против
$158{,}7/199{,}9$; и raw не догоняет ни при $2\times$, ни при
$4\times$ бюджета ($159{,}3/201{,}1$ всюду). Дефицит raw --- не
нехватка хранилища, которую чинит больший кап, а отсутствие
консолидации и допуска. (На континуальном субстрате raw не просто
хуже, а вычислительно нерешаем на полном горизонте --- $\mathrm{O}
(E^2)$ матчинг над несконсолидированным стором --- стоимость, которую
учёт байтов сам по себе не ловит.)

\paragraph{Adversarial provenance (P1).} При одном отравляющем пире,
который искажает свои песни и отмывает мутированные чужие записи под
их оригинальным origin, полный метод деградирует на $+1{,}2\%$, а
порядок обороны держится (full $\le$ без-провенанса $\le$
без-допуска). Поучительный негатив: незащищённое плечо почти не
деградирует под ядом, \emph{потому что уже насыщено фантомами
обычного дрейфа} ($159{,}7$ чистый). Главный противник коллективной
памяти --- энтропия мира, а не злоумышленник; допуск, гейтящий по
продемонстрированной полезности на собственном визите приёмника,
поглощает и лжеца, и устаревшее; origin-bound приём блокирует
отмывание без измеримой цены.

\paragraph{Безопасность под шумом (N1, пересмотр).} Планка fail-open
$<1\%$ не достигается смягчённой косинусной маржой (до $0{,}21$).
Названное лекарство --- требовать консенсус $\ge 2$ взаимно
сопоставленных якорей, чьи попарные смещения согласуются с цепочкой
тактов (loop-closure правило слоя идентичности, которое одно-якорное
потребление обходило), --- доводит \textbf{fail-open до $0{,}000$ по
всему конверту $20\%{\times}20\%$}, ценой покрытия ($0{,}88 \to
0{,}58$ при $10\%/10\%$). Безопасность покупается отказами, как и
требует fail-closed контракт.

\paragraph{Второй субстрат (C1).} Формирование и потребление
перенесены в непрерывный ящик с вещественными позициями признаков,
пообзорным джиттером, мягким сопоставлением констелляций
(позиционный допуск вместо точного равенства ключей), унимодальной
кластеризацией якорей с центроидным уточнением и loop-closure по
допуску --- рантайм в остальном не менялся. Преимущество переносится
на обеих ролях и против обоих бейзлайнов --- на согласованном
выполнимом горизонте full $245{,}5/71{,}4 <$ raw $256{,}9/101{,}2$ и
$<$ независимость $255{,}3/75{,}9$ (C1.3), а на полном горизонте full
$244{,}7/74{,}7 <$ независимость $256{,}8/79{,}6$ (C1.1), --- что
подтверждает: эффект не сеточный артефакт. Континуальная безопасность
на социальном масштабе требует \emph{трёхслойной} обороны, где каждый
слой ловит пропущенное предыдущим: консенсус якорей ($\ge 3$ взаимно
сопоставленных, замкнутых по циклу) снижает пер-транспортный алиасинг
($0{,}79 \to 0{,}15$); гейт единственного коммита (идти только к самой
заякоренной цели, а не по фантомному списку) доводит до $0{,}04$; а
\emph{верификация безопасного префикса} --- путник физически проходит
первые пять куплетов песни и требует совпадения каждой наблюдаемой
констелляции, так что чужесемейный фантом, расходящийся с реальной
местностью, отвергается до коммита, --- доводит fail-open до
$\mathbf{0{,}0014}$ на held-out test-сидах (в $554\times$ ниже сырого
континуального уровня, $10/12$ сидов ровно ноль), \emph{ниже} планки
безопасности $1\%$, причём успех первого лока \emph{растёт} (фантомы,
что раньше коммитились, теперь отфильтрованы). Проверка префикса ---
это <<успешное использование на коротком безопасном префиксе>> самого
рецензента, поднятое до жёсткого гейта; она же позволяет ослабить
порог консенсуса (вернув покрытие), ловя одиночные неверные якоря,
недоступные счётчику консенсуса. Калибровка заморожена на
dev-семействах ($5001{+}$); test-сиды ($100{+}$) не трогались.

\paragraph{Внешний субстрат (физика VMAS).}
Чтобы ответить на сильнейшее возражение рецензента --- что grid и
синтетический континуальный мир суть <<домашние>> миры авторов --- мы
перенесли на VMAS (сторонний симулятор многоагентной непрерывной
физики) не диагностический логгер, а \emph{весь рантайм}
(формирование, обмен сертификатами, карантин, приём на стороне
получателя, безкадровое непрерывное выравнивание, потребление по
безопасному префиксу). Голономные агенты движутся под
пропорциональным контроллером в боксе с ландмарками; получатель без
общего кадра восстанавливает засвидетельствованную цель-<<воду>>
исключительно сопоставлением констелляций ландмарок. Против
memory-free базлайна, которому дан плотный перебор всего бокса и
больший горизонт --- чтобы он \emph{мог} найти цель исчерпывающим
поиском, и контраст был честным ускорением, а не полом --- полный
рантайм достигает цели за $254{,}5$ против $543{,}5$ командных шагов
($-53{,}2\%$, $12/12$ парных сидов, знаковый тест $p=0{,}0005$), при
transport fail-open $\mathbf{0{,}0000}$: агенты, не сумевшие
заанкориться, \emph{отказываются} и переходят к тому же слепому
поиску, а не коммитят фантом. Перекалибрована лишь плотность субстрата
(число ландмарок, частота сэмплинга свидетеля, порог консенсуса под
геометрию констелляций); все флаги механизмов те же, что на grid.
Эффект и его fail-closed-безопасность выживают на третьем, внешне
созданном субстрате.

\section{Связанные работы}
\label{p5:sec:related}

\citet{rdd2026} утверждают, что память должна сохранять различия,
релевантные решению, а не описательное сходство, и гейтить слияние по
decision distortion; мы берём интуицию и добавляем координационную
сторону: гетерогенные приёмники, профили полезности по ролям и
сертификаты, пересчитываемые приёмником. Generative
Agents~\citep{genagents2023} оценивают recency, importance и
relevance --- наша важность есть измеренный контрфакт, и сравнение U7
квантифицирует разрыв. MemRL~\citep{memrl2026} учит полезность
\emph{использования} записанных эпизодов;
Mem-$\alpha$~\citep{memalpha2025} учит операции построения с
RL-наградой на вопросах-ответах --- наш контроллер учит
\emph{формирование} под координационной наградой с пробами на
коррупцию. MemGAS~\citep{memgas2025} выбирает гранулярность;
TrustMem~\citep{trustmem2026} верифицирует переходы консолидации
(coverage, preservation, faithfulness) --- наша операция исключения и
неизменяемый эпизодный слой атакуют тот же отказ со стороны
механизма, а наш выученный контроллер переоткрывает
append-not-overwrite из одной награды. \citet{faulty2026} показывают,
что непрерывная перезапись со временем портит память; наш U1 меряет
ту же коррупцию у similarity-бейзлайна и расценивает лекарство.
Насколько нам известно, ни одна система прежде не объединяла
структурную аналогию, причинную полезность, гетерогенные воплощения,
координационные последствия, эмерджентный статус ориентиров и
неизменяемый провенанс в одном измеренном цикле.

\section{Ограничения и открытые проблемы}
\label{p5:sec:limits}

(1) Сеточные субстраты и детерминированный replay; здесь контрфакт
точен, в других средах он должен стать оценщиком. (2) У аналогии теперь есть измеренный graph-matching движок
(\S\ref{p5:sec:echelon3}); роли и причинные отношения внутри
выравнивания остаются открытыми. (3) Причинный цикл
повышения замкнут над предзаданным суперсетом признаков
(\S\ref{p5:sec:gaps}); открытое порождение признаков из выученных
представлений --- новые перцептивные концепты, именуемые и проверяемые
между модальностями --- остаётся открытым. (4) Гейт доверие$\times$давность теперь вшит в потребление
(\S\ref{p5:sec:gaps}) и закрывает живучую роль; остаток хрупкой роли
($-0{,}9\%$) --- открытый хвост. (5) У кросс-семейной абстракции есть
измеренный первый шаг (кодбук сигнатур, \S\ref{p5:sec:gaps}); слоты и
роли уровня схем остаются открытыми. (6) Критерий оценки всюду --- \emph{экологическая
пригодность}: исполнимость на свежих мирах под координационным
контрактом, а не прямая верификация содержимого; формулировка
закрыта, её полное поколенческое измерение --- программа meta-RL.

\section{Заключение}

Формирование памяти, поставленное как ограниченный пятивариантный
выбор над двумя независимыми осями, измеримо (U1), выучиваемо (U2) и
эволюционируемо (U3) --- и, будучи выученным или отобранным, а не
спроектированным, возвращает уроки проектировщику: добавляй вместо
перезаписи, конечные бюджеты бесплатно, пороги у аттрактора селекции.
Сравнение семи политик честно расценивает всё пространство дизайна:
структура доминирует над эвристиками, консолидация --- по битам,
свежесть --- по успеху, а остаточный режим отказа --- устаревание при
дрейфе --- указывает ровно на тот закон, который остальная часть
этого тома уже доказала. Честная статусная строка: онтогенез единицы
памяти --- и её индивидуальное, гетерогенное, LLM-опосредованное
потребление --- закрыт в масштабе прототипа; а устойчивый долгосрочный
\emph{социальный} обмен в эволюционирующем мире теперь закрыт в том
же масштабе --- трёхступенчатой измеренной лестницей: обмен без
контрактов вреден (S1); staleness на часах мира необходим, но
недостаточен (S2); допуск по продемонстрированной полезности на
собственных визитах заставляет коллектив бить независимость на
каждом сиде и обеих ролях (S3, $-20\%$). Весь рантайм, а не только
эффект, переносится на сторонний симулятор непрерывной физики (VMAS,
$-53{,}2\%$ командных шагов, fail-open $0{,}0000$). Открытым остаётся
внешнее в ином смысле: открытое порождение признаков из выученных
представлений, социальный процесс на языковых (LLM) субстратах и
валидность на реальных роботах.



\FloatBarrier
\begin{center}\rule{0.55\linewidth}{0.4pt}\\[2pt]\textbf{Приложения к этой части}\end{center}


\section{Воспроизводимость}
\label{p5:app:repro}
\begin{verbatim}
PYTHONPATH=. python experiments/song_grammar/exp_u1_ucsm_smoke.py
PYTHONPATH=. python experiments/song_grammar/exp_u2_bandit.py \
    --train-seeds 0 24 --eval-seeds 100 110 --episodes 40
PYTHONPATH=. python experiments/song_grammar/exp_u3_evolution.py \
    --generations 30 --pop 16 --episodes 25
PYTHONPATH=. python experiments/song_grammar/exp_u7_seven_arms.py \
    --seeds 0 12 --episodes 1000
\end{verbatim}
Зарегистрированные предсказания, построчные результаты и вердикты
хранятся рядом с выводом каждого эксперимента; обе провалившиеся
первые регистрации (U1 v1, U2 v1) сохранены с примечаниями о
ревизиях. Дизайн-документы:
\texttt{docs/FRONTIER\_SONG\_GRAMMAR\_2026-07-25.md},
\texttt{docs/FRONTIER\_UCSM\_2026-07-27.md},
\texttt{docs/MEMORY\_REFRAMES\_2026-07-26.md}.
\FloatBarrier

\clearpage
\begin{thebibliography}{99}
\bibitem[Andrychowicz et~al.(2017)]{andrychowicz2017hindsight}
M.~Andrychowicz, F.~Wolski, A.~Ray, J.~Schneider, R.~Fong, P.~Welinder, B.~McGrew, J.~Tobin, P.~Abbeel, and W.~Zaremba.
\newblock Hindsight experience replay.
\newblock In \emph{NeurIPS}, 2017.

\bibitem[B\v{e}lohl\'avek(1999)]{belohlavek1999fuzzy}
R.~B\v{e}lohl\'avek.
\newblock Fuzzy Galois connections.
\newblock \emph{Mathematical Logic Quarterly}, 45(4):497--504, 1999.

\bibitem[Bettini et~al.(2022)]{bettini2022vmas}
M.~Bettini, R.~Kortvelesy, J.~Blumenkamp, and A.~Prorok.
\newblock VMAS: A vectorized multi-agent simulator for collective robot learning.
\newblock In \emph{16th International Symposium on Distributed Autonomous Robotic Systems (DARS)}, 2022.
\newblock arXiv:2207.03530.

\bibitem[Boyd et~al.(2006)]{boyd2006gossip}
S.~Boyd, A.~Ghosh, B.~Prabhakar, and D.~Shah.
\newblock Randomized gossip algorithms.
\newblock \emph{IEEE Trans.\ Information Theory}, 52(6):2508--2530, 2006.

\bibitem[Chevalier-Boisvert et~al.(2018)]{chevalier2018babyai}
M.~Chevalier-Boisvert, D.~Bhoopchand, H.~Abdolmaleki, and A.~Mordatch.
\newblock BabyAI: A platform to study the sample efficiency of grounded language learning.
\newblock In \emph{ICLR}, 2019.

\bibitem[Chevalier-Boisvert et~al.(2023)]{MinigridMiniworld23}
M.~Chevalier-Boisvert, B.~Lahlou, S.~Willems, and P.~S. Castro.
\newblock MiniGrid and MiniWorld: Modular reinforcement learning environments.
\newblock \emph{arXiv:2306.13831}, 2023.

\bibitem[Chhikara et~al.(2025)]{chhikara2025mem0}
P.~Chhikara, D.~Khant, S.~Aryan, T.~Singh, and D.~Yadav.
\newblock Mem0: Building production-ready AI agents with scalable long-term memory.
\newblock \emph{arXiv:2504.19413}, 2025.

\bibitem[Dayan(1993)]{dayan1993improving}
P.~Dayan.
\newblock Improving generalization for temporal difference learning: The successor representation.
\newblock \emph{Neural Computation}, 5(4):613--624, 1993.

\bibitem[Deitke et~al.(2022)]{deitke2022procthor}
M.~Deitke, E.~V. Wijmans, D.~Schwenk, O.~Salvador, K.~M. Ehsani, et al.
\newblock ProcTHOR: Large-scale embodied AI using procedural generation.
\newblock In \emph{NeurIPS Datasets and Benchmarks}, 2022.

\bibitem[Efron(1979)]{efron1979bootstrap}
B.~Efron.
\newblock Bootstrap methods: Another look at the jackknife.
\newblock \emph{Annals of Statistics}, 7(1):1--26, 1979.

\bibitem[Foerster et~al.(2016)]{foerster2016learning}
J.~N. Foerster, Y.~M. Assael, N.~de~Freitas, and S.~Whiteson.
\newblock Learning to communicate with deep multi-agent reinforcement learning.
\newblock In \emph{NeurIPS}, 2016.

\bibitem[Foerster et~al.(2018)]{foerster2018ctde}
J.~N. Foerster, G.~Farquhar, T.~Afouras, N.~Nardelli, and S.~Whiteson.
\newblock Counterfactual multi-agent policy gradients.
\newblock In \emph{AAAI}, 2018.

\bibitem[Ganter and Wille(1999)]{ganter1999fca}
B.~Ganter and R.~Wille.
\newblock \emph{Formal Concept Analysis: Mathematical Foundations}.
\newblock Springer, 1999.

\bibitem[Ghosh et~al.(2021)]{ghosh2021learning}
D.~Ghosh, A.~Gupta, A.~Reddy, J.~Fu, S.~Levine, and C.~Finn.
\newblock Learning to reach goals via iterated supervised learning.
\newblock In \emph{ICLR}, 2021.

\bibitem[Gigerenzer and Goldstein(1996)]{gigerenzer1996reasoning}
G.~Gigerenzer and D.~G. Goldstein.
\newblock Reasoning the fast and frugal way: Models of bounded rationality.
\newblock \emph{Psychological Review}, 103(4):650--669, 1996.

\bibitem[Gupta et~al.(2017)]{gupta2017cognitive}
S.~Gupta, J.~Davidson, S.~Levine, R.~Sukthankar, and J.~Malik.
\newblock Cognitive mapping and planning for visual navigation.
\newblock In \emph{CVPR}, 2017.

\bibitem[Holm(1979)]{holm1979simple}
S.~Holm.
\newblock A simple sequentially rejective multiple test procedure.
\newblock \emph{Scandinavian Journal of Statistics}, 6(2):65--70, 1979.

\bibitem[Kolve et~al.(2017)]{kolve2017ai2thor}
E.~Kolve, R.~Mottaghi, D.~Gordon, Y.~Zhu, A.~Gupta, A.~Farhadi, and D.~Fox.
\newblock AI2-THOR: An interactive 3D environment for visual AI.
\newblock \emph{arXiv:1712.05474}, 2017.

\bibitem[Kuipers(2000)]{kuipers2000spatial}
B.~Kuipers.
\newblock The spatial semantic hierarchy.
\newblock \emph{Artificial Intelligence}, 119(1--2):191--233, 2000.

\bibitem[Lieder and Griffiths(2020)]{lieder2020resource}
F.~Lieder and T.~L. Griffiths.
\newblock Resource-rational analysis: Understanding human cognition as the optimal use of limited computational resources.
\newblock \emph{Behavioral and Brain Sciences}, 43:e1, 2020.

\bibitem[Lowe et~al.(2017)]{lowe2017maddpg}
R.~Lowe, Y.~Wu, A.~Tamar, J.~Harb, P.~Abbeel, and I.~Mordatch.
\newblock Multi-agent actor-critic for mixed cooperative-competitive environments.
\newblock In \emph{NeurIPS}, 2017.

\bibitem[Mac~Lane(1998)]{maclane1998categories}
S.~Mac~Lane.
\newblock \emph{Categories for the Working Mathematician}.
\newblock Springer, 2nd edition, 1998.

\bibitem[Momennejad et~al.(2017)]{momennejad2017successor}
I.~Momennejad, E.~M. Russek, J.~H. Cheong, M.~M. Botvinick, N.~Daw, and S.~J. Gershman.
\newblock The successor representation in human reinforcement learning.
\newblock \emph{Nature Human Behaviour}, 1:680--692, 2017.

\bibitem[O'Keefe and Nadel(1978)]{okeefe1978hippocampus}
J.~O'Keefe and L.~Nadel.
\newblock \emph{The Hippocampus as a Cognitive Map}.
\newblock Oxford University Press, 1978.

\bibitem[Packer et~al.(2023)]{packer2023memgpt}
C.~Packer, V.~Fang, S.~Patil, K.~Lin, S.~Wooders, and J.~Gonzalez.
\newblock MemGPT: Towards LLMs as operating systems.
\newblock \emph{arXiv:2310.08560}, 2023.

\bibitem[Park et~al.(2023)]{park2023generative}
J.~S. Park, J.~C. O'Brien, C.~J. Cai, M.~R. Morris, P.~Liang, and M.~S. Bernstein.
\newblock Generative agents: Interactive simulacra of human behavior.
\newblock \emph{arXiv:2304.03442}, 2023.

\bibitem[Rashid et~al.(2018)]{rashid2018qmix}
T.~Rashid, M.~Samvelyan, C.~Schroeder de Witt, G.~Farquhar, J.~Foerster, and S.~Whiteson.
\newblock QMIX: Monotonic value function factorisation for deep multi-agent reinforcement learning.
\newblock In \emph{ICML}, 2018.

\bibitem[Savinov et~al.(2018)]{savinov2018semi}
N.~Savinov, A.~Dosovitskiy, and V.~Koltun.
\newblock Semi-parametric topological memory for navigation.
\newblock In \emph{ICLR}, 2018.

\bibitem[Savva et~al.(2019)]{savva2019habitat}
M.~Savva, A.~Kadian, O.~Maksymets, Y.~Zhao, E.~Wijmans, B.~Jain, J.~Straub, J.~Liu, V.~Koltun, J.~Malik, et al.
\newblock Habitat: A platform for embodied AI research.
\newblock In \emph{ICCV}, 2019.

\bibitem[Schulman et~al.(2016)]{schulman2016gae}
J.~Schulman, P.~Moritz, S.~Levine, M.~Jordan, and P.~Abbeel.
\newblock High-dimensional continuous control using generalized advantage estimation.
\newblock In \emph{ICLR}, 2016.
\newblock arXiv:1506.02438.

\bibitem[Schulman et~al.(2017)]{schulman2017proximal}
J.~Schulman, F.~Wolski, P.~Dhariwal, A.~Radford, and O.~Klimov.
\newblock Proximal policy optimization algorithms.
\newblock \emph{arXiv:1707.06347}, 2017.

\bibitem[Shinn et~al.(2023)]{shinn2023reflexion}
N.~Shinn, F.~Cassano, A.~Berman, and A.~G. Labash.
\newblock Reflexion: Language agents with verbal reinforcement learning.
\newblock In \emph{NeurIPS}, 2023.

\bibitem[Simon(1955)]{simon1955behavioral}
H.~A. Simon.
\newblock A behavioral model of rational choice.
\newblock \emph{Quarterly Journal of Economics}, 69(1):99--118, 1955.

\bibitem[Simon(1956)]{simon1956rational}
H.~A. Simon.
\newblock Rational choice and the structure of the environment.
\newblock \emph{Psychological Review}, 63(2):129--138, 1956.

\bibitem[Stachenfeld et~al.(2017)]{stachenfeld2017hippocampus}
K.~L. Stachenfeld, M.~M. Botvinick, and S.~J. Gershman.
\newblock The hippocampus as a predictive map.
\newblock \emph{Nature Neuroscience}, 20(11):1643--1653, 2017.

\bibitem[Sukhbaatar et~al.(2016)]{sukhbaatar2016commnet}
S.~Sukhbaatar, A.~Szlam, and R.~Fergus.
\newblock Learning multiagent communication with backpropagation.
\newblock In \emph{NeurIPS}, 2016.

\bibitem[Sutton et~al.(1999)]{sutton1999between}
R.~S. Sutton, D.~Precup, and S.~Singh.
\newblock Between MDPs and semi-MDPs: A framework for temporal abstraction in reinforcement learning.
\newblock \emph{Artificial Intelligence}, 112(1--2):181--211, 1999.

\bibitem[Tolman(1948)]{tolman1948cognitive}
E.~C. Tolman.
\newblock Cognitive maps in rats and men.
\newblock \emph{Psychological Review}, 55(4):189--208, 1948.

\bibitem[Wang et~al.(2023)]{wang2023voyager}
G.~Wang, Y.~Xie, Z.~Jiang, A.~Mandlekar, C.~Xie, A.~Yu, and P.~Liang.
\newblock Voyager: An open-ended embodied agent with large language models.
\newblock \emph{arXiv:2305.16291}, 2023.

\bibitem[Wu et~al.(2023)]{wu2023autogen}
Q.~Wu, G.~Bansal, J.~Zhang, Y.~Wu, B.~Li, E.~Zhu, L.~Jiang, X.~Zhang, S.~Zhang, J.~Liu, A.~H. Awadallah, R.~W. White, D.~Burger, and C.~Wang.
\newblock AutoGen: Enabling next-gen LLM applications via multi-agent conversation.
\newblock \emph{arXiv:2308.08155}, 2023.

\bibitem[Wilcoxon(1945)]{wilcoxon1945individual}
F.~Wilcoxon.
\newblock Individual comparisons by ranking methods.
\newblock \emph{Biometrics Bulletin}, 1(6):80--83, 1945.

\bibitem[Williams(1992)]{williams1992simple}
R.~J. Williams.
\newblock Simple statistical gradient-following algorithms for connectionist reinforcement learning.
\newblock \emph{Machine Learning}, 8(3--4):229--256, 1992.

\bibitem[Yao et~al.(2023)]{yao2022react}
S.~Yao, J.~Zhao, D.~Yu, N.~Du, I.~Shafran, K.~Narasimhan, and Y.~Cao.
\newblock ReAct: Synergizing reasoning and acting in language models.
\newblock In \emph{ICLR}, 2023.

\bibitem[Xiao et~al.(2007)]{xiao2007distributed}
L.~Xiao, S.~Boyd, and S.~Lall.
\newblock A scheme for robust distributed sensor fusion based on average consensus.
\newblock \emph{IPSN}, 2007.

\bibitem[Yu et~al.(2022)]{yu2022mappo}
C.~Yu, A.~Velu, E.~Vinitsky, J.~Gao, Y.~Wang, A.~Bayen, and Y.~Wu.
\newblock The surprising effectiveness of PPO in cooperative multi-agent games.
\newblock In \emph{NeurIPS Datasets and Benchmarks Track}, 2022.

\bibitem[Das et~al.(2019)]{das2019tarmac}
A.~Das, T.~Gervet, J.~Romoff, D.~Batra, D.~Parikh, M.~Rabbat, and
J.~Pineau.
\newblock TarMAC: Targeted multi-agent communication.
\newblock In \emph{ICML}, 2019.

\bibitem[Jefferson(1985)]{jefferson1985}
D.~R. Jefferson.
\newblock Virtual time.
\newblock \emph{ACM Transactions on Programming Languages and Systems},
7(3):404--425, 1985.

\bibitem[Cadena et~al.(2016)]{cadena2016slam}
C.~Cadena, L.~Carlone, H.~Carrillo, Y.~Latif, D.~Scaramuzza,
J.~Neira, I.~Reid, and J.~J. Leonard.
\newblock Past, present, and future of simultaneous localization and
mapping: Toward the robust-perception age.
\newblock \emph{IEEE Transactions on Robotics}, 32(6):1309--1332, 2016.

\bibitem[Lowry et~al.(2016)]{lowry2016vpr}
S.~Lowry, N.~S\"underhauf, P.~Newman, J.~J. Leonard, D.~Cox,
P.~Corke, and M.~J. Milford.
\newblock Visual place recognition: A survey.
\newblock \emph{IEEE Transactions on Robotics}, 32(1):1--19, 2016.

\bibitem[Aczel(1988)]{aczel1988nonwellfounded}
P.~Aczel.
\newblock \emph{Non-Well-Founded Sets}.
\newblock CSLI Lecture Notes 14, 1988.

\bibitem[Barwise and Seligman(1997)]{barwise1997information}
J.~Barwise and J.~Seligman.
\newblock \emph{Information Flow: The Logic of Distributed Systems}.
\newblock Cambridge University Press, 1997.

\bibitem[Jacobs(2016)]{jacobs2016coalgebra}
B.~Jacobs.
\newblock \emph{Introduction to Coalgebra: Towards Mathematics of States and Observation}.
\newblock Cambridge University Press, 2016.

\bibitem[Rao and Georgeff(1995)]{rao1995bdi}
A.~S. Rao and M.~P. Georgeff.
\newblock BDI agents: From theory to practice.
\newblock In \emph{ICMAS}, 1995.

\bibitem[Rutten(2000)]{rutten2000universal}
J.~J.~M.~M. Rutten.
\newblock Universal coalgebra: a theory of systems.
\newblock \emph{Theoretical Computer Science}, 249(1):3--80, 2000.

\bibitem[Sangiorgi(2011)]{sangiorgi2011bisimulation}
D.~Sangiorgi.
\newblock \emph{Introduction to Bisimulation and Coinduction}.
\newblock Cambridge University Press, 2011.

\bibitem[Sperber and Wilson(1995)]{sperber1995relevance}
D.~Sperber and D.~Wilson.
\newblock \emph{Relevance: Communication and Cognition}.
\newblock Blackwell, 2nd edition, 1995.

\bibitem[Anonymous(2026a)]{rdd2026}
Anonymous.
\newblock Remember the decision, not the description: A
rate-distortion framework for agent memory.
\newblock \emph{arXiv preprint arXiv:2605.10870}, 2026.

\bibitem[Park et~al.(2023)]{genagents2023}
J.~S. Park, J.~O'Brien, C.~J. Cai, M.~R. Morris, P.~Liang, and
M.~S. Bernstein.
\newblock Generative agents: Interactive simulacra of human behavior.
\newblock In \emph{UIST}, 2023.

\bibitem[Anonymous(2026b)]{memrl2026}
Anonymous.
\newblock {MemRL}: Self-evolving agents via runtime reinforcement
learning on episodic memory.
\newblock \emph{arXiv preprint arXiv:2601.03192}, 2026.

\bibitem[Anonymous(2025a)]{memalpha2025}
Anonymous.
\newblock {Mem-$\alpha$}: Learning memory construction via
reinforcement learning.
\newblock \emph{arXiv preprint arXiv:2509.25911}, 2025.

\bibitem[Anonymous(2025b)]{memgas2025}
Anonymous.
\newblock From single to multi-granularity: Toward long-term memory
association and selection of conversational agents.
\newblock \emph{arXiv preprint arXiv:2505.19549}, 2025.

\bibitem[Anonymous(2026c)]{trustmem2026}
Anonymous.
\newblock {TrustMem}: Learning trustworthy memory consolidation for
{LLM} agents with long-term memory.
\newblock \emph{arXiv preprint arXiv:2606.25161}, 2026.

\bibitem[Anonymous(2026d)]{faulty2026}
Anonymous.
\newblock Useful memories become faulty when continuously updated by
{LLMs}.
\newblock \emph{arXiv preprint arXiv:2605.12978}, 2026.
\end{thebibliography}
\end{document}
